\documentclass[11pt]{article}
% Compile with XeLaTeX or Tectonic (Unicode + Korean).
\usepackage[margin=1in]{geometry}
\usepackage{amsmath,amssymb,amsthm}
\usepackage{tcolorbox}
\tcbuselibrary{breakable}
\usepackage{enumitem}
\usepackage{tabularx}
\usepackage[hidelinks]{hyperref}
\setlength{\parindent}{0pt}
\setlength{\parskip}{4pt}
\setlist[enumerate,1]{label=\arabic*.,leftmargin=*,itemsep=1em}
\setlist[enumerate,2]{leftmargin=*,itemsep=.35em}
\newcommand{\R}{\mathbb R}
\newcommand{\C}{\mathbb C}
\newcommand{\Ric}{\operatorname{Ric}}
\newcommand{\supp}{\operatorname{supp}}
\newcommand{\id}{\operatorname{Id}}
\newcommand{\ddbar}{\partial\bar\partial}
\newcommand{\cR}{c_1(M)_{\R}}
\definecolor{theoremcream}{RGB}{247,241,228}
\definecolor{theorembrown}{RGB}{106,85,56}
\definecolor{theorembody}{RGB}{254,254,254}
\definecolor{theoremequationborder}{RGB}{110,110,110}
\newcommand{\theoremboxed}[1]{%
  {\color{theoremequationborder}\boxed{\color{black}#1}}%
}
\tcbset{namedstatement/.style={
  colback=theorembody,
  colframe=theorembrown,
  colbacktitle=theoremcream,
  coltitle=black,
  fonttitle=\normalsize\rmfamily\bfseries,
  fontupper=\normalfont,
  boxrule=0.5pt,
  boxsep=0pt,
  arc=1.5mm,
  left=4mm,
  right=4mm,
  top=2mm,
  bottom=2mm,
  toptitle=0.8mm,
  bottomtitle=0.8mm,
  before skip=12pt,
  after skip=14pt,
  before upper={\setlength{\parindent}{0pt}\setlength{\parskip}{6pt}}
}}
\newtcolorbox{theorem}{
  namedstatement,
  title={Calabi--Yau Theorem (Calabi Conjecture)}
}
\newcounter{definition}
\newcounter{subdefinition}[definition]
\renewcommand{\thesubdefinition}{\thedefinition.\arabic{subdefinition}}
\newtcolorbox[use counter*=definition]{definitionbox}[1]{
  namedstatement,
  breakable,
  title={\thetcbcounter.\ #1},
  title after break={\thetcbcounter.\ #1\ (continued)}
}
\newtcolorbox[use counter*=subdefinition]{subdefinitionbox}[1]{
  namedstatement,
  breakable,
  title={\thetcbcounter.\ #1},
  title after break={\thetcbcounter.\ #1\ (continued)}
}

\usepackage{fontspec}
\usepackage{kotex}
\IfFontExistsTF{Malgun Gothic}{\setmainhangulfont{Malgun Gothic}}{\IfFontExistsTF{Noto Serif CJK KR}{\setmainhangulfont{Noto Serif CJK KR}}{}}
\usepackage{longtable}
\newcommand{\CP}{\mathbf{CP}}
\newcommand{\Q}{\mathbb Q}
\newcommand{\Z}{\mathbb Z}
\newtcolorbox{atlasdefinition}[1]{namedstatement,breakable,title={#1},title after break={#1\ (continued)}}
\newcommand{\exampleheading}{\par\medskip\textbf{구체적 예시 / Example on $\CP^1$ or $S^2$}\par\smallskip}
\setlength{\emergencystretch}{3em}
\hypersetup{pdftitle={Definitions with CP1 and S2 examples},pdfauthor={Definition study notes}}
\begin{document}

\begin{center}
{\LARGE\bfseries 정의와 구체적 예시}\par\medskip
{\Large Definitions on $\CP^1$ and $S^2\subset\R^3$}\par\medskip
받은 TeX 원고 + 통합 정의집의 고정 태그 + 모든 정의의 예시
\end{center}
\begin{tcolorbox}[namedstatement,breakable,title={읽는 순서 / Reading routes}]
받은 원고의 정의 25개(하위 정의 포함)를 보존하고 각각 예시를 붙였습니다. 뒤에는 우리 정의집 161개를 기존 분야와 고정 태그로 수록하며, 각 항목마다 별도의 예시와 직접 선행 정의를 제공합니다. 중복은 원고의 맥락을 보존하기 위해 유지하고 상호 참조합니다.

이번 발표의 정의 경로는 \hyperref[def:B10]{B10 정준선다발의 Chern 곡률}에서 거꾸로 읽으면 됩니다. Calabi--Yau 정리 부분은 받은 원고의 참고 내용으로 보존했습니다. 그 정리의 증명을 이번 발표 목표로 설정하는 것은 아닙니다.

일반 정의와 아래의 구체적 예시는 구별합니다. 한 예시를 확인하는 것으로 모든 공간에 대한 정리를 증명한 것은 아닙니다. 특히 fine acyclicity, Cartan B, de Rham 비교 등은 사용한 정리로 명시합니다.
\end{tcolorbox}
\tableofcontents
\clearpage
\section{공통 예시와 부호 규약 / Fixed examples and conventions}
\label{sec:conventions}
\[
S^2=\{(X,Y,Z)\in\R^3:X^2+Y^2+Z^2=1\},\quad
N=(0,0,1),\quad S=(0,0,-1).
\]
The sphere has the subspace topology and the outward orientation. Identify it with $\CP^1$ by the charts
\[
z=\frac{X+iY}{1+Z}\quad\text{on }U_0=S^2\setminus\{S\},\qquad
w=\frac{X-iY}{1-Z}\quad\text{on }U_\infty=S^2\setminus\{N\}.
\]
Then $w=1/z$. In homogeneous coordinates $[S:T]$ (the homogeneous symbol S is distinct from the south pole), $z=T/S$, $w=S/T$, and $\infty=[0:1]$.
Put $D=1+x^2+y^2$ for $z=x+iy$. The inverse of the z-chart is
\[
\Phi(z)=\frac{(2x,2y,1-x^2-y^2)}{D}.
\]
The unit round metric and its area form are
\[
g=\frac4{D^2}(dx^2+dy^2),\quad
\Omega=\frac4{D^2}dx\wedge dy=\frac{2i}{D^2}dz\wedge d\bar z,
\quad\int_{S^2}\Omega=4\pi.
\]
Here $g_{z\bar z}=g_{\C}(\partial_z,\partial_{\bar z})=2/D^2$.
The tangent-line Hermitian metric has $h_T(\partial_z,\partial_z)=2/D^2$; its dual on $K$ has $h_K(dz,dz)=D^2/2$.
\[
A_K=\frac{2\bar z}{D}dz,\quad F_K=-\frac2{D^2}dz\wedge d\bar z=i\Omega,
\quad \frac{i}{2\pi}\int_{\CP^1}F_K=-2.
\]
We take Hermitian forms linear in the first argument, $\Theta=\nabla^2$, and $c_1(L,h)=iF/(2\pi)$. The Gaussian curvature is the scalar $K_G=1$; it is not the two-form $F_K$.

\textbf{Analytic versus algebraic.} Unless explicitly stated otherwise, $\mathcal O_X$ on $\CP^1$ means the analytic holomorphic-function sheaf. Thus $\mathcal O_X(U_0)=\mathcal O(\C)$, not $\C[z]$. Algebraic examples explicitly say so. A skyscraper $i_*A$ at p has sections A when $p\in U$ and zero otherwise; these examples are in abelian sheaves unless stated to be module sheaves.

\textbf{A fixed triangulated sphere.} Use the four unit vectors proportional to $(1,1,1)$, $(1,-1,-1)$, $(-1,1,-1)$, $(-1,-1,1)$. Radial projection identifies the boundary K of their tetrahedron with $S^2$. The open star of a vertex is the locus where its barycentric coordinate is positive. Its nonempty multiple intersections are the open stars of faces. They are contractible, and the fourfold intersection is empty. These four sets give the good cover used in the homology and Čech examples. Reorder vertex labels if necessary so that the displayed fundamental cycle agrees with the outward orientation.

\textbf{A fixed smooth cutoff.} Let $\eta(t)=e^{-1/t}$ for $t>0$ and $\eta(t)=0$ otherwise. Put
\[
\rho_0=\frac{\eta(Z+1/2)}{\eta(Z+1/2)+\eta(1/2-Z)},\qquad
\rho_\infty=1-\rho_0.
\]
The denominator is positive, $\supp\rho_0\subset\{Z\ge-1/2\}\subset U_0$, and $\supp\rho_\infty\subset\{Z\le1/2\}\subset U_\infty$.
\clearpage
\section{받은 원고 / Supplied manuscript with examples}


% Replacement for the supplied sections; retain your existing preamble,
% theorem environment, and namedstatement tcolorbox style.

\section*{Calabi--Yau Theorem (Calabi Conjecture)}

\begin{theorem}
Let $(M,\omega_0)$ be a compact K\"ahler manifold with $\dim_{\mathbb C} M=n$.
\[
\left(\Rightarrow[\operatorname{Ric}(\omega_0)]
=:2\pi c_1(M)_{\mathbb R}
\in H_{\mathrm{dR}}^2(M;\mathbb R)\right)
\]
Then for any smooth closed real $(1,1)$-form $\rho$ satisfying
\[
[\rho]=2\pi c_1(M)_{\mathbb R},
\]
there exists a unique K\"ahler metric $\omega$ such that
\[
\begin{cases}
\;\operatorname{Ric}(\omega)=\rho,\\
\;[\omega]=[\omega_0]\in H_{\mathrm{dR}}^2(M;\mathbb R).
\end{cases}
\]
\end{theorem}
\begin{tcolorbox}[namedstatement,breakable,title={정리의 구체적 입력 / A concrete input on $\CP^1$}]
Take $M=\CP^1$, $\omega_0=\Omega$ from the unit sphere, and prescribe $\rho=\Omega$. The coefficient is $g_{z\bar z}=2/D^2$, so $\operatorname{Ric}(\omega_0)=\Omega$. Choose $F=0$ and $\varphi=0$; then the normalized volume equation and Ricci equation hold. The existence and uniqueness for arbitrary admissible $\rho$ are the theorem, not proved by this single example. On $\CP^1$ one cannot prescribe $\rho=0$, because $\int\rho$ would be 0 whereas $\int\operatorname{Ric}(\omega_0)=4\pi$.
\end{tcolorbox}

\begin{tcolorbox}[
  namedstatement,
  title={Local-Coordinate Formulation}
]
In local holomorphic coordinates $(z^1,\ldots,z^n)$, write
\[
\begin{aligned}
\omega_0
&=\sqrt{-1}\sum_{i,j=1}^{n}
g_{i\bar j}\,dz^i\wedge d\bar z^j,\\[6pt]
\omega
&=\sqrt{-1}\sum_{i,j=1}^{n}
\left(
g_{i\bar j}
+\frac{\partial^2\varphi}
{\partial z^i\,\partial\bar z^j}
\right)
dz^i\wedge d\bar z^j.\\[6pt]
&(\Longleftrightarrow
[\omega]=[\omega_0]
\;\text{by the }\partial\bar\partial\text{-lemma.})
\end{aligned}
\]
\[
\begin{gathered}
\operatorname{Ric}(\omega)=\rho
\;\Longleftrightarrow\;
\det\left(
g_{i\bar j}
+\frac{\partial^2\varphi}
{\partial z^i\,\partial\bar z^j}
\right)
=\boxed{e^F\det\left(g_{i\bar j}\right)}
\end{gathered}
\]

\textbf{Note}: \(F\in C^\infty(M,\mathbb R)\) is the normalized
smooth function determined by \(\rho\) relative to \(\omega_0\).

\end{tcolorbox}

\newpage
\begingroup
\renewcommand{\arraystretch}{1.25}
\begin{tabularx}{\textwidth}{@{}>{\raggedright\arraybackslash}p{0.32\textwidth}>{\raggedright\arraybackslash}X@{}}
\textbf{Notation} & \textbf{Meaning} \\
\hline
$\omega_0$
& ``Omega naught'': the reference K\"ahler form. \\

$(M,\omega_0)$
& ``M equipped with the K\"ahler form $\omega_0$.'' \\

$\dim_{\mathbb C} M=n$
& ``The complex dimension of M is n.'' \\

$\operatorname{Ric}(\omega_0)$
& ``The Ricci form of $\omega_0$.'' \\

$[\operatorname{Ric}(\omega_0)]$
& ``The de Rham cohomology class of the Ricci form of $\omega_0$.'' \\

$c_1(M)_{\mathbb R}$
& ``The first Chern class of M with real coefficients.'' \\

$H_{\mathrm{dR}}^2(M;\mathbb R)$
& ``The second de Rham cohomology group of M with real coefficients.'' \\

$\rho$
& ``Rho'': the prescribed smooth, closed, real differential form
of type one-one that is to be the Ricci form. \\

$(1,1)\text{-form}$
& ``A differential form of type one-one.'' \\

$[\rho]$
& ``The de Rham cohomology class of rho.'' \\

$[\omega]=[\omega_0]$
& ``The de Rham cohomology class of omega equals the de Rham
cohomology class of $\omega_0$.'' \\

$\displaystyle
\frac{\partial^2\varphi}{\partial z^i\,\partial\bar z^j}$
& ``The mixed second partial derivative of phi with respect to
z superscript i and the complex conjugate of z superscript j.'' \\

$g$
& ``g'': the reference K\"ahler metric tensor associated with $\omega_0$. \\

$g_{i\bar j}$
& ``The corresponding component
of the reference K\"ahler metric in the chosen holomorphic coordinates.'' \\

$(g_{i\bar j})$
& ``The matrix whose entry in row i and column j is g with
subscripts i and barred j'': the Hermitian coordinate matrix
of the reference K\"ahler metric.'' \\

$\det(g_{i\bar j})$
& ``The determinant of the matrix whose entry in row i and column j
is g with subscripts i and barred j.'' \\

$\displaystyle
\det\!\left(
g_{i\bar j}
+\frac{\partial^2\varphi}{\partial z^i\,\partial\bar z^j}
\right)$
& ``The determinant of the matrix whose entry in row i and column j
is the sum of g with subscripts i and barred j and the mixed second
partial derivative of phi with respect to z superscript i and the
complex conjugate of z superscript j.'' \\

$F\in C^\infty(M,\mathbb R)$
& ``F is a smooth function from M to $\mathbb R$.'' \\
\end{tabularx}
\endgroup

\newpage

\subsection*{Explanations}

\textbf{1. The cohomological condition.}

We use the Ricci-form convention
\[
\operatorname{Ric}(\omega_0)
=-\sqrt{-1}\,\partial\bar\partial
\log\det(g_{i\bar j}),
\]
where $(g_{i\bar j})$ is the coefficient matrix of $\omega_0$.
With this convention,
$[\operatorname{Ric}(\omega_0)]=2\pi c_1(M)_{\mathbb R}$,
where $c_1(M)_{\mathbb R}$ is the image of the integral first Chern class
in real de Rham cohomology. Changing the metric preserves this class,
so the prescribed form $\rho$ must be closed and represent it.
The form $\rho$ need not be positive.

\textbf{2. Why the potential describes a fixed K\"ahler class.}

The term $\sqrt{-1}\,\partial\bar\partial\varphi$ is exact, so adding
it to $\omega_0$ preserves the de Rham class. Conversely, on a compact
K\"ahler manifold, the $\partial\bar\partial$-lemma says that the
difference of two cohomologous K\"ahler forms has this form for a
globally defined smooth real function $\varphi$.

The condition $\omega>0$ means that the Hermitian coefficient matrix
in the Local-Coordinate Formulation is positive-definite at every point.
It ensures that the resulting form defines a K\"ahler metric;
a positive determinant alone is not sufficient in general.
Thus the unknown metric is encoded by one real-valued potential
$\varphi$, subject to positivity.

\textbf{3. How $F$ is chosen.}

The function $F$ in the Local-Coordinate Formulation is chosen
from the prescribed form $\rho$. The cohomological hypothesis gives
$[\operatorname{Ric}(\omega_0)-\rho]=0$, so the
$\partial\bar\partial$-lemma provides a smooth real function $F$
satisfying
\[
\operatorname{Ric}(\omega_0)-\rho
=\sqrt{-1}\,\partial\bar\partial F.
\]
This determines $F$ up to a constant on each connected component.
Choose these constants so that
\[
\int_{M_\alpha}e^F\omega_0^n
=\int_{M_\alpha}\omega_0^n
\qquad
\text{for every connected component }M_\alpha.
\]
For connected $M$, this is a single condition integrated over $M$.

This normalization is necessary because cohomologous K\"ahler forms
have the same total volume on each component:
\[
\int_{M_\alpha}\omega^n=\int_{M_\alpha}\omega_0^n.
\]
Once normalized, $F$ cannot be shifted by a nonzero constant
while retaining solvability of the volume equation.
Adding a constant to $\varphi$ leaves $\omega$ unchanged
and cannot compensate for such a shift in $F$.

\textbf{4. Why the Ricci equation becomes a determinant equation.}

The local determinant ratio is the globally defined ratio of volume
forms, so the change in the Ricci form can be written as
\[
\operatorname{Ric}(\omega)-\operatorname{Ric}(\omega_0)
=-\sqrt{-1}\,\partial\bar\partial
\log\frac{\omega^n}{\omega_0^n}.
\]

If $\operatorname{Ric}(\omega)=\rho$, this identity and the defining
equation for $F$ imply that
$\log(\omega^n/\omega_0^n)-F$ has vanishing
$\partial\bar\partial$. On a compact connected component this
function is constant; the volume normalization forces that constant
to be zero. Conversely, the volume equation below immediately gives
$\operatorname{Ric}(\omega)=\rho$ by the same identity. Hence
\[
\operatorname{Ric}(\omega)=\rho
\quad\Longleftrightarrow\quad
\left(
\omega_0+\sqrt{-1}\,\partial\bar\partial\varphi
\right)^n
=e^F\omega_0^n.
\]
Expanding the top wedge powers gives exactly the determinant equation
in the Local-Coordinate Formulation. On the positive-definite branch
specified there, this complex Monge--Amp\`ere equation is elliptic.

\textbf{5. What is unique.}\par\nopagebreak

The K\"ahler metric $\omega$ is unique. Its potential $\varphi$
is unique up to addition of a constant on each connected component.
One may fix this freedom by imposing
\[
\int_{M_\alpha}\varphi\,\omega_0^n=0
\qquad
\text{for every connected component }M_\alpha.
\]
In particular, if $c_1(M)_{\mathbb R}=0$ and one chooses $\rho=0$,
then every K\"ahler class contains a unique Ricci-flat K\"ahler metric.

\newpage
\section*{Definitions}

\begin{definitionbox}{위상공간 / Topological Space}
\label{source:S01}

A \emph{topological space} is a pair $(X,\mathcal T)$, where $X$ is a set
and $\mathcal T\subseteq\mathcal P(X)$ satisfies:

\begin{enumerate}[label=(\alph*),leftmargin=*,itemsep=.35em]

\item $\varnothing,X\in\mathcal T$;

\item any arbitrary union of members of $\mathcal T$ belongs to $\mathcal T$;

\item any finite intersection of members of $\mathcal T$ belongs to $\mathcal T$.

\end{enumerate}

The members of $\mathcal T$ are the \emph{open sets}.


\exampleheading
Let $S^2=\{(X,Y,Z)\in\R^3:X^2+Y^2+Z^2=1\}$. Its standard topology is $\mathcal T=\{S^2\cap V:V\subset\R^3\text{ open}\}$. Thus $\{Z>0\}\cap S^2$ is open, while the equator is closed. Arbitrary unions and finite intersections remain intersections with ambient open sets. This explicitly answers the supplied [PROMPT].

\textbf{연결 태그:} \hyperref[def:01-01]{\texttt{01-01} 위상공간}.
\end{definitionbox}
% Supplied PROMPT fulfilled explicitly in Example S01.




\begin{definitionbox}{거리와 거리 위상 / Metric and Metric Topology}
\label{source:S02}

Let \(X\) be a set. A \emph{metric} on \(X\) is a function
\[
d\colon X\times X\to\mathbb{R}
\]
such that, for all \(x,y,z\in X\),

\begin{enumerate}[label=(\alph*), leftmargin=*, itemsep=.35em]
    \item \(d(x,y)\geq 0\), and \(d(x,y)=0\) if and only if \(x=y\);
    \item \(d(x,y)=d(y,x)\);
    \item \(d(x,z)\leq d(x,y)+d(y,z)\).
\end{enumerate}

The pair \((X,d)\) is called a \emph{metric space}. 

For \(p\in X\) and
\(\varepsilon>0\), the \emph{open ball} centered at \(p\) with radius
\(\varepsilon\) is
\[
B_{\varepsilon}(p)
=
\{x\in X:d(x,p)<\varepsilon\}.
\]

A subset \(U\subseteq X\) is called \emph{open} if, for every \(p\in U\),
there exists \(\varepsilon>0\) such that
\[
B_{\varepsilon}(p)\subseteq U.
\]
The collection
\[
\mathcal{T}_d
=
\left\{
U\subseteq X:
\text{for every }p\in U,\text{ there exists }\varepsilon>0
\text{ such that }B_{\varepsilon}(p)\subseteq U
\right\}
\]
is a topology on \(X\), called the \emph{metric topology induced by \(d\)}.


\exampleheading
For $p,q\in S^2$, the chordal metric is $d(p,q)=\|p-q\|_{\R^3}$. Its balls are $B_r^{\R^3}(p)\cap S^2$, so its metric topology is exactly the subspace topology. For the poles N,S, d(N,S)=2. The intrinsic round distance is instead $d_g(p,q)=\arccos(p\cdot q)$, with $d_g(N,S)=\pi$; it induces the same topology, not the same numerical distance.

\textbf{연결 태그:} \hyperref[def:01-01]{\texttt{01-01} 위상공간}.
\end{definitionbox}

\begin{definitionbox}{콤팩트성 / Compactness}
\label{source:S03}
A topological space $X$ is \emph{compact} if every open cover
$X=\bigcup_{\alpha\in A}U_\alpha$ has a finite subcover:
\[
X=U_{\alpha_1}\cup\cdots\cup U_{\alpha_N}.
\]


\exampleheading
The subset $S^2\subset\R^3$ is closed, as the inverse image of $\{1\}$ under $X^2+Y^2+Z^2$, and bounded. Heine--Borel therefore makes it compact. In particular every relative open cover, not only the two stereographic charts, has a finite subcover.

\textbf{연결 태그:} \hyperref[def:12-03]{\texttt{12-03} Paracompact 공간}.
\end{definitionbox}

\begin{definitionbox}{파라콤팩트성 / Paracompactness}
\label{source:S04}
A topological space $X$ is \emph{paracompact} if every open cover of $X$ admits a locally finite open refinement.
An open cover $\{V_\beta\}_{\beta\in B}$ \emph{refines} an open cover
$\{U_\alpha\}_{\alpha\in A}$ if for every $\beta\in B$, there exists
$\alpha\in A$ such that $V_\beta \subseteq U_\alpha.$

The collection $\{V_\beta\}_{\beta\in B}$ is \emph{locally finite} if every point $x\in X$ has a neighborhood that intersects only finitely many of the sets $V_\beta$.
Some authors include Hausdorffness in the definition of paracompactness. In our setting, manifolds are assumed to be Hausdorff.


\exampleheading
The unit sphere is compact because it is closed and bounded in $\R^3$ (Heine--Borel). Every open cover therefore has a finite subcover, which is itself a locally finite open refinement. This verifies paracompactness of this particular $S^2$.

\textbf{연결 태그:} \hyperref[def:12-03]{\texttt{12-03} Paracompact 공간}.
\end{definitionbox}


\begin{definitionbox}{단위분할 / Partition of unity}
\label{source:S05}
Let $\{U_\alpha\}_{\alpha\in A}$ be an open cover of a smooth manifold $M$.
A \emph{smooth partition of unity subordinate to this cover} is a family
$\{\varphi_\alpha\}_{\alpha\in A}$ of smooth functions
$\varphi_\alpha:M\to[0,1]$ such that
\[
\supp\varphi_\alpha\subseteq U_\alpha,
\qquad
\sum_{\alpha\in A}\varphi_\alpha(x)=1\quad(x\in M),
\]
and the family $\{\supp\varphi_\alpha\}$ is locally finite. Here
\[
\supp\varphi_\alpha
=\overline{\{x\in M:\varphi_\alpha(x)\ne0\}}.
\]
Local finiteness ensures that the sum is locally a finite sum of smooth
functions. Zero functions are allowed.


\exampleheading
On $S^2$, put $\rho_0=\eta(Z+1/2)/(\eta(Z+1/2)+\eta(1/2-Z))$ and $\rho_\infty=1-\rho_0$. The denominator is everywhere positive. These smooth functions sum to 1, with supports in $\{Z\ge-1/2\}\subset U_0$ and $\{Z\le1/2\}\subset U_\infty$. This is an explicit partition subordinate to the two standard charts.

\textbf{연결 태그:} \hyperref[def:12-05]{\texttt{12-05} 매끄러운 단위분할}.
\end{definitionbox}

\begin{definitionbox}{위상다양체 / Topological manifold}
\label{source:S06}
An $m$-dimensional \emph{topological manifold} is a Hausdorff,
second-countable topological space $M$ in which each point has an open
neighborhood homeomorphic to an open subset of $\R^m$.
Such a homeomorphism
\[
\varphi:U\longrightarrow\varphi(U)\subseteq\R^m
\]
defines a \emph{coordinate chart} $(U,\varphi)$.
Hausdorff means that any two distinct points have disjoint open
neighborhoods; second-countable means that the topology has a countable basis.


\exampleheading
On $S^2\setminus\{S\}$, $z=(X+iY)/(1+Z)$ is a homeomorphism onto $\C\simeq\R^2$ with inverse $\Phi(x+iy)=(2x,2y,1-x^2-y^2)/(1+x^2+y^2)$. Together with the w-chart at S it verifies local Euclidean dimension two. The subspace topology is Hausdorff and second countable.

\textbf{연결 태그:} \hyperref[def:06-04]{\texttt{06-04} 매끄러운 다양체}.
\end{definitionbox}

\begin{definitionbox}{매끄러운 다양체 / Smooth manifold}
\label{source:S07}
A \emph{smooth manifold} is a topological manifold equipped with a smooth
structure. A smooth structure can be specified by an atlas of charts
$\{(U_\alpha,\varphi_\alpha)\}$ covering $M$, whose transition maps
\[
\varphi_\beta\circ\varphi_\alpha^{-1}:
\varphi_\alpha(U_\alpha\cap U_\beta)
\longrightarrow\varphi_\beta(U_\alpha\cap U_\beta)
\]
are $C^\infty$. Equivalently, one specifies a maximal smooth atlas.
Two smooth atlases define the same smooth structure if their union is a
smooth atlas. Thus an atlas need not itself be maximal to specify a
smooth manifold.


\exampleheading
The fixed charts $z=x+iy$ and $w=1/z$ identify the underlying $S^2$ with a smooth real surface. On the overlap, $(x,y)\mapsto(x/(x^2+y^2),-y/(x^2+y^2))$ is smooth with smooth inverse. The Hausdorff and countable-basis conditions follow from the subspace topology in $\R^3$.

\textbf{연결 태그:} \hyperref[def:06-04]{\texttt{06-04} 매끄러운 다양체}.
\end{definitionbox}

\newpage
\begin{definitionbox}{리만다양체 / Riemannian manifold}
\label{source:S08}
A \emph{Riemannian manifold} is a pair $(M,g)$ consisting of a smooth
manifold and a Riemannian metric.

\exampleheading
The pair $(S^2,g)$ with $g_p(u,v)=u\cdot v$ for $u,v\in T_pS^2=p^\perp$ is a Riemannian manifold. In the fixed z-chart, $g=4(1+|z|^2)^{-2}(dx^2+dy^2)$. Here juxtaposition denotes symmetric metric tensors, not wedge products.

\textbf{연결 태그:} \hyperref[def:06-04]{\texttt{06-04} 매끄러운 다양체}.
\end{definitionbox}

\begin{subdefinitionbox}{리만계량 / Riemannian metric}
\label{source:S09}
A Riemannian metric is a smooth section $g$ of
$\operatorname{Sym}^2(T^*M)$ such that at every point
\[
g_p:T_pM\times T_pM\to\R
\]
is symmetric, bilinear, and positive-definite. In particular,
\[
g_p(X,Y)=g_p(Y,X),\qquad g_p(X,X)>0\quad\text{for }X\ne0.
\]


\exampleheading
On $S^2$, restrict the Euclidean dot product to each tangent plane. It is symmetric and bilinear, and $g_p(u,u)=\|u\|^2>0$ for $u\ne0$. Its coordinate matrix in the z-chart is $4(1+|z|^2)^{-2}I_2$, a smooth positive-definite matrix. This checks the smooth family condition as well as pointwise positivity.

\textbf{연결 태그:} \hyperref[def:06-06]{\texttt{06-06} 여접공간}.
\end{subdefinitionbox}

\begin{subdefinitionbox}{Levi--Civita 연결 / Levi--Civita connection}
\label{source:S10}
A connection on $TM$ is an $\R$-bilinear operation $(X,Y)\mapsto\nabla_XY$
on smooth vector fields satisfying
\[
\nabla_{fX}Y=f\nabla_XY,\qquad
\nabla_X(fY)=X(f)Y+f\nabla_XY.
\]
The \emph{Levi--Civita connection} is the unique connection with
\begin{align*}
\nabla_XY-\nabla_YX&=[X,Y], &&\text{(torsion-free)}\\
X\bigl(g(Y,Z)\bigr)&=g(\nabla_XY,Z)+g(Y,\nabla_XZ),
&&\text{(metric compatible)}
\end{align*}
where $[X,Y]$ is the Lie bracket, defined by
$[X,Y](f)=X(Y(f))-Y(X(f))$.


\exampleheading
For tangent vector fields U,V on the unit sphere, $\nabla_UV=P_p(D_UV)=D_UV+\langle U,V\rangle p$, where $P_p(a)=a-\langle a,p\rangle p$ and D is the ambient Euclidean derivative. Differentiating $\langle V,p\rangle=0$ gives $\langle D_UV,p\rangle=-\langle V,U\rangle$. Projection makes the output tangent, and the Euclidean product rule verifies metric compatibility. The normal correction is symmetric in U,V, so torsion is zero.

\textbf{연결 태그:} \hyperref[def:N29]{\texttt{N29} 복소 connection}.
\end{subdefinitionbox}

\begin{subdefinitionbox}{실수 접공간과 접다발 / Real tangent space and tangent bundle}
\label{source:S11}
A tangent vector at $p$ can be defined as an $\R$-linear map
$v:C^\infty(M,\R)\to\R$ satisfying
\[
v(fh)=f(p)v(h)+h(p)v(f).
\]
These derivations form the real vector space $T_pM$ of dimension $m$.
The tangent bundle is
\[
TM=\bigsqcup_{p\in M}T_pM,
\qquad \pi:TM\to M,\quad v\in T_pM\longmapsto p,
\]
with the smooth vector-bundle structure induced by coordinate charts.

\exampleheading
At $p\in S^2$, $T_pS^2=p^\perp$. If $v\in p^\perp$ and $\widetilde f$ extends a smooth sphere function near p, the derivation is $v(f)=D\widetilde f_p(v)$, independent of the extension. At N, $v=e_1$ takes the coordinate restrictions X,Y,Z to 1,0,0, respectively. The total tangent bundle consists of pairs $(p,v)$ with $\|p\|=1$ and $p\cdot v=0$.

\textbf{연결 태그:} \hyperref[def:06-05]{\texttt{06-05} 접벡터와 접공간}.
\end{subdefinitionbox}

\newpage
\begin{definitionbox}{복소다양체와 복소구조 / Complex manifold and complex structure}
\label{source:S12}
A \emph{complex manifold} of complex dimension $n$ is a Hausdorff,
second-countable space with an atlas of homeomorphisms
\[
\varphi_\alpha:U_\alpha\longrightarrow
\varphi_\alpha(U_\alpha)\subseteq\C^n
\]
onto open sets, covering $M$, with holomorphic transition maps.
The underlying smooth manifold has real dimension $2n$.

An \emph{almost complex structure} is a smooth real vector-bundle
endomorphism $J:TM\to TM$ with $J^2=-\id$.
It is \emph{integrable} if there is a complex atlas inducing $J$;
then it is called a \emph{complex structure}.
In holomorphic coordinates $z^j=x^j+iy^j$,
\[
J\frac{\partial}{\partial x^j}=\frac{\partial}{\partial y^j},
\qquad
J\frac{\partial}{\partial y^j}=-\frac{\partial}{\partial x^j}.
\]


\exampleheading
On the oriented unit sphere put $J_pv=p\times v$ for $v\perp p$. The vector triple-product identity gives $J_p^2v=-v$. In our z-chart, $J\partial_x=\partial_y$, and the second chart satisfies w=1/z. Thus this J is induced by a complex atlas, not merely an almost complex structure.

\textbf{연결 태그:} \hyperref[def:B04]{\texttt{B04} 복소다양체}.
\end{definitionbox}

\begin{definitionbox}{텐서곱 / Tensor product}
\label{source:S13}
For vector spaces $V,W$ over a field $\mathbb F$, their tensor product
is a vector space $V\otimes_{\mathbb F}W$ with a bilinear map
$(v,w)\mapsto v\otimes w$ such that every bilinear map
$B:V\times W\to U$ into an $\mathbb F$-vector space factors uniquely as
\[
B(v,w)=\widetilde B(v\otimes w)
\]
for a linear map $\widetilde B:V\otimes_{\mathbb F}W\to U$.
This universal property determines the tensor product up to a unique
isomorphism respecting the bilinear map.


\exampleheading
At $N\in S^2$, take the tangent basis $e_1,e_2$ and dual basis $\epsilon^1,\epsilon^2$. Then $(\epsilon^1\otimes\epsilon^2)(e_1,e_2)=1$, while $(\epsilon^1\wedge\epsilon^2)(e_1,e_2)=1$ and its value on $(e_2,e_1)$ is $-1$. The first tensor is not alternating, and the dual covector satisfies $\epsilon^1(e_1)=1$.

\textbf{연결 태그:} \hyperref[def:B02]{\texttt{B02} 쌍대·텐서곱·외대수}.
\end{definitionbox}

\begin{definitionbox}{복소화 접다발 / Complexified tangent bundle}
\label{source:S14}
The \emph{complexified tangent bundle} is
\[
T_{\C}M=TM\otimes_{\R}\C,
\qquad (T_{\C}M)_p=T_pM\otimes_{\R}\C.
\]
If $M$ has real dimension $m$, its complexification has complex rank $m$.
An almost complex structure extends $\C$-linearly to $T_{\C}M$ and gives
\[
T_{\C}M=T^{1,0}M\oplus T^{0,1}M.
\]
When $\dim_{\C}M=n$, each summand has complex rank $n$, whereas
$T_{\C}M$ has complex rank $2n$.


\exampleheading
At N on $S^2$, $T_NS^2=\R e_1\oplus\R e_2$, so its complexification is $\C e_1\oplus\C e_2$, of complex dimension two. The extended J has eigenvectors $e_1-i e_2$ and $e_1+i e_2$, with eigenvalues i and -i. These are the one-dimensional (1,0) and (0,1) spaces. Complexification and taking the holomorphic tangent line are different operations.

\textbf{연결 태그:} \hyperref[def:B07]{\texttt{B07} 정칙 접다발과 여접다발}.
\end{definitionbox}

\newpage
\begin{definitionbox}{정칙 접다발 / Holomorphic tangent bundle}
\label{source:S15}
For a complex manifold $(M,J)$, define the eigenbundles
\[
T^{1,0}M=\ker(J-i\id),\qquad T^{0,1}M=\ker(J+i\id)
\]
inside $T_{\C}M$. In holomorphic coordinates,
\begin{align*}
\frac{\partial}{\partial z^j}
&=\frac12\left(\frac{\partial}{\partial x^j}
-i\frac{\partial}{\partial y^j}\right),\\
\frac{\partial}{\partial\bar z^j}
&=\frac12\left(\frac{\partial}{\partial x^j}
+i\frac{\partial}{\partial y^j}\right)
\end{align*}
form local frames for $T^{1,0}M$ and $T^{0,1}M$, respectively.
The holomorphic coordinate frames give $T^{1,0}M$ its natural
\emph{holomorphic vector-bundle structure}. A smooth section
$\sum_j a^j\partial/\partial z^j$ is a holomorphic section precisely when
the functions $a^j$ are holomorphic. Thus not every smooth section of the
holomorphic tangent bundle is holomorphic.


\exampleheading
On $\CP^1$, $\partial_w=-z^2\partial_z$ and $dw=-z^{-2}dz$. Pairing gives $dw(\partial_w)=1$ on the overlap. Thus the tangent and cotangent transition functions are reciprocal, and each bundle has complex rank one.

\textbf{연결 태그:} \hyperref[def:B07]{\texttt{B07} 정칙 접다발과 여접다발}.
\end{definitionbox}

\begin{definitionbox}{Chern 연결 / Chern connection}
\label{source:S16}
A holomorphic vector bundle $E\to M$ has local trivializations
$E|_U\simeq U\times\C^r$ with holomorphic, fiberwise linear transition maps.
A \emph{Hermitian metric} $h$ on $E$ is a smooth family of positive-definite
Hermitian inner products on its fibers. We take $h$ to be complex-linear
in its first argument and conjugate-linear in its second.

The \emph{Chern connection} is the unique complex connection
\[
\nabla:\Gamma(E)\longrightarrow\Omega^1(M;E)
\]
such that $\nabla h=0$ and $\nabla^{0,1}=\bar\partial_E$.
Here $\Gamma(E)$ denotes smooth sections, $\Omega^1(M;E)$ denotes
$E$-valued one-forms, and $\bar\partial_E$ applies $\bar\partial$ to the
coefficient functions in a local holomorphic frame.
A complex connection satisfies $\nabla(fs)=df\otimes s+f\nabla s$.
For real vector fields $X$, compatibility with $h$ means
\[
X(h(s,t))=h(\nabla_Xs,t)+h(s,\nabla_Xt).
\]
We use the curvature convention
\[
\Theta=\nabla^2,\qquad
\Theta(X,Y)=\nabla_X\nabla_Y-\nabla_Y\nabla_X-\nabla_{[X,Y]}.
\]
The curvature of a Chern connection is of type $(1,1)$.


\exampleheading
The Hermitian metric $h_z=(1+|z|^2)^2/2$ on $K_{\CP^1}$ gives the Chern connection $A_z=\partial\log h_z=2\bar z(1+|z|^2)^{-1}dz$. Its curvature is $F_K=-2(1+|z|^2)^{-2}dz\wedge d\bar z$, a global (1,1)-form. The factor $1/2$ in h disappears when differentiating its logarithm.

\textbf{연결 태그:} \hyperref[def:N30]{\texttt{N30} Chern connection과 곡률}.
\end{definitionbox}

\begin{definitionbox}{심플렉틱 형식 / Symplectic form}
\label{source:S17}
A \emph{symplectic form} is a real differential two-form
$\omega\in\Omega^2(M;\R)$ such that $d\omega=0$ and $\omega$ is
nondegenerate:
\[
\omega_p(X,Y)=0\ \text{for every }Y\in T_pM
\quad\Longrightarrow\quad X=0.
\]
A pair $(M,\omega)$ is then a \emph{symplectic manifold}.


\exampleheading
On $S^2$, $\Omega_p(u,v)=p\cdot(u\times v)$ is closed for dimension reasons. If $u\ne0$, choose $v=J_pu=p\times u$; then $\Omega_p(u,J_pu)=\|u\|^2>0$. Hence it is nondegenerate and gives a symplectic form.

\textbf{연결 태그:} \hyperref[def:06-07]{\texttt{06-07} 미분 q-형식}.
\end{definitionbox}

\begin{definitionbox}{호환성·Hermitian 계량·켈러 다양체 / Compatibility, Hermitian metrics, and K\"ahler manifolds}
\label{source:S18}
A Riemannian metric $g$ on an almost complex manifold $(M,J)$ is
\emph{$J$-compatible} if
\[
g(JX,JY)=g(X,Y).
\]
Its associated real two-form, with our convention, is
\[
\theoremboxed{\omega(X,Y)=g(JX,Y).}
\]
It is alternating because $g(JX,Y)=-g(X,JY)$, and
\[
\omega(X,JX)=g(X,X)>0\quad(X\ne0).
\]
If $J$ is integrable, $g$ is called a \emph{Hermitian metric} on $(M,J)$.
If, additionally, $d\omega=0$, the tuple $(M,J,g,\omega)$ is a
\emph{K\"ahler manifold}, and $\omega$ is its \emph{K\"ahler form}.
Equivalently, on a complex manifold a K\"ahler form is a closed positive
real $(1,1)$-form; see Definition 20 for the type decomposition.


\exampleheading
For the unit sphere, take $g_p(u,v)=u\cdot v$ and $J_pu=p\times u$. Then $g(Ju,Jv)=g(u,v)$ and $g(Ju,v)=\Omega(u,v)$. The chosen complex atlas makes J integrable, while $d\Omega=0$, so this gives a K\"ahler structure on $\CP^1$. In the z-chart, $\Omega=2i(1+|z|^2)^{-2}dz\wedge d\bar z$ and $g_{z\bar z}=2(1+|z|^2)^{-2}$.

\textbf{연결 태그:} \hyperref[def:N28]{\texttt{N28} 선다발의 Hermitian 계량}.
\end{definitionbox}

\begin{definitionbox}{리만 곡률 텐서 / Riemann curvature tensor}
\label{source:S19}
For the Levi--Civita connection we define
\[
R(X,Y)Z=\nabla_X\nabla_YZ-\nabla_Y\nabla_XZ-\nabla_{[X,Y]}Z.
\]
This is a tensor measuring the failure of covariant derivatives to
commute, with the Lie-bracket correction included.


\exampleheading
For the unit round $S^2$ with the stated sign convention, $R(U,V)W=\langle V,W\rangle U-\langle U,W\rangle V$. For an oriented orthonormal pair $e_1,e_2$ at p, $R(e_1,e_2)e_2=e_1$, and the Gaussian curvature is $g(R(e_1,e_2)e_2,e_1)=1$. This is tangent-bundle curvature, not the canonical-line curvature form.

\textbf{연결 태그:} \hyperref[def:B09]{\texttt{B09} 선다발 연결의 곡률}.
\end{definitionbox}

\begin{definitionbox}{Ricci 곡률 텐서 / Ricci curvature tensor}
\label{source:S20}
The \emph{Ricci tensor} is the contraction
\[
\Ric_g(X,Y)=\operatorname{tr}\bigl(Z\longmapsto R(Z,X)Y\bigr).
\]
If $\{e_1,\ldots,e_m\}$ is a real $g$-orthonormal basis of $T_pM$, then
\[
\theoremboxed{
\begin{gathered}
\Ric_g(X,Y)=\sum_{a=1}^{m}g(R(e_a,X)Y,e_a),\\[4pt]
m=\dim_{\R}M.
\end{gathered}
}
\]
In particular, this sum has $2n$ terms on a complex $n$-dimensional
manifold. The Ricci tensor is a symmetric real $(0,2)$-tensor; the
Ricci form in Definition 21 is an alternating two-form.


\exampleheading
On the round $S^2$, contract the preceding R with an orthonormal pair. The term for $e_1$ in $\Ric_g(e_1,e_1)$ is zero; the term for $e_2$ is $g(R(e_2,e_1)e_1,e_2)=1$. Likewise $\Ric_g(e_2,e_2)=1$ and the mixed term vanishes, so $\Ric_g=g$.

\textbf{연결 태그:} \hyperref[def:B09]{\texttt{B09} 선다발 연결의 곡률}.
\end{definitionbox}

\newpage
\begin{definitionbox}{미분형식 / Differential forms}
\label{source:S21}
A real differential $k$-form is a smooth section of $\Lambda^kT^*M$:
\[
\Omega^k(M;\R)=\Gamma(\Lambda^kT^*M).
\]
At $p$, its value is an alternating $k$-linear map $(T_pM)^k\to\R$.
Complex-valued forms instead lie in
\[
\Omega^k(M;\C)=\Gamma(\Lambda^kT^*M\otimes_{\R}\C).
\]
We abbreviate $\Omega^k(M;\R)$ to $\Omega^k(M)$ when there is no ambiguity.
The wedge product $\alpha\wedge\beta$ is the exterior product of forms;
it has degree $k+\ell$ when $\alpha$ has degree $k$ and $\beta$ degree $\ell$.


\exampleheading
The area form of the unit sphere is $\Omega_p(u,v)=p\cdot(u\times v)$ for $u,v\in T_pS^2$. This is alternating and depends smoothly on p. In the z-chart it is $\Omega=4(1+|z|^2)^{-2}\,dx\wedge dy$, so its coordinate coefficient is explicitly smooth.

\textbf{연결 태그:} \hyperref[def:06-07]{\texttt{06-07} 미분 q-형식}.
\end{definitionbox}

\begin{definitionbox}{외미분과 드람 코호몰로지 / Exterior derivative and de Rham cohomology}
\label{source:S22}
The exterior derivative consists of $\R$-linear maps
$d:\Omega^k(M)\to\Omega^{k+1}(M)$ characterized by
\begin{align*}
df(X)&=X(f),\\
d^2&=0,\\
d(\alpha\wedge\beta)&=d\alpha\wedge\beta+(-1)^k\alpha\wedge d\beta,
\\[-2pt]
&\qquad \alpha\in\Omega^k(M).
\end{align*}
It extends complex-linearly to complex-valued forms.
A form $\alpha$ is \emph{closed} if $d\alpha=0$ and \emph{exact} if
$\alpha=d\beta$ for some form $\beta$ of one lower degree.
The de Rham cohomology is
\[
H^k_{\mathrm{dR}}(M;\R)=
\frac{\ker(d:\Omega^k(M)\to\Omega^{k+1}(M))}
{\operatorname{im}(d:\Omega^{k-1}(M)\to\Omega^k(M))}.
\]
Set $\Omega^{-1}(M)=0$ for the case $k=0$.
For a closed form $\alpha$, its cohomology class is denoted $[\alpha]$.
Closed forms $\alpha,\beta$ have the same class exactly when
$\alpha-\beta$ is exact.


\exampleheading
Integration induces $H^2_{\mathrm{dR}}(S^2;\R)\simeq\R$, $[\alpha]\mapsto\int_{S^2}\alpha$. The class $[\Omega/(4\pi)]$ maps to 1. In degrees 0 and 1 the groups are $\R$ and 0, respectively. These are the computed groups of this sphere, not part of the general definition. The global height function on $S^2$ is $H=(1-x^2-y^2)/(1+x^2+y^2)$ in the z-chart. Its differential is $dH=-4(x\,dx+y\,dy)/(1+x^2+y^2)^2$. Applying d again gives zero: the mixed partials cancel in the coefficient of $dx\wedge dy$.

\textbf{연결 태그:} \hyperref[def:06-14]{\texttt{06-14} de Rham 코호몰로지}.
\end{definitionbox}

\begin{definitionbox}{형식차수와 Dolbeault 분해 / Forms of type $(p,q)$ and $d=\partial+\bar\partial$}
\label{source:S23}
On a complex manifold, forms of type $(p,q)$ are locally sums of terms
\[
a_{I\bar J}(z)\,dz^{i_1}\wedge\cdots\wedge dz^{i_p}
\wedge d\bar z^{j_1}\wedge\cdots\wedge d\bar z^{j_q},
\]
with smooth complex-valued coefficients. Their space is
$\Omega^{p,q}(M)$, and
\[
\Omega^k(M;\C)=\bigoplus_{p+q=k}\Omega^{p,q}(M).
\]
The exterior derivative decomposes as $d=\partial+\bar\partial$, where
\[
\partial:\Omega^{p,q}\to\Omega^{p+1,q},\qquad
\bar\partial:\Omega^{p,q}\to\Omega^{p,q+1}.
\]
For a function $f$, the operators are
\[
\partial f=\sum_j\frac{\partial f}{\partial z^j}dz^j,
\qquad \bar\partial f=\sum_j\frac{\partial f}{\partial\bar z^j}d\bar z^j.
\]
They satisfy $\partial^2=\bar\partial^2=0$ and
$\partial\bar\partial+\bar\partial\partial=0$.
The notation $\partial\bar\partial f$ means $\partial(\bar\partial f)$.
This two-term decomposition of $d$ uses integrability; an arbitrary
almost complex structure need not have it.


\exampleheading
On the z-chart of $\CP^1$, take $f=z\bar z$. Then $\partial f=\bar z\,dz$, $\bar\partial f=z\,d\bar z$, and $\partial\bar\partial f=dz\wedge d\bar z$. These have types $(1,0),(0,1),(1,1)$ respectively. The function f is used only on this chart, not as a global function on the sphere.

\textbf{연결 태그:} \hyperref[def:B06]{\texttt{B06} 형식의 형식차수와 Dolbeault 연산자}.
\end{definitionbox}

\begin{definitionbox}{Ricci 실수 (1,1)-형식 / Ricci real (1,1)-form}
\label{source:S24}
For a K\"ahler metric, let $g_{\C}$ be the complex-bilinear extension of
the real metric and set
\[
g_{i\bar j}=g_{\C}\left(\frac{\partial}{\partial z^i},
\frac{\partial}{\partial\bar z^j}\right).
\]
Then, consistently with Definition 15,
\[
\omega=i\sum_{i,j=1}^{n}g_{i\bar j}\,dz^i\wedge d\bar z^j.
\]
The \emph{Ricci form} is
\[
\theoremboxed{\rho_\omega=-i\ddbar\log\det(g_{i\bar j}).}
\]
Equivalently,
\begin{align*}
\rho_\omega&=-i\sum_{i,j=1}^{n}
\frac{\partial^2\log\det(g_{k\bar\ell})}{\partial z^i\partial\bar z^j}
\,dz^i\wedge d\bar z^j,\\[4pt]
\rho_\omega(X,Y)&=\Ric_g(JX,Y).
\end{align*}
These local expressions define a global closed real $(1,1)$-form.
The Ricci form is not necessarily positive.
Some authors denote this form by $\Ric(\omega)$; that notation refers
to the two-form, whereas $\Ric_g$ above denotes the symmetric tensor.


\exampleheading
For the round metric on $\CP^1$, $g_{z\bar z}=2(1+|z|^2)^{-2}$. Hence $\rho=-i\partial\bar\partial\log g_{z\bar z}=2i(1+|z|^2)^{-2}dz\wedge d\bar z=\Omega$. Its integral is $4\pi$, and $\rho(U,V)=\Ric_g(JU,V)=g(JU,V)$. The tensor $\Ric_g=g$ is symmetric, whereas $\rho=\Omega$ is alternating.

\textbf{연결 태그:} \hyperref[def:B11]{\texttt{B11} 첫째 Chern 형식}.
\end{definitionbox}

\newpage
\begin{definitionbox}{첫째 Chern 류 / First Chern class}
\label{source:S25}
For a smooth complex vector bundle $E\to M$, the first Chern class is
an integral characteristic class
\[
c_1(E)\in H^2(M;\mathbb Z),\qquad c_1(E)=c_1(\det E).
\]
Write $c_1(E)_{\R}$ for its image under the natural map to
$H^2_{\mathrm{dR}}(M;\R)$. Choose a Hermitian metric on $E$ and a
connection compatible with that metric. With $\Theta=\nabla^2$ as in
Definition 13, Chern--Weil theory gives
\[
\theoremboxed{c_1(E)_{\R}
=\left[\frac{i}{2\pi}\operatorname{Tr}\Theta\right].}
\]
The displayed form is real and closed, and its cohomology class is
independent of the chosen Hermitian metric and compatible connection.
If $E$ is holomorphic, one may use its Chern connection; then the form
also has type $(1,1)$.
Requiring only that a connection be complex-linear would not guarantee
that this particular representative is a real-valued form.
The curvature formula determines the real image of $c_1(E)$; it does
not detect any torsion in the integral class. See \cite{demailly}.

For a complex manifold, $c_1(M):=c_1(T^{1,0}M)$. On a K\"ahler manifold,
\[
\theoremboxed{\left[\frac{\rho_\omega}{2\pi}\right]=\cR,
\qquad [\rho_\omega]=2\pi\cR.}
\]
These equalities use real de Rham cohomology. With the definition of
$\rho_\omega$ in Definition 21, the factor $2\pi$ must be retained.
If one defines a different, normalized form
$\widehat\rho_\omega=\rho_\omega/(2\pi)$, then
$[\widehat\rho_\omega]=\cR$.


\exampleheading
Let u be the integral generator of $H^2(\CP^1;\Z)$ evaluating to 1 on the complex-oriented sphere. Then $c_1(T^{1,0}\CP^1)=2u$ and $c_1(K_{\CP^1})=-2u$. Their real images are represented by $\Omega/(2\pi)$ and $-\Omega/(2\pi)$, respectively. Here there is no torsion in $H^2$, so the integral class is determined by this integer pairing; this is special to this example.

\textbf{연결 태그:} \hyperref[def:B12]{\texttt{B12} 실수 첫째 Chern 류}.
\end{definitionbox}

\clearpage
\section{통합 정의 체계 / Complete tagged definition atlas}
각 태그는 웹 정의집과 같습니다. 본문은 원래 정의를 보존하고, 예시를 덧붙였습니다. B01--B12도 모두 포함합니다. 기존 TeX와 겹치는 정의도 이 부에서는 태그로 다시 찾을 수 있습니다.

\subsection{공간과 위상 / Spaces and topology}

\begin{atlasdefinition}{\texttt{01-01}: 위상공간 / Topological space}
\label{def:01-01}
집합 \(X\)와 부분집합들의 모임 \(\tau\)의 쌍으로, \(\varnothing,X\in\tau\)이고 \(\tau\)는 임의 합집합과 유한 교집합에 닫혀 있다. \(\tau\)의 원소를 열린집합 (open set)이라 한다.

\textbf{조건 확인 / Conditions.} 교집합은 유한 개, 합집합은 임의 개다.

\textbf{직접 선행 / Prerequisites.} 집합·함수·실수·복소수 등 공통 배경.

\exampleheading
Let $S^2=\{(X,Y,Z)\in\R^3:X^2+Y^2+Z^2=1\}$. Its standard topology is $\mathcal T=\{S^2\cap V:V\subset\R^3\text{ open}\}$. Thus $\{Z>0\}\cap S^2$ is open, while the equator is closed. Arbitrary unions and finite intersections remain intersections with ambient open sets. This explicitly answers the supplied [PROMPT].
\end{atlasdefinition}

\begin{atlasdefinition}{\texttt{01-03}: 열린덮개 / Open cover}
\label{def:01-03}
위상공간 (topological space) \(X\)의 열린덮개 (open cover)는 각 \(U_i\subset X\)가 열려 있고 \(X=\bigcup_{i\in I}U_i\)인 집합족 \(\{U_i\}_{i\in I}\)다.

\textbf{조건 확인 / Conditions.} 합집합이 공간 전체여야 한다.

\textbf{직접 선행 / Prerequisites.} \hyperref[def:01-01]{\texttt{01-01} 위상공간}

\exampleheading
The two open sets $U_0=S^2\setminus\{S\}$ and $U_\infty=S^2\setminus\{N\}$ cover $S^2$. Under the fixed identification with $\CP^1$, these are precisely the $z$- and $w$-charts, with overlap $\C^\times$ and $w=1/z$.
\end{atlasdefinition}

\begin{atlasdefinition}{\texttt{01-02}: 열린 근방 / Open neighborhood}
\label{def:01-02}
X를 위상공간 (topological space)이라 하자. \(x\in X\)의 열린 근방 (open neighborhood)은 \(x\in U\)를 만족하는 열린집합 (open set) \(U\subset X\)다.

\textbf{조건 확인 / Conditions.} 점을 포함하고 열려 있어야 한다.

\textbf{직접 선행 / Prerequisites.} \hyperref[def:01-01]{\texttt{01-01} 위상공간}

\exampleheading
At $N=(0,0,1)\in S^2$, the cap $U=S^2\cap\{Z>1/2\}$ is an open neighbourhood: $N\in U$, and $\{Z>1/2\}$ is open in $\R^3$. The equator is not a neighbourhood of $N$.
\end{atlasdefinition}

\begin{atlasdefinition}{\texttt{06-01}: 연속사상 / Continuous map}
\label{def:06-01}
위상공간 (topological space) 사이의 함수 \(f:X\to Y\)가 연속이라는 것은 모든 열린 \(V\subset Y\)에 대해 \(f^{-1}(V)\subset X\)가 열린 것이라는 뜻이다.

\textbf{조건 확인 / Conditions.} 상이 아니라 역상 조건이다.

\textbf{직접 선행 / Prerequisites.} \hyperref[def:01-01]{\texttt{01-01} 위상공간}

\exampleheading
The height function $H:S^2\to\R$, $H(X,Y,Z)=Z$, is continuous. For example $H^{-1}((0,2))=S^2\cap\{Z>0\}$ is open. More generally every inverse image of an open interval is an ambient open slab intersected with $S^2$.
\end{atlasdefinition}

\begin{atlasdefinition}{\texttt{06-02}: Hausdorff 공간 / Hausdorff space}
\label{def:06-02}
모든 \(x\ne y\)에 대해 \(x\in U\), \(y\in V\), \(U\cap V=\varnothing\)인 열린 근방 (open neighborhood) \(U,V\)가 존재하는 위상공간 (topological space)이다.

\textbf{조건 확인 / Conditions.} 서로 다른 모든 두 점에 대해 요구한다.

\textbf{직접 선행 / Prerequisites.} \hyperref[def:01-02]{\texttt{01-02} 열린 근방}

\exampleheading
For distinct $p,q\in S^2$, put $r=\|p-q\|/3$. The relative open balls $B_r^{\R^3}(p)\cap S^2$ and $B_r^{\R^3}(q)\cap S^2$ are disjoint neighbourhoods, by the triangle inequality. Thus the subspace topology of $S^2$ is Hausdorff.
\end{atlasdefinition}

\begin{atlasdefinition}{\texttt{06-03}: 열린 기저 / Basis for a topology}
\label{def:06-03}
위상공간 (topological space) X를 고정한다. 열린집합 (open set)들의 모임 \(\mathcal B\)가 기저라는 것은 모든 열린집합이 \(\mathcal B\) 원소들의 합집합으로 표현된다는 뜻이다. 가산 열린 기저는 이 모임을 가산으로 잡을 수 있다는 조건이다.

\textbf{조건 확인 / Conditions.} 기저 원소 자체도 열려 있어야 한다.

\textbf{직접 선행 / Prerequisites.} \hyperref[def:01-01]{\texttt{01-01} 위상공간}

\exampleheading
The nonempty sets $S^2\cap B_r(a)$, with $a\in\Q^3$ and $r\in\Q_{>0}$, form a countable basis. Given $p\in U$ open in $S^2$, choose a small ambient ball around $p$ inside the ambient open set defining $U$, then a rational-centred, rational-radius ball containing $p$ inside it.
\end{atlasdefinition}

\begin{atlasdefinition}{\texttt{01-04}: 닫힌집합 / Closed set}
\label{def:01-04}
X를 위상공간 (topological space)이라 하자. \(A\subset X\)가 닫혔다는 것은 \(X\setminus A\)가 열린집합 (open set)이라는 뜻이다.

\textbf{조건 확인 / Conditions.} 여집합은 주어진 공간 X 안에서 취한다.

\textbf{직접 선행 / Prerequisites.} \hyperref[def:01-01]{\texttt{01-01} 위상공간}

\exampleheading
The equator $E=S^2\cap\{Z=0\}$ is closed in $S^2$: its complement is the union of the two open hemispheres $\{Z>0\}$ and $\{Z<0\}$, intersected with $S^2$.
\end{atlasdefinition}

\begin{atlasdefinition}{\texttt{01-05}: 폐포 / Closure}
\label{def:01-05}
X를 위상공간 (topological space)이라 하자. \(A\subset X\)의 폐포 (closure) \(\overline A\)는 \(A\)를 포함하는 모든 닫힌집합 (closed set)의 교집합이다.

\textbf{조건 확인 / Conditions.} A를 포함하는 모든 닫힌집합 (closed set)을 사용한다.

\textbf{직접 선행 / Prerequisites.} \hyperref[def:01-04]{\texttt{01-04} 닫힌집합}

\exampleheading
For the open northern hemisphere $H=\{p\in S^2:Z(p)>0\}$, its closure in $S^2$ is $\overline H=\{Z\ge0\}$. Every neighbourhood of an equatorial point meets $H$; a point with $Z<0$ has a small ambient ball disjoint from $H$.
\end{atlasdefinition}

\begin{atlasdefinition}{\texttt{12-01}: 국소유한 집합족 / Locally finite family}
\label{def:12-01}
X를 위상공간 (topological space), \((A_i)_{i\in I}\)를 X의 부분집합족이라 하자. 점과 근방은 X 안에서 취한다. 집합족 \(\{A_i\}\)가 국소유한 (locally finite)이라는 것은 모든 점 x에 대해 어떤 열린 근방 (open neighborhood) V가 존재하여 \(\{i:V\cap A_i\ne\varnothing\}\)가 유한하다는 뜻이다.

\textbf{조건 확인 / Conditions.} 집합족 전체가 유한할 필요는 없다.

\textbf{직접 선행 / Prerequisites.} \hyperref[def:01-02]{\texttt{01-02} 열린 근방}

\exampleheading
The family $\{S^2\setminus\{S\},S^2\setminus\{N\}\}$ is locally finite: every neighbourhood meets at most two members. In contrast $V_n=\{p\in S^2:Z(p)>-1/n\}$ is not locally finite at $N$, since every $V_n$ contains $N$.
\end{atlasdefinition}

\begin{atlasdefinition}{\texttt{12-02}: 열린덮개의 세분 / Refinement of an open cover}
\label{def:12-02}
X를 위상공간 (topological space), \((U_i)_{i\in I}\)와 \((V_a)_{a\in A}\)를 X의 열린덮개 (open cover)라 하자. \(\{V_a\}\)가 \(\{U_i\}\)의 세분이라는 것은 새 열린집합족도 같은 공간을 덮고, 모든 a에 대해 어떤 i가 존재하여 \(V_a\subset U_i\)인 것이다.

\textbf{조건 확인 / Conditions.} 새 집합이 원래 집합 안에 들어간다.

\textbf{직접 선행 / Prerequisites.} \hyperref[def:01-03]{\texttt{01-03} 열린덮개}

\exampleheading
For the standard two-chart cover, $V_0=\{Z>-3/4\}$ and $V_\infty=\{Z<3/4\}$ form an open refinement: $V_0\subset U_0$, $V_\infty\subset U_\infty$, and $V_0\cup V_\infty=S^2$. All sets here are relative subsets of $S^2$.
\end{atlasdefinition}

\begin{atlasdefinition}{\texttt{12-03}: Paracompact 공간 / Paracompact space}
\label{def:12-03}
모든 열린덮개 (open cover)가 국소유한 (locally finite) 열린 세분을 갖는 위상공간 (topological space)을 paracompact라 한다.

\textbf{조건 확인 / Conditions.} 이 정의 자체에는 Hausdorff를 별도로 넣지 않는다.

\textbf{직접 선행 / Prerequisites.} \hyperref[def:12-01]{\texttt{12-01} 국소유한 집합족}; \hyperref[def:12-02]{\texttt{12-02} 열린덮개의 세분}

\exampleheading
The unit sphere is compact because it is closed and bounded in $\R^3$ (Heine--Borel). Every open cover therefore has a finite subcover, which is itself a locally finite open refinement. This verifies paracompactness of this particular $S^2$.
\end{atlasdefinition}

\subsection{대수와 범주의 언어 / Algebra and categories}

\begin{atlasdefinition}{\texttt{02-01}: 아벨군 / Abelian group}
\label{def:02-01}
집합 \(A\)에 닫힌 이항연산 \(+\)가 주어지고, 모든 \(a,b,c\in A\)에 대해 \((a+b)+c=a+(b+c)\), \(a+b=b+a\)다. 영원소 \(0\)가 있어 \(a+0=a\)이고, 각 \(a\)에 대해 \(a+(-a)=0\)인 역원 \(-a\)가 존재한다.

\textbf{조건 확인 / Conditions.} 가환성이 포함된다.

\textbf{직접 선행 / Prerequisites.} 집합·함수·실수·복소수 등 공통 배경.

\exampleheading
The smooth real functions $C^\infty(S^2,\R)$ form an abelian group under pointwise addition. With $H(p)=Z(p)$, the additive inverse of $H$ is $-H$, and $H+(-H)=0$ is the zero function on the whole sphere.
\end{atlasdefinition}

\begin{atlasdefinition}{\texttt{26-01}: 환 / Ring (commutative, with identity)}
\label{def:26-01}
단위원을 가진 가환환 (commutative ring)은 집합 R과 두 이항연산 \(+,\cdot:R\times R\to R\)로 이루어진다. \((R,+)\)는 아벨군 (abelian group)이며, 모든 \(a,b,c\in R\)에 대해 \((ab)c=a(bc)\), \(ab=ba\), \(a(b+c)=ab+ac\), \((a+b)c=ac+bc\)를 만족한다. 또한 어떤 \(1_R\in R\)가 존재하여 모든 \(a\in R\)에 대해 \(1_Ra=a1_R=a\)다.

\textbf{조건 확인 / Conditions.} 가환성과 단위원의 존재를 포함한 정의다. 일반적인 비가환환의 정의와 범위를 구별한다.

\textbf{직접 선행 / Prerequisites.} \hyperref[def:02-01]{\texttt{02-01} 아벨군}

\exampleheading
On the chart $U_0\subset\CP^1$, $\mathcal O(U_0)=\mathcal O(\C)$ is a unital commutative ring under pointwise operations. It contains both $z$ and $e^z$; the latter shows why this analytic ring is larger than $\C[z]$.
\end{atlasdefinition}

\begin{atlasdefinition}{\texttt{N01}: 가군 / Module}
\label{def:N01}
R을 단위원을 가진 가환환 (commutative ring)이라 하자. R-가군은 아벨군 (abelian group) M과 스칼라 곱 \(R\times M\to M,(r,m)\mapsto rm\)으로, 모든 \(r,s\in R\), \(m,n\in M\)에 대해 \[r(m+n)=rm+rn,\quad(r+s)m=rm+sm,\quad(rs)m=r(sm),\quad1_Rm=m\]을 만족하는 자료다. R-가군 준동형 (homomorphism)은 덧셈과 R의 스칼라 곱을 보존하는 함수다.

\textbf{조건 확인 / Conditions.} 환 R을 고정하며, \(1_R\)이 항등적으로 작용해야 한다.

\textbf{직접 선행 / Prerequisites.} \hyperref[def:02-01]{\texttt{02-01} 아벨군}; \hyperref[def:26-01]{\texttt{26-01} 환}

\exampleheading
On $U_0\subset\CP^1$, the holomorphic one-forms $\mathcal K(U_0)=\mathcal O(\C)\,dz$ form a module over $\mathcal O(\C)$. For instance $(1+z)(z\,dz)=(z+z^2)\,dz$. The unit function acts as the identity.
\end{atlasdefinition}

\begin{atlasdefinition}{\texttt{28-01}: 아이디얼 / Ideal}
\label{def:28-01}
가환환 (commutative ring) R의 아이디얼 (ideal) J는 덧셈 부분군 (subgroup)으로, 모든 \(r\in R\), \(a\in J\)에 대해 \(ra\in J\)를 만족한다.

\textbf{조건 확인 / Conditions.} J 안 원소끼리의 곱만 검사하는 것이 아니다.

\textbf{직접 선행 / Prerequisites.} \hyperref[def:26-01]{\texttt{26-01} 환}

\exampleheading
In the analytic local ring $\mathcal O_{\CP^1,0}=\C\{z\}$, the set $(z)=\{zf(z):f\in\C\{z\}\}$ is an ideal. It contains $z+z^2$ and is closed under multiplication by every convergent germ.
\end{atlasdefinition}

\begin{atlasdefinition}{\texttt{26-02}: 환준동형 / Ring homomorphism}
\label{def:26-02}
R,S는 단위원을 가진 가환환 (commutative ring)이고 아래의 a,b는 R의 임의의 원소다. 환 사이의 함수 \(f:R\to S\)로 모든 a,b에 대해 \(f(a+b)=f(a)+f(b)\), \(f(ab)=f(a)f(b)\), \(f(1_R)=1_S\)를 만족한다.

\textbf{조건 확인 / Conditions.} 단위원 보존도 요구한다.

\textbf{직접 선행 / Prerequisites.} \hyperref[def:26-01]{\texttt{26-01} 환}

\exampleheading
At $p=[1:0]\in\CP^1$, evaluation $\mathcal O_{X,p}\to\C$, $[f]_p\mapsto f(p)$, is a unital ring homomorphism: $[fg]_p$ maps to $f(p)g(p)$ and $[1]_p$ maps to $1$.
\end{atlasdefinition}

\begin{atlasdefinition}{\texttt{N02}: 환의 국소화 / Localization of a ring}
\label{def:N02}
R을 단위원을 가진 가환환 (commutative ring), S를 \(1\in S\), \(s,t\in S\Rightarrow st\in S\)를 만족하는 부분집합이라 하자. \(S^{-1}R\)의 원소는 \((a,s)\in R\times S\)의 동치류 (equivalence class) \(a/s\)이며 \[(a,s)\sim(b,t)\iff\exists u\in S:\ u(at-bs)=0.\]연산은 \(a/s+b/t=(at+bs)/(st)\), \((a/s)(b/t)=ab/(st)\)다. 표준 환준동형은 \(R\to S^{-1}R,a\mapsto a/1\)이다. 소아이디얼 \(\mathfrak p\)에 대해 \(R_{\mathfrak p}:=(R\setminus\mathfrak p)^{-1}R\)라 쓴다. 소아이디얼은 \(\mathfrak p\ne R\)이고 \(ab\in\mathfrak p\Rightarrow a\in\mathfrak p\) 또는 \(b\in\mathfrak p\)인 아이디얼 (ideal)이다.

\textbf{조건 확인 / Conditions.} 영인수가 있을 수 있으므로 동치관계 (equivalence relation)를 단순히 at=bs로 줄이지 않는다. 0\ensuremath{\in}S이면 국소화는 영환이다.

\textbf{직접 선행 / Prerequisites.} \hyperref[def:28-01]{\texttt{28-01} 아이디얼}; \hyperref[def:26-02]{\texttt{26-02} 환준동형}

\exampleheading
On the algebraic chart of $\CP^1$, take $R=\C[z]$ and $S=\{b:b(0)\ne0\}$. Then $R_{(z)}=\{a(z)/b(z):b(0)\ne0\}$ is the algebraic local ring at $0$; $1/(1-z)$ belongs to it but $1/z$ does not. It embeds in the analytic local ring $\C\{z\}$, and these rings must not be identified: $e^z$ is an analytic germ but not a rational function.
\end{atlasdefinition}

\begin{atlasdefinition}{\texttt{02-02}: 아벨군 준동형 / Homomorphism of abelian groups}
\label{def:02-02}
아벨군 (abelian group) 사이의 함수 \(f:A\to B\)로, 모든 \(a,a'\in A\)에 대해 \(f(a+a')=f(a)+f(a')\)를 만족한다.

\textbf{조건 확인 / Conditions.} 모든 두 원소의 덧셈을 보존한다.

\textbf{직접 선행 / Prerequisites.} \hyperref[def:02-01]{\texttt{02-01} 아벨군}

\exampleheading
Evaluation at $N$ gives an abelian-group homomorphism $\operatorname{ev}_N:C^\infty(S^2,\R)\to\R$, $f\mapsto f(N)$. In particular $\operatorname{ev}_N(2+3H)=5$ and $\operatorname{ev}_N(f+g)=f(N)+g(N)$.
\end{atlasdefinition}

\begin{atlasdefinition}{\texttt{28-02}: 극대아이디얼 / Maximal ideal}
\label{def:28-02}
R을 단위원을 가진 가환환 (commutative ring)이라 하자. \(\mathfrak m\subsetneq R\)인 아이디얼 (ideal)로, \(\mathfrak m\subset J\subset R\)인 모든 아이디얼 J가 \(J=\mathfrak m\) 또는 \(J=R\)이면 극대아이디얼 (maximal ideal)이다.

\textbf{조건 확인 / Conditions.} 진아이디얼이어야 한다.

\textbf{직접 선행 / Prerequisites.} \hyperref[def:28-01]{\texttt{28-01} 아이디얼}

\exampleheading
The ideal $(z)\subset\mathcal O_{\CP^1,0}$ is maximal because evaluation gives $\mathcal O_{\CP^1,0}/(z)\simeq\C$, a field. It consists exactly of germs vanishing at $0$.
\end{atlasdefinition}

\begin{atlasdefinition}{\texttt{28-03}: 국소환 / Local ring}
\label{def:28-03}
단위원을 가진 가환환 (commutative ring) R에 대한 정의다. 극대아이디얼 (maximal ideal)을 정확히 하나 갖는 환이다.

\textbf{조건 확인 / Conditions.} 존재와 유일성 모두 필요하다.

\textbf{직접 선행 / Prerequisites.} \hyperref[def:28-02]{\texttt{28-02} 극대아이디얼}

\exampleheading
A germ $f\in\mathcal O_{\CP^1,0}$ is a unit exactly when $f(0)\ne0$, since its reciprocal is then holomorphic near $0$. Thus all nonunits form the single maximal ideal $(z)$, making this ring local.
\end{atlasdefinition}

\begin{atlasdefinition}{\texttt{13-01}: 범주 / Category}
\label{def:13-01}
대상들과 각 대상쌍의 사상 집합 \(\operatorname{Hom}(A,B)\), 합성 \(\operatorname{Hom}(B,C)\times\operatorname{Hom}(A,B)\to\operatorname{Hom}(A,C)\), 항등사상 (identity morphism) \(\mathrm{id}_A\)로 이루어진다. 합성은 결합적이고 \(\mathrm{id}_Bf=f=f\mathrm{id}_A\)다.

\textbf{조건 확인 / Conditions.} 합성 가능한 사상에 대해서만 합성한다.

\textbf{직접 선행 / Prerequisites.} 집합·함수·실수·복소수 등 공통 배경.

\exampleheading
Fix $X=\CP^1$. In $\operatorname{Sh}(X,\mathbf{Ab})$, objects include $\underline\Z$, $\mathcal O_X$ under addition, and the skyscraper $i_*\C$ at $p=0$. Morphisms are compatible families of homomorphisms on open sets. Evaluation $\mathcal O_X\to i_*\C$ is one such morphism.
\end{atlasdefinition}

\begin{atlasdefinition}{\texttt{13-02}: 동형사상 / Isomorphism}
\label{def:13-02}
범주 \(\mathcal C\)의 대상 A,B와 사상 \(f:A\to B\)를 고정한다. \(\mathrm{id}_A,\mathrm{id}_B\)는 각 대상의 항등사상 (identity morphism)이고 사상의 곱은 합성을 뜻한다. \(f:A\to B\)에 대해 \(g:B\to A\)가 존재하여 \(gf=\mathrm{id}_A\), \(fg=\mathrm{id}_B\)이면 f를 동형사상 (isomorphism)이라 한다.

\textbf{조건 확인 / Conditions.} 양쪽 합성이 각각의 항등사상 (identity morphism)이어야 한다.

\textbf{직접 선행 / Prerequisites.} \hyperref[def:13-01]{\texttt{13-01} 범주}

\exampleheading
On $\CP^1$, the map $\mathcal O_X(-2[\infty])\to\mathcal K_X$, $f\mapsto f\,dz$, is a sheaf isomorphism. On $U_0$ it sends $1$ to $dz$; on $U_\infty$ it sends $w^2$ to $-dw$. Both images are local frames, giving local inverses which agree on overlaps.
\end{atlasdefinition}

\begin{atlasdefinition}{\texttt{B01}: 실수·복소 벡터공간 / Real and complex vector space}
\label{def:B01}
체 \(k=\mathbb R\) 또는 \(\mathbb C\)를 고정한다. k-벡터공간은 아벨군 V와 스칼라곱 \(k\times V\to V\)로, 모든 \(a,b\in k\), \(u,v\in V\)에 대해 \(a(u+v)=au+av\), \((a+b)v=av+bv\), \((ab)v=a(bv)\), \(1v=v\)를 만족한다. k-선형사상은 덧셈과 이 스칼라곱을 보존한다.

\textbf{조건 확인 / Conditions.} 실수 선형성과 복소 선형성을 구별한다. 체 자체와 집합·함수는 기초 배경으로 사용한다.

\textbf{직접 선행 / Prerequisites.} \hyperref[def:02-01]{\texttt{02-01} 아벨군}

\exampleheading
At $N=(0,0,1)\in S^2$, $T_NS^2=\{(a,b,0):a,b\in\R\}$ is a real vector space with basis $e_1,e_2$. Its complexification has basis $e_1,e_2$ over $\C$, while $T_N^{1,0}\CP^1$ has complex dimension one. The scalar field changes the dimension statement.
\end{atlasdefinition}

\begin{atlasdefinition}{\texttt{B02}: 쌍대·텐서곱·외대수 / Dual, tensor product and exterior power}
\label{def:B02}
k-벡터공간 V의 쌍대는 \(V^*=\operatorname{Hom}_k(V,k)\)이다. \(V\otimes_kW\)는 쌍선형사상 \(b:V\times W\to V\otimes_kW\)와 함께, 모든 k-벡터공간 Z와 쌍선형사상 f에 대해 유일한 선형사상 \(\widetilde f:V\otimes_kW\to Z\)가 \(f=\widetilde f\circ b\)를 만족하는 공간이다. \(\bigwedge^qV\)는 같은 보편성에서 입력을 \(V^q\), 사상을 교대 q-다중선형사상으로 바꾼 공간이다. \(\bigwedge^0V=k\).

\textbf{조건 확인 / Conditions.} 교대란 두 입력이 같으면 값이 0이라는 뜻이다. 여기서 k는 실수체 또는 복소수체이다. 이 항목은 공통 선형대수 배경을 모은 참조 항목이다.

\textbf{직접 선행 / Prerequisites.} \hyperref[def:B01]{\texttt{B01} 실수·복소 벡터공간}

\exampleheading
At $N\in S^2$, take the tangent basis $e_1,e_2$ and dual basis $\epsilon^1,\epsilon^2$. Then $(\epsilon^1\otimes\epsilon^2)(e_1,e_2)=1$, while $(\epsilon^1\wedge\epsilon^2)(e_1,e_2)=1$ and its value on $(e_2,e_1)$ is $-1$. The first tensor is not alternating, and the dual covector satisfies $\epsilon^1(e_1)=1$.
\end{atlasdefinition}

\begin{atlasdefinition}{\texttt{13-03}: 공변함자 / Covariant functor}
\label{def:13-03}
\(\mathcal C,\mathcal D\)를 범주라 하자. 아래에서 A,B는 \(\mathcal C\)의 대상이고 f,g는 \(\mathcal C\)의 합성 가능한 사상이다. \(T:\mathcal C\to\mathcal D\)는 대상 A를 T(A)로, 사상 \(f:A\to B\)를 \(T(f):T(A)\to T(B)\)로 보내며 \(T(\mathrm{id}_A)=\mathrm{id}_{T(A)}\), \(T(gf)=T(g)T(f)\)를 만족한다.

\textbf{조건 확인 / Conditions.} 합성과 항등사상 (identity morphism)을 모두 보존한다.

\textbf{직접 선행 / Prerequisites.} \hyperref[def:13-01]{\texttt{13-01} 범주}

\exampleheading
On $X=\CP^1$, $\Gamma(X,-)$ is covariant: a sheaf map $\phi:\mathcal F\to\mathcal G$ gives $\phi_X:\mathcal F(X)\to\mathcal G(X)$. The sheaf map $\times2:\underline\Z\to\underline\Z$ is sent to $\times2:\Z\to\Z$.
\end{atlasdefinition}

\begin{atlasdefinition}{\texttt{14-02}: 유한 직합 / Finite direct sum / Biproduct}
\label{def:14-02}
가법적 문맥에서 \(A\oplus B\)는 포함 \(i_A:A\to A\oplus B\), \(i_B:B\to A\oplus B\)와 사영 \(p_A:A\oplus B\to A\), \(p_B:A\oplus B\to B\)를 가지며 \[p_Ai_A=1_A,\quad p_Bi_B=1_B,\quad p_Ai_B=p_Bi_A=0,\]\[i_Ap_A+i_Bp_B=1_{A\oplus B}\]을 만족하는 대상이다. 유한 개 대상에 반복하고 빈 직합은 영대상 (zero object)이다.

\textbf{조건 확인 / Conditions.} 마지막 항등사상 (identity morphism) 분해식도 필요하다.

\textbf{직접 선행 / Prerequisites.} \hyperref[def:13-01]{\texttt{13-01} 범주}; \hyperref[def:02-01]{\texttt{02-01} 아벨군}

\exampleheading
On $\CP^1$, $\mathcal O_X\oplus\mathcal K_X$ has local sections $(f,\alpha)$. The inclusions $f\mapsto(f,0)$, $\alpha\mapsto(0,\alpha)$ and coordinate projections satisfy $p_i\iota_j=\delta_{ij}$ and $\iota_1p_1+\iota_2p_2=\id$.
\end{atlasdefinition}

\begin{atlasdefinition}{\texttt{14-01}: 영대상 / Zero object}
\label{def:14-01}
범주의 대상 0이 영대상 (zero object)이라는 것은 모든 대상 A에 대해 \(0\to A\)와 \(A\to0\)인 사상이 각각 정확히 하나 존재한다는 뜻이다.

\textbf{조건 확인 / Conditions.} 나가는 사상과 들어오는 사상 모두 유일하다.

\textbf{직접 선행 / Prerequisites.} \hyperref[def:13-01]{\texttt{13-01} 범주}

\exampleheading
The zero sheaf on $S^2$ assigns the zero abelian group to every open set, including the empty set. There is exactly one sheaf map $0\to\mathcal F$ and exactly one $\mathcal F\to0$ for any abelian sheaf $\mathcal F$.
\end{atlasdefinition}

\begin{atlasdefinition}{\texttt{14-03}: 범주론적 핵 / Categorical kernel}
\label{def:14-03}
사상 집합마다 아벨군 (abelian group) 구조가 있고 합성이 쌍선형인 범주에서 사상 \(f:A\to B\)를 고정한다. 0은 해당 사상 집합의 영사상 (zero morphism)이다. \(f:A\to B\)의 핵은 \(k:K\to A\)로, \(fk=0\)이고 모든 대상 T와 모든 \(g:T\to A\)에 대해 \(fg=0\)이면 유일한 \(\widetilde g:T\to K\)가 존재하여 \(g=k\widetilde g\)가 되는 것이다.

\textbf{조건 확인 / Conditions.} 핵으로 들어가는 분해 사상의 존재와 유일성이 필요하다.

\textbf{직접 선행 / Prerequisites.} \hyperref[def:14-01]{\texttt{14-01} 영대상}

\exampleheading
For $i:\{0\}\hookrightarrow\CP^1$, the categorical kernel of evaluation $\mathcal O_X\to i_*\C$ is the ideal sheaf $\mathcal I_0$: near $0$ its sections are $z f(z)$, and away from $0$ it equals $\mathcal O_X$. A sheaf map killed by evaluation factors uniquely through this inclusion.
\end{atlasdefinition}

\begin{atlasdefinition}{\texttt{14-04}: 범주론적 여핵 / Categorical cokernel}
\label{def:14-04}
사상 집합마다 아벨군 (abelian group) 구조가 있고 합성이 쌍선형인 범주에서 사상 \(f:A\to B\)를 고정한다. 0은 해당 사상 집합의 영사상 (zero morphism)이다. \(f:A\to B\)의 여핵은 \(q:B\to Q\)로, \(qf=0\)이고 모든 대상 T와 모든 \(g:B\to T\)에 대해 \(gf=0\)이면 유일한 \(\widetilde g:Q\to T\)가 존재하여 \(g=\widetilde gq\)가 되는 것이다.

\textbf{조건 확인 / Conditions.} 여핵에서 나가는 분해 사상이다.

\textbf{직접 선행 / Prerequisites.} \hyperref[def:14-01]{\texttt{14-01} 영대상}

\exampleheading
The cokernel of $\mathcal I_0\hookrightarrow\mathcal O_X$ on $\CP^1$ is $i_*\C$, with quotient map evaluation. A map from $\mathcal O_X$ vanishing on $\mathcal I_0$ factors uniquely through the value at $0$.
\end{atlasdefinition}

\begin{atlasdefinition}{\texttt{14-05}: 범주론적 상 / Categorical image}
\label{def:14-05}
핵과 여핵이 있는 가법적 범주에서 \(f:A\to B\)의 상은 \(\operatorname{im}f:=\ker(B\to\operatorname{coker}f)\)다.

\textbf{조건 확인 / Conditions.} 공역 B에서 여핵으로 가는 사상의 핵이다.

\textbf{직접 선행 / Prerequisites.} \hyperref[def:14-03]{\texttt{14-03} 범주론적 핵}; \hyperref[def:14-04]{\texttt{14-04} 범주론적 여핵}

\exampleheading
On $S^2$, the image sheaf of $\times2:\underline\Z\to\underline\Z$ is $2\underline\Z$. Its sections are locally constant even integers. It is the kernel of the quotient to $\underline{\Z/2\Z}$.
\end{atlasdefinition}

\begin{atlasdefinition}{\texttt{14-06}: 범주론적 여상 / Coimage}
\label{def:14-06}
핵과 여핵이 있는 가법적 범주에서 \(f:A\to B\)의 여상은 \(\operatorname{coim}f:=\operatorname{coker}(\ker f\to A)\)다.

\textbf{조건 확인 / Conditions.} 핵에서 정의역 A로 가는 사상의 여핵이다.

\textbf{직접 선행 / Prerequisites.} \hyperref[def:14-03]{\texttt{14-03} 범주론적 핵}; \hyperref[def:14-04]{\texttt{14-04} 범주론적 여핵}

\exampleheading
The coimage of $\times2:\underline\Z\to\underline\Z$ on $S^2$ is $\underline\Z/0$. The canonical coimage-to-image map sends $n\mapsto2n$ and is an isomorphism onto $2\underline\Z$, even though the original map is not onto $\underline\Z$.
\end{atlasdefinition}

\begin{atlasdefinition}{\texttt{15-01}: 단사상 - 범주론 / Monomorphism}
\label{def:15-01}
범주 \(\mathcal C\)에서 사상 \(j:A\to B\)를 고정한다. T는 \(\mathcal C\)의 임의의 대상이다. \(j:A\to B\)가 단사상 (monomorphism)이라는 것은 모든 대상 T와 모든 \(g,h:T\to A\)에 대해 \(jg=jh\Rightarrow g=h\)라는 뜻이다.

\textbf{조건 확인 / Conditions.} g와 h의 출발 대상은 같고 도착 대상은 A다.

\textbf{직접 선행 / Prerequisites.} \hyperref[def:13-01]{\texttt{13-01} 범주}

\exampleheading
Multiplication by $2$ on $\underline\Z$ over $S^2$ is a monomorphism: if $2a=2b$ for locally constant integer sections then $a=b$. Its stalk at every point is the injective map $\Z\to\Z$, $n\mapsto2n$.
\end{atlasdefinition}

\begin{atlasdefinition}{\texttt{14-07}: 아벨범주 / Abelian category}
\label{def:14-07}
범주 \(\mathcal C\)가 다음을 만족하면 아벨범주 (abelian category)라 한다. 모든 대상쌍 A,B의 \(\operatorname{Hom}_{\mathcal C}(A,B)\)가 아벨군 (abelian group)이고 합성은 각 변수에 대해 가법적이다. 영대상 (zero object)과 유한 직합이 존재한다. 또한 모든 사상 \(f:A\to B\)에 핵 \(k_f:K_f\to A\)와 여핵 \(q_f:B\to Q_f\)가 존재한다. \[p_f:A\to\operatorname{coim}f:=\operatorname{coker}k_f,\qquad i_f:\operatorname{im}f:=\ker q_f\to B\]를 각각 그 여핵사상과 핵사상이라 하자. 핵·여핵의 보편성으로 정해지는 사상 \[\overline f:\operatorname{coim}f\to\operatorname{im}f,\qquad f=i_f\circ\overline f\circ p_f\]이 모든 f에 대해 동형이어야 한다.

\textbf{조건 확인 / Conditions.} 핵·여핵의 존재만으로는 부족하고 여상\ensuremath{\to}상의 자연사상이 동형이어야 한다.

\textbf{직접 선행 / Prerequisites.} \hyperref[def:14-02]{\texttt{14-02} 유한 직합}; \hyperref[def:14-03]{\texttt{14-03} 범주론적 핵}; \hyperref[def:14-04]{\texttt{14-04} 범주론적 여핵}; \hyperref[def:14-05]{\texttt{14-05} 범주론적 상}; \hyperref[def:14-06]{\texttt{14-06} 범주론적 여상}; \hyperref[def:15-01]{\texttt{15-01} 단사상 - 범주론}

\exampleheading
The category of abelian sheaves on $S^2$ is abelian. The concrete morphism $\times2:\underline\Z\to\underline\Z$ has kernel $0$, image $2\underline\Z$, and cokernel $\underline{\Z/2\Z}$, and its coimage maps isomorphically onto its image. These illustrate the axioms; the assertion for every morphism is the general abelian-sheaf theorem.
\end{atlasdefinition}

\begin{atlasdefinition}{\texttt{02-03}: 핵 - 아벨군 / Kernel of a homomorphism}
\label{def:02-03}
아벨군 (abelian group) 준동형 (homomorphism) \(f:A\to B\)에 대해 \(\ker f=\{a\in A:f(a)=0_B\}\)다.

\textbf{조건 확인 / Conditions.} 핵은 정의역 A의 부분군 (subgroup)이다.

\textbf{직접 선행 / Prerequisites.} \hyperref[def:02-02]{\texttt{02-02} 아벨군 준동형}

\exampleheading
For $\operatorname{ev}_N$ on $S^2$, $\ker\operatorname{ev}_N=\{f:f(N)=0\}$. The nonzero function $H-1$ belongs to this kernel, since $H(N)=1$.
\end{atlasdefinition}

\begin{atlasdefinition}{\texttt{02-04}: 상 - 아벨군 / Image of a homomorphism}
\label{def:02-04}
아벨군 (abelian group) 준동형 (homomorphism) \(f:A\to B\)에 대해 \(\operatorname{im}f=\{f(a):a\in A\}\subset B\)다.

\textbf{조건 확인 / Conditions.} 상은 공역 B의 부분군 (subgroup)이다.

\textbf{직접 선행 / Prerequisites.} \hyperref[def:02-02]{\texttt{02-02} 아벨군 준동형}

\exampleheading
The image of $\operatorname{ev}_N:C^\infty(S^2,\R)\to\R$ is all of $\R$: every $a\in\R$ is the image of the constant function $f\equiv a$.
\end{atlasdefinition}

\begin{atlasdefinition}{\texttt{02-05}: 몫군 / Quotient group}
\label{def:02-05}
아벨군 (abelian group) \(A\)의 부분군 (subgroup) \(K\)에 대해 \(A/K=\{a+K:a\in A\}\)이며, \(a+K=b+K\iff a-b\in K\)다. 덧셈은 \((a+K)+(b+K)=(a+b)+K\)다.

\textbf{조건 확인 / Conditions.} K는 A의 부분군 (subgroup)이어야 한다.

\textbf{직접 선행 / Prerequisites.} \hyperref[def:02-01]{\texttt{02-01} 아벨군}

\exampleheading
The quotient $C^\infty(S^2,\R)/\ker\operatorname{ev}_N$ is isomorphic to $\R$ by $[f]\mapsto f(N)$. Thus $H$ and the constant function $1$ define the same quotient class although they are different functions on $S^2$.
\end{atlasdefinition}

\begin{atlasdefinition}{\texttt{B03}: Hermitian 내적 / Hermitian inner product}
\label{def:B03}
복소벡터공간 V 위 형식 \(h:V\times V\to\mathbb C\)가 첫 인자에 복소선형, 둘째 인자에 켤레선형이고, 모든 u,v에 대해 \(h(u,v)=\overline{h(v,u)}\), 모든 \(v\ne0\)에 대해 \(h(v,v)>0\)이면 Hermitian 내적이다.

\textbf{조건 확인 / Conditions.} 이 정의집은 첫 인자에 선형인 규약을 사용한다. 양의 정부호 조건을 생략하지 않는다.

\textbf{직접 선행 / Prerequisites.} \hyperref[def:B01]{\texttt{B01} 실수·복소 벡터공간}

\exampleheading
In the complex tangent line $T_0^{1,0}\CP^1=\C\,\partial_z$, set $h(a\partial_z,b\partial_z)=2a\overline b$. Thus $h(\partial_z,i\partial_z)=-2i$, $h(i\partial_z,\partial_z)=2i$, and $h(a\partial_z,a\partial_z)=2|a|^2>0$ for $a\ne0$.
\end{atlasdefinition}

\begin{atlasdefinition}{\texttt{02-06}: 여핵 - 아벨군 / Cokernel of a homomorphism}
\label{def:02-06}
아벨군 (abelian group) 준동형 (homomorphism) \(f:A\to B\)의 여핵은 \(\operatorname{coker}f=B/\operatorname{im}f\)다.

\textbf{조건 확인 / Conditions.} 정의역이 아니라 공역을 상으로 나눈다.

\textbf{직접 선행 / Prerequisites.} \hyperref[def:02-04]{\texttt{02-04} 상 - 아벨군}; \hyperref[def:02-05]{\texttt{02-05} 몫군}

\exampleheading
On the connected sphere, $\Gamma(S^2,\underline\Z)=\Z$. The map induced by multiplication by $2$ has cokernel $\Z/2\Z$; the class of the constant section $1$ is nonzero and twice that class is zero.
\end{atlasdefinition}

\begin{atlasdefinition}{\texttt{13-04}: 반변함자 / Contravariant functor}
\label{def:13-04}
\(\mathcal C,\mathcal D\)를 범주라 하자. 반변함자 (contravariant functor) \(T:\mathcal C^{\mathrm{op}}\to\mathcal D\)는 다음 자료다. A,B는 \(\mathcal C\)의 대상이고 f,g는 \(\mathcal C\)의 합성 가능한 사상이다. \(\mathcal C^{\mathrm{op}}\)는 \(\mathcal C\)의 모든 사상 방향을 뒤집은 범주다. 대상 A를 T(A)로, 사상 \(f:A\to B\)를 \(T(f):T(B)\to T(A)\)로 보내고, 항등사상 (identity morphism)을 보존하며 \(T(gf)=T(f)T(g)\)를 만족하는 자료다.

\textbf{조건 확인 / Conditions.} 방향을 뒤집으면 합성 순서도 뒤집힌다.

\textbf{직접 선행 / Prerequisites.} \hyperref[def:13-01]{\texttt{13-01} 범주}

\exampleheading
On $X=\CP^1$, the functor $\operatorname{Hom}(-,\mathcal O_X)$ on $\mathcal O_X$-modules is contravariant. Precomposition with $\times2:\mathcal O_X\to\mathcal O_X$ sends a map $\psi$ to $\psi\circ(\times2)$; the arrow between Hom spaces goes in the reversed direction.
\end{atlasdefinition}

\subsection{미분과 복소기하 / Differential and complex geometry}

\begin{atlasdefinition}{\texttt{27-01}: 정칙함수 / Holomorphic function}
\label{def:27-01}
열린 \(U\subset\mathbb C\)의 함수 \(f:U\to\mathbb C\)가 모든 \(z\in U\)에서 복소미분 가능하면 정칙이라 한다.

\textbf{조건 확인 / Conditions.} 열린집합 (open set)의 모든 점에서 복소미분 가능해야 한다.

\textbf{직접 선행 / Prerequisites.} 집합·함수·실수·복소수 등 공통 배경.

\exampleheading
On $U_0\simeq\C\subset\CP^1$, $f(z)=z^2$ is holomorphic, with derivative $2z$. The function $\bar z$ is not holomorphic. A holomorphic function on all of the compact connected $\CP^1$ must be constant; a chart function need not extend globally.
\end{atlasdefinition}

\begin{atlasdefinition}{\texttt{B04}: 복소다양체 / Complex manifold}
\label{def:B04}
복소 n차원 다양체는 Hausdorff 제2가산 위상공간 X와 최대 복소 좌표계 모음이다. 좌표 \(\varphi_i:U_i\to V_i\subset\mathbb C^n\)는 열린집합 사이의 위상동형이고 \(\bigcup_iU_i=X\)이며, 모든 겹침에서 \(\varphi_j\circ\varphi_i^{-1}\)는 쌍정칙이다. 쌍정칙은 사상과 역이 정칙이라는 뜻이다. 여러 복소변수에서 정칙은 각 성분이 각 점 근방에서 수렴하는 복소 거듭제곱급수로 표현된다는 뜻이다.

\textbf{조건 확인 / Conditions.} 정칙사상은 이 좌표들에서 성분별 정칙인 사상이다. N04에 포함된 다양체 정의를 독립적으로 찾을 수 있게 분리했다.

\textbf{직접 선행 / Prerequisites.} \hyperref[def:27-01]{\texttt{27-01} 정칙함수}; \hyperref[def:06-01]{\texttt{06-01} 연속사상}; \hyperref[def:06-02]{\texttt{06-02} Hausdorff 공간}; \hyperref[def:06-03]{\texttt{06-03} 열린 기저}

\exampleheading
For $\CP^1=\{[S:T]\}$, the charts $z=T/S$ and $w=S/T$ cover the space and satisfy $w=1/z$ on $\C^\times$. They define a complex manifold of complex dimension one and underlying real dimension two, namely the fixed sphere.
\end{atlasdefinition}

\begin{atlasdefinition}{\texttt{27-02}: 리만곡면 / Riemann surface}
\label{def:27-02}
Hausdorff이고 가산 열린 기저를 갖는 공간에, 공간을 덮는 열린집합 (open set)에서 \(\mathbb C\)의 열린집합으로 가는 위상동형 좌표들이 주어지고 모든 겹침의 좌표변환이 정칙인 복소 1차원 다양체다.

\textbf{조건 확인 / Conditions.} 여기서 복소곡선은 리만곡면 (Riemann surface)을 뜻한다.

\textbf{직접 선행 / Prerequisites.} \hyperref[def:06-01]{\texttt{06-01} 연속사상}; \hyperref[def:06-02]{\texttt{06-02} Hausdorff 공간}; \hyperref[def:06-03]{\texttt{06-03} 열린 기저}; \hyperref[def:27-01]{\texttt{27-01} 정칙함수}

\exampleheading
The sphere with coordinates $z=(X+iY)/(1+Z)$ and $w=(X-iY)/(1-Z)$ is a Riemann surface. The overlap relation $w=1/z$ is holomorphic with nonzero derivative for $z\ne0$. This complex structure identifies it with $\CP^1$.
\end{atlasdefinition}

\begin{atlasdefinition}{\texttt{06-04}: 매끄러운 다양체 / Smooth manifold}
\label{def:06-04}
n\ensuremath{\ge}0을 정수라 하자. 매끄러운 n-다양체는 Hausdorff이고 가산 열린 기저를 갖는 위상공간 (topological space) M과 다음 매끄러운 좌표계 모음으로 정해지는 매끄러운 구조의 쌍이다. 각 좌표 \(\varphi_\alpha:U_\alpha\to V_\alpha\subset\mathbb R^n\)는 열린집합 (open set) 사이의 위상동형이고 \(M=\bigcup_\alpha U_\alpha\)다. 모든 \ensuremath{\alpha},\ensuremath{\beta}에 대해 좌표변환 \[\varphi_\beta\circ\varphi_\alpha^{-1}:\varphi_\alpha(U_\alpha\cap U_\beta)\to\varphi_\beta(U_\alpha\cap U_\beta)\]은 C\ensuremath{\infty}이어야 한다. 위상동형은 연속인 전단사로 역함수도 연속인 사상이다. 두 좌표계 모음의 합집합도 위 조건을 만족하면 같은 매끄러운 구조를 정하는 것으로 본다.

\textbf{조건 확인 / Conditions.} Hausdorff·가산 기저·좌표의 역함수 연속성을 함께 기억한다.

\textbf{직접 선행 / Prerequisites.} \hyperref[def:06-01]{\texttt{06-01} 연속사상}; \hyperref[def:06-02]{\texttt{06-02} Hausdorff 공간}; \hyperref[def:06-03]{\texttt{06-03} 열린 기저}

\exampleheading
The fixed charts $z=x+iy$ and $w=1/z$ identify the underlying $S^2$ with a smooth real surface. On the overlap, $(x,y)\mapsto(x/(x^2+y^2),-y/(x^2+y^2))$ is smooth with smooth inverse. The Hausdorff and countable-basis conditions follow from the subspace topology in $\R^3$.
\end{atlasdefinition}

\begin{atlasdefinition}{\texttt{06-05}: 접벡터와 접공간 / Tangent vector and tangent space}
\label{def:06-05}
M을 매끄러운 n-다양체, x\ensuremath{\in}M을 고정한 점이라 하자. 어떤 \ensuremath{\varepsilon}>0에 대해 정의된 매끄러운 곡선 \(\gamma:(-\varepsilon,\varepsilon)\to M\), \(\gamma(0)=x\)를 생각한다. x의 좌표 근방 \(\varphi:U\to V\subset\mathbb R^n\)에서 두 곡선 \ensuremath{\gamma},\ensuremath{\eta}가 \((\varphi\circ\gamma)'(0)=(\varphi\circ\eta)'(0)\)을 만족하면 동치로 본다. 곡선은 0 근처로 제한하여 U 안에 있게 한다. 이 동치조건은 좌표 선택에 무관하다. 곡선의 동치류 (equivalence class)를 x에서의 접벡터 (tangent vector)라 하고 그 집합을 \(T_xM\)이라 쓴다. 전단사 \(T_xM\to\mathbb R^n,\ [\gamma]\mapsto(\varphi\circ\gamma)'(0)\)로 실수벡터공간의 덧셈과 스칼라곱을 옮긴 것이 접공간 (tangent space)의 벡터공간 구조다. 이 연산들도 좌표 선택에 무관하다.

\textbf{조건 확인 / Conditions.} 곡선의 점도 같고 좌표 속도도 같아야 한다.

\textbf{직접 선행 / Prerequisites.} \hyperref[def:06-04]{\texttt{06-04} 매끄러운 다양체}; \hyperref[def:B01]{\texttt{B01} 실수·복소 벡터공간}

\exampleheading
The curves $\gamma_a(t)=(\sin(at),0,\cos(at))$ in $S^2$ pass through $N$ and have velocity $(a,0,0)$. Thus $\gamma_1$ represents $e_1\in T_NS^2$. A reparametrization $\gamma_1(t+t^2)$ has the same velocity and hence the same tangent vector at $t=0$.
\end{atlasdefinition}

\begin{atlasdefinition}{\texttt{B07}: 정칙 접다발과 여접다발 / Holomorphic tangent and cotangent bundles}
\label{def:B07}
복소 n차원 다양체 X의 정칙 접다발 \(T^{1,0}X\)는 정칙좌표에서 \(\partial/\partial z^1,\ldots,\partial/\partial z^n\)을 기저로 하고 좌표변환의 복소 Jacobian에 의한 연쇄법칙으로 붙인 rank n 정칙 벡터다발이다. 각 섬유의 복소선형 쌍대를 취하면 정칙 여접다발이며 쌍대 기저는 \(dz^1,\ldots,dz^n\)이다. 정칙 벡터다발은 N08의 국소자명화 정의에서 섬유를 \(\mathbb C^n\)으로 바꾼 것이다.

\textbf{조건 확인 / Conditions.} 실수 접다발 TX와 복소 rank n의 T\ensuremath{{}^{1}}\ensuremath{{}^{0}}X를 구별한다. 여기서는 Jacobian으로 붙이는 구성을 사용한다.

\textbf{직접 선행 / Prerequisites.} \hyperref[def:B04]{\texttt{B04} 복소다양체}; \hyperref[def:06-05]{\texttt{06-05} 접벡터와 접공간}; \hyperref[def:B02]{\texttt{B02} 쌍대·텐서곱·외대수}

\exampleheading
On $\CP^1$, $\partial_w=-z^2\partial_z$ and $dw=-z^{-2}dz$. Pairing gives $dw(\partial_w)=1$ on the overlap. Thus the tangent and cotangent transition functions are reciprocal, and each bundle has complex rank one.
\end{atlasdefinition}

\begin{atlasdefinition}{\texttt{N23}: 복소사영공간 / Complex projective space}
\label{def:N23}
n\ensuremath{\ge}0에 대해 \(\mathbf{CP}^n\)은 \(\mathbb C^{n+1}\setminus\{0\}\)에서 \(Z\sim W\iff W=\lambda Z\)인 \(\lambda\in\mathbb C^*\)가 존재한다는 동치관계 (equivalence relation)의 몫이다. 원소를 \([Z_0:\cdots:Z_n]\)으로 쓴다. U\ensuremath{{}_{i}}=\{Z\ensuremath{{}_{i}}\ensuremath{\ne}0\}의 좌표는 \((Z_0/Z_i,\ldots,\widehat{Z_i/Z_i},\ldots,Z_n/Z_i)\in\mathbb C^n\)이며 이 좌표계로 복소다양체 (complex manifold) 구조를 준다.

\textbf{조건 확인 / Conditions.} 동차좌표는 하나의 비영 벡터가 아니라 비영 복소수 배까지의 동치류 (equivalence class)다.

\textbf{직접 선행 / Prerequisites.} \hyperref[def:B04]{\texttt{B04} 복소다양체}; \hyperref[def:B01]{\texttt{B01} 실수·복소 벡터공간}

\exampleheading
Every point of $\CP^1$ is a complex line $\C(S,T)\subset\C^2$, with $(S,T)\ne(0,0)$. For example $[1:2]=[3:6]$, while $[1:2]\ne[1:3]$. The point $[0:1]$ is omitted from the z-chart and is w=0 in the second chart.
\end{atlasdefinition}

\begin{atlasdefinition}{\texttt{06-06}: 여접공간 / Cotangent space}
\label{def:06-06}
M을 매끄러운 다양체 (smooth manifold), \(x\in M\)을 한 점, \(T_xM\)을 x에서의 실수 접공간 (tangent space)이라 하자. \(T_x^*M=\operatorname{Hom}_{\mathbb R}(T_xM,\mathbb R)\)로, 접공간의 실수 선형쌍대다.

\textbf{조건 확인 / Conditions.} 입력은 x에서의 접벡터 (tangent vector)다.

\textbf{직접 선행 / Prerequisites.} \hyperref[def:06-05]{\texttt{06-05} 접벡터와 접공간}; \hyperref[def:B02]{\texttt{B02} 쌍대·텐서곱·외대수}

\exampleheading
At $N\in S^2$, the covectors $\epsilon^1(a,b,0)=a$ and $\epsilon^2(a,b,0)=b$ form a basis of $T_N^*S^2$. Under our $z$-chart, $dx_N=\epsilon^1/2$ and $dy_N=\epsilon^2/2$, because $d\Phi_0(\partial_x)=2e_1$.
\end{atlasdefinition}

\begin{atlasdefinition}{\texttt{06-07}: 미분 q-형식 / Differential q-form}
\label{def:06-07}
M을 매끄러운 n-다양체, q를 양의 정수라 하자. 매끄러운 q-형식은 각 \(p\in M\)에 실수 다중선형사상 \(\omega_p:(T_pM)^q\to\mathbb R\)를 주는 족 \(\omega=(\omega_p)_{p\in M}\)으로, 다음 두 조건을 만족한다. \textbf{(1) 교대성.} 각 \(\omega_p\)는 두 입력을 교환하면 부호가 바뀐다. 다중선형성은 다른 입력들을 고정했을 때 각 입력에 대해 실수 선형이라는 뜻이다. \textbf{(2) 족의 매끄러움.} 임의의 매끄러운 좌표계 \(\varphi:U\to V\subset\mathbb R^n\)를 잡는다. 표준기저벡터 \(e_i\in\mathbb R^n\)에 대해 p에서의 좌표 접벡터 (tangent vector)를 \[E_i(p):=\left.\frac{d}{dt}\right|_{t=0}\varphi^{-1}(\varphi(p)+te_i)\in T_pM\]로 둔다. 이는 충분히 작은 t에 대해 정의되는 곡선의 속도다. 모든 지표 \(i_1,\ldots,i_q\in\{1,\ldots,n\}\)에 대해 함수 \[a_{i_1\cdots i_q}:V\longrightarrow\mathbb R,\]\[y\longmapsto\omega_{\varphi^{-1}(y)}\bigl(E_{i_1}(\varphi^{-1}(y)),\ldots,E_{i_q}(\varphi^{-1}(y))\bigr)\]가 \(C^\infty\), 즉 모든 차수의 편도함수가 존재하고 연속이어야 한다. 이것을 모든 매끄러운 좌표계에 대해 요구한다. \textbf{0차.} 매끄러운 0-형식은 함수 \(f:M\to\mathbb R\)로, 모든 매끄러운 좌표계에서 \(f\circ\varphi^{-1}:V\to\mathbb R\)가 \(C^\infty\)인 것이다. 매끄러운 q-형식들의 실수벡터공간을 \(\Omega^q(M)\)라 쓴다.

\textbf{조건 확인 / Conditions.} 점마다 교대 다중선형사상을 고르는 것만으로는 부족하다. 위에서 정의한 좌표 성분 함수들이 점에 따라 C\ensuremath{\infty}으로 변해야 한다.

\textbf{직접 선행 / Prerequisites.} \hyperref[def:06-06]{\texttt{06-06} 여접공간}

\exampleheading
The area form of the unit sphere is $\Omega_p(u,v)=p\cdot(u\times v)$ for $u,v\in T_pS^2$. This is alternating and depends smoothly on p. In the z-chart it is $\Omega=4(1+|z|^2)^{-2}\,dx\wedge dy$, so its coordinate coefficient is explicitly smooth.
\end{atlasdefinition}

\begin{atlasdefinition}{\texttt{B05}: 복소값 미분형식 / Complex-valued differential form}
\label{def:B05}
매끄러운 실수다양체 M에서 복소값 k-형식은 각 p에 교대 실수 k-다중선형사상 \(\omega_p:(T_pM)^k\to\mathbb C\)를 주되 모든 좌표 성분이 매끄러운 족이다. 유일하게 \(\omega=\alpha+i\beta\)로 쓰며 \ensuremath{\alpha},\ensuremath{\beta}는 실수값 k-형식이다. \(d\omega=d\alpha+i\,d\beta\)로 외미분을 확장한다. 0-형식은 매끄러운 복소값 함수이다.

\textbf{조건 확인 / Conditions.} 복소값이라는 말이 실수 접공간에 복소선형성을 요구한다는 뜻은 아니다. L값 형식은 N29에서 정의한다.

\textbf{직접 선행 / Prerequisites.} \hyperref[def:06-07]{\texttt{06-07} 미분 q-형식}

\exampleheading
On the z-chart of $\CP^1$, $dz=dx+i\,dy$ is a complex-valued real-linear one-form. It takes $\partial_x$ to $1$ and $\partial_y$ to $i$. The global form $i\Omega$ is a complex-valued two-form on all of $S^2$; its imaginary part is $\Omega$.
\end{atlasdefinition}

\begin{atlasdefinition}{\texttt{06-08}: 좌표 1-형식 / Coordinate one-form}
\label{def:06-08}
M을 매끄러운 n-다양체, \(\varphi=(x^1,\ldots,x^n):U\to V\subset\mathbb R^n\)를 매끄러운 좌표계라 하자. \(p\in U\), \(v=[\gamma]\in T_pM\)에 대해 \[dx^i_p:T_pM\to\mathbb R,\qquad dx^i_p(v):=(x^i\circ\gamma)'(0)\quad(1\le i\le n)\]로 정의한다. 여기서 \(\gamma(0)=p\)이고 \ensuremath{\gamma}는 0 근처에서 U 안에 있다. 접벡터 (tangent vector)의 동치관계 (equivalence relation)에 의해 이 값은 대표 곡선에 무관하다. 좌표 속도 \((\varphi\circ\gamma)'(0)=(v^1,\ldots,v^n)\)로 쓰면 \(dx^i_p(v)=v^i\)다. p\ensuremath{\in}U에 대한 이 선형함수들의 족이 U 위의 좌표 1-형식 \(dx^i\)다.

\textbf{조건 확인 / Conditions.} dx\ensuremath{{}^{i}}는 i번째 좌표 속도를 읽는 선형함수다.

\textbf{직접 선행 / Prerequisites.} \hyperref[def:06-07]{\texttt{06-07} 미분 q-형식}

\exampleheading
In the chart $z=x+iy$ of $\CP^1$, the real coordinate forms satisfy $dx(\partial_x)=1$, $dx(\partial_y)=0$, $dy(\partial_x)=0$, $dy(\partial_y)=1$. The complex coordinate form is $dz=dx+i\,dy$, and $dw=-z^{-2}dz$ on the overlap.
\end{atlasdefinition}

\begin{atlasdefinition}{\texttt{06-09}: Wedge 곱 / Wedge product / Exterior product}
\label{def:06-09}
M을 매끄러운 다양체 (smooth manifold), p,q\ensuremath{\ge}0을 정수, \(\alpha\in\Omega^p(M)\), \(\beta\in\Omega^q(M)\)라 하자. 외적 \(\alpha\wedge\beta\in\Omega^{p+q}(M)\)는 각 x\ensuremath{\in}M과 \(v_1,\ldots,v_{p+q}\in T_xM\)에서 \[(\alpha\wedge\beta)_x(v_1,\ldots,v_{p+q})=\frac{1}{p!q!}\sum_{\sigma\in S_{p+q}}\operatorname{sgn}(\sigma)\,\alpha_x(v_{\sigma(1)},\ldots,v_{\sigma(p)})\,\beta_x(v_{\sigma(p+1)},\ldots,v_{\sigma(p+q)})\]로 정의한다. \(S_{p+q}\)는 \{1,\ldots{},p+q\}의 순열들의 집합이고 \(\operatorname{sgn}(\sigma)\)는 순열의 부호다. 0-형식은 그 점에서의 함수값으로 취급하며 0!=1이다. 좌표계의 기본 형식은 \[(dx^{i_1}\wedge\cdots\wedge dx^{i_q})_x(v_1,\ldots,v_q)=\det\bigl(dx^{i_a}_x(v_b)\bigr)_{a,b=1}^q\]를 만족한다.

\textbf{조건 확인 / Conditions.} 행렬식의 행은 좌표 지표, 열은 입력 벡터에 대응한다.

\textbf{직접 선행 / Prerequisites.} \hyperref[def:06-07]{\texttt{06-07} 미분 q-형식}

\exampleheading
On the z-chart of $\CP^1$, $dz\wedge d\bar z=(dx+i\,dy)\wedge(dx-i\,dy)=-2i\,dx\wedge dy$. In particular $dz\wedge dz=0$, whereas the tensor $dz\otimes dz$ is nonzero. Wedge and tensor products are different operations.
\end{atlasdefinition}

\begin{atlasdefinition}{\texttt{06-10}: 외미분 / Exterior derivative}
\label{def:06-10}
M을 매끄러운 n-다양체, q\ensuremath{\ge}0을 정수라 하자. 외미분 (exterior derivative)은 실수 선형사상 \(d:\Omega^q(M)\to\Omega^{q+1}(M)\)으로 다음 국소식으로 정의한다. 매끄러운 좌표계 \(\varphi=(x^1,\ldots,x^n):U\to V\subset\mathbb R^n\)에서 \[\omega|_U=\sum_{I}(a_I\circ\varphi)\,dx^I,\qquad a_I:V\to\mathbb R\]로 쓴다. I는 \(1\le i_1<\cdots<i_q\le n\)을 만족하는 지표이고 \(dx^I=dx^{i_1}\wedge\cdots\wedge dx^{i_q}\), 각 a\ensuremath{{}_{i}}는 C\ensuremath{\infty} 함수다. V의 좌표를 \(y=(y_1,\ldots,y_n)\)이라 쓰면 \[(d\omega)|_U:=\sum_{j=1}^n\sum_I\left(\frac{\partial a_I}{\partial y_j}\circ\varphi\right)dx^j\wedge dx^I.\]q=0이면 빈 지표 I 하나를 사용하고 \(dx^\varnothing=1\)로 둔다. 각 좌표에서의 식은 겹침에서 일치하여 하나의 전역 (q+1)-형식을 정한다.

\textbf{조건 확인 / Conditions.} 새 dx\ensuremath{{}^{j}}는 dx\ensuremath{{}^{I}}의 앞에 놓인다.

\textbf{직접 선행 / Prerequisites.} \hyperref[def:06-08]{\texttt{06-08} 좌표 1-형식}; \hyperref[def:06-09]{\texttt{06-09} Wedge 곱}

\exampleheading
The global height function on $S^2$ is $H=(1-x^2-y^2)/(1+x^2+y^2)$ in the z-chart. Its differential is $dH=-4(x\,dx+y\,dy)/(1+x^2+y^2)^2$. Applying d again gives zero: the mixed partials cancel in the coefficient of $dx\wedge dy$.
\end{atlasdefinition}

\begin{atlasdefinition}{\texttt{B06}: 형식의 형식차수와 Dolbeault 연산자 / Bidegree and Dolbeault operators}
\label{def:B06}
복소 n차원 다양체에서 매끄러운 (p,q)-형식은 정칙좌표마다 \(\sum_{|I|=p,|J|=q}a_{IJ}\,dz^I\wedge d\bar z^J\)로 표현되는 복소값 형식이다. 계수는 매끄럽다. 외미분의 형식차수 (p+1,q) 성분을 \ensuremath{\partial}, (p,q+1) 성분을 \ensuremath{\bar\partial}라 한다. 따라서 \(d=\partial+\bar\partial\). 함수에 대해 \(\partial f=\sum_j\frac12(f_{x_j}-if_{y_j})dz^j\), \(\bar\partial f=\sum_j\frac12(f_{x_j}+if_{y_j})d\bar z^j\)이다.

\textbf{조건 확인 / Conditions.} (p,q)는 미분형식의 두 종류 입력에 대한 차수이다. 선다발의 degree와 구별한다.

\textbf{직접 선행 / Prerequisites.} \hyperref[def:B04]{\texttt{B04} 복소다양체}; \hyperref[def:B05]{\texttt{B05} 복소값 미분형식}; \hyperref[def:06-09]{\texttt{06-09} Wedge 곱}; \hyperref[def:06-10]{\texttt{06-10} 외미분}

\exampleheading
On the z-chart of $\CP^1$, take $f=z\bar z$. Then $\partial f=\bar z\,dz$, $\bar\partial f=z\,d\bar z$, and $\partial\bar\partial f=dz\wedge d\bar z$. These have types $(1,0),(0,1),(1,1)$ respectively. The function f is used only on this chart, not as a global function on the sphere.
\end{atlasdefinition}

\begin{atlasdefinition}{\texttt{06-11}: Closed 형식 / Closed differential form}
\label{def:06-11}
M을 매끄러운 다양체 (smooth manifold), \(q\ge0\)을 정수라 하자. \(\Omega^q(M)\)는 M 위 매끄러운 q-형식의 공간이며 \(d:\Omega^q(M)\to\Omega^{q+1}(M)\)는 외미분 (exterior derivative)이다. \(\alpha\in\Omega^q(M)\)가 closed라는 것은 \(d\alpha=0\)이라는 뜻이다.

\textbf{조건 확인 / Conditions.} 주어진 공간 위에서의 조건이다.

\textbf{직접 선행 / Prerequisites.} \hyperref[def:06-10]{\texttt{06-10} 외미분}

\exampleheading
The global area form $\Omega$ on $S^2$ is closed: $d\Omega$ has degree three, but the sphere has real dimension two, so every alternating three-form on its tangent spaces vanishes.
\end{atlasdefinition}

\begin{atlasdefinition}{\texttt{06-12}: Exact 형식 / Exact differential form}
\label{def:06-12}
매끄러운 다양체 (smooth manifold) M과 정수 \(q\ge1\)에 대해, \(\alpha\in\Omega^q(M)\)가 exact라는 것은 \(d\beta=\alpha\)인 \(\beta\in\Omega^{q-1}(M)\)가 존재한다는 뜻이다. 0차에서는 de Rham 복합체를 \(\Omega^{-1}(M)=0\), \(d^{-1}:0\to\Omega^0(M)\)으로 확장하는 관례에 따라, exact 0-형식을 \(\operatorname{im}d^{-1}=\{0\}\)의 원소로 정의한다.

\textbf{조건 확인 / Conditions.} \ensuremath{\beta}는 같은 공간 M 전체에 정의되어야 한다. 0차에서는 영함수만 exact다.

\textbf{직접 선행 / Prerequisites.} \hyperref[def:06-10]{\texttt{06-10} 외미분}

\exampleheading
The one-form $dH$ is globally exact on $S^2$, with global primitive H. The area form $\Omega$ is not exact because $\int_{S^2}\Omega=4\pi$, while the integral of a global exact two-form on a closed oriented surface is zero by Stokes.
\end{atlasdefinition}

\begin{atlasdefinition}{\texttt{06-13}: 원시형식 / Primitive of a differential form}
\label{def:06-13}
M을 매끄러운 다양체 (smooth manifold), \(q\ge1\)을 정수라 하자. \(\Omega^r(M)\)는 M 위 매끄러운 r-형식의 공간이고 d는 외미분 (exterior derivative)이다. \(\alpha\in\Omega^q(M)\)에 대해 \(d\beta=\alpha\)를 만족하는 \(\beta\in\Omega^{q-1}(M)\)를 \ensuremath{\alpha}의 원시형식 (primitive)이라 한다.

\textbf{조건 확인 / Conditions.} \ensuremath{\beta}의 차수는 \ensuremath{\alpha}보다 하나 낮다.

\textbf{직접 선행 / Prerequisites.} \hyperref[def:06-12]{\texttt{06-12} Exact 형식}

\exampleheading
For the form $\alpha=dH$ on $S^2$, both H and $H+7$ are global primitives. On the connected sphere any two function primitives differ by a constant. On $U_0\cap U_\infty\simeq\C^\times$, $dz/z$ has local logarithmic primitives but no global single-valued primitive.
\end{atlasdefinition}

\subsection{층과 국소 자료 / Sheaves and local data}

\begin{atlasdefinition}{\texttt{07-01}: 아벨군의 준층 / Presheaf of abelian groups}
\label{def:07-01}
X를 위상공간 (topological space)이라 하자. 각 열린 \(U\subset X\)에 아벨군 (abelian group) \(\mathcal F(U)\), 각 열린 포함 \(V\subset U\)에 준동형 (homomorphism) \(\rho_{UV}:\mathcal F(U)\to\mathcal F(V)\)를 주고 \[\rho_{UU}=\mathrm{id},\qquad\rho_{UW}=\rho_{VW}\circ\rho_{UV}\quad(W\subset V\subset U)\]를 만족시키는 자료다.

\textbf{조건 확인 / Conditions.} 제한은 큰 열린집합 (open set)에서 작은 열린집합 방향이다.

\textbf{직접 선행 / Prerequisites.} \hyperref[def:01-03]{\texttt{01-03} 열린덮개}; \hyperref[def:02-02]{\texttt{02-02} 아벨군 준동형}

\exampleheading
On $X=\CP^1$, assign to every open U the additive group $\mathcal O_X(U)$ of holomorphic functions. For $V\subset U$, use ordinary restriction. Restricting $z^2$ first from $U_0$ to $|z|<2$ and then to $|z|<1$ equals direct restriction. This is an abelian presheaf (and in fact a sheaf).
\end{atlasdefinition}

\begin{atlasdefinition}{\texttt{08-01}: 층 / Sheaf}
\label{def:08-01}
X를 위상공간 (topological space), \(\mathcal F\)를 X 위 아벨군 준층 (presheaf of abelian groups)이라 하자. \(\mathcal F\)가 층이라는 것은 모든 열린 \(U\subset X\), 모든 열린덮개 (open cover) \(U=\bigcup_{i\in I}U_i\), 모든 단면족 (family of sections) \((s_i\in\mathcal F(U_i))_{i\in I}\)에 대해 다음 조건이 성립하는 것이다. 모든 i,j\ensuremath{\in}I에서 \(s_i|_{U_i\cap U_j}=s_j|_{U_i\cap U_j}\)이면, 모든 i\ensuremath{\in}I에서 \(s|_{U_i}=s_i\)인 \(s\in\mathcal F(U)\)가 유일하게 존재한다. \(s|_V\)는 준층 (presheaf)의 제한사상 (restriction map)에 의한 상이다. 빈 덮개도 포함하므로 \(\mathcal F(\varnothing)=0\)이다.

\textbf{조건 확인 / Conditions.} 모든 열린 U의 모든 열린덮개 (open cover)에 대해 존재와 유일성을 요구한다.

\textbf{직접 선행 / Prerequisites.} \hyperref[def:07-01]{\texttt{07-01} 아벨군의 준층}

\exampleheading
On the two charts of $\CP^1$, take $f_0=3$ and $f_\infty=3$. Their restrictions agree, so the sheaf axiom glues them uniquely to the global constant function 3. Choosing $f_0=z$ and $f_\infty=0$ does not meet the overlap condition, so the axiom does not assert that these functions glue.
\end{atlasdefinition}

\begin{atlasdefinition}{\texttt{26-03}: 환층 / Sheaf of rings}
\label{def:26-03}
X를 위상공간 (topological space)이라 하자. 아래의 U,V는 X의 열린집합 (open set)이며 환은 단위원을 가진 가환환 (commutative ring)을 뜻한다. 각 열린 U의 단면 (section)집합 \(\mathcal O(U)\)가 환이고 모든 열린 포함 \(V\subset U\)의 제한이 환준동형이며 층의 유일성·붙임 공리를 만족하는 층이다.

\textbf{조건 확인 / Conditions.} 제한사상 (restriction map)은 단위원도 보존한다.

\textbf{직접 선행 / Prerequisites.} \hyperref[def:08-01]{\texttt{08-01} 층}; \hyperref[def:26-02]{\texttt{26-02} 환준동형}

\exampleheading
The holomorphic function sheaf $\mathcal O_X$ on $\CP^1$ is a sheaf of unital commutative rings: restriction preserves addition, multiplication, and the constant function 1. For example $(fg)|_V=(f|_V)(g|_V)$ for $V\subset U_0$.
\end{atlasdefinition}

\begin{atlasdefinition}{\texttt{07-02}: 단면 / Section}
\label{def:07-02}
X를 위상공간 (topological space), \(\mathcal F\)를 X 위 아벨군 준층 (presheaf of abelian groups), U를 X의 열린집합 (open set)이라 하자. \(\rho_{UV}:\mathcal F(U)\to\mathcal F(V)\)는 열린 포함 \(V\subset U\)에 주어진 제한사상 (restriction map)이다. 준층 (presheaf) \(\mathcal F\)에 대해 \(s\in\mathcal F(U)\)를 U 위 단면 (section)이라 하며, 열린 \(V\subset U\)에 대한 제한은 \(s|_V=\rho_{UV}(s)\)다.

\textbf{조건 확인 / Conditions.} 단면 (section)에는 정의된 열린집합 (open set) U가 따라다닌다.

\textbf{직접 선행 / Prerequisites.} \hyperref[def:07-01]{\texttt{07-01} 아벨군의 준층}

\exampleheading
The function $s(z)=e^z$ belongs to $\mathcal O_X(U_0)$ on $\CP^1$: it is a section over the chart $U_0$. No assertion that it is a section over all of X is being made; it has an essential singularity at infinity.
\end{atlasdefinition}

\begin{atlasdefinition}{\texttt{09-01}: Germ - 싹 / Germ}
\label{def:09-01}
X를 위상공간 (topological space), \(\mathcal F\)를 X 위 아벨군 준층 (presheaf of abelian groups), \(x\in X\)를 고정한 점이라 하자. 열린 근방 (open neighborhood) \(U\ni x\)와 단면 (section) \(s\in\mathcal F(U)\)의 쌍 \((U,s)\)들을 생각한다. \((U,s)\sim_x(V,t)\)는 어떤 열린 \(W\)가 존재하여 \(x\in W\subset U\cap V\)이고 \(s|_W=t|_W\)라는 뜻이다. \(s|_W\)는 준층 (presheaf)의 제한사상 (restriction map) \(\mathcal F(U)\to\mathcal F(W)\)에 의한 상이다. 이 동치관계 (equivalence relation)에 대한 \((U,s)\)의 동치류 (equivalence class) \([U,s]_x\)를 s의 x에서의 germ이라 하며 \(s_x\)로 쓴다.

\textbf{조건 확인 / Conditions.} 어떤 공통 열린 근방 (open neighborhood)에서 제한 단면 (section)들이 같아야 한다.

\textbf{직접 선행 / Prerequisites.} \hyperref[def:07-02]{\texttt{07-02} 단면}; \hyperref[def:01-02]{\texttt{01-02} 열린 근방}

\exampleheading
At $0\in U_0\subset\CP^1$, the function $1/(1-z)$ on $|z|<1/2$ and the sum $\sum_{n\ge0}z^n$ on $|z|<1$ define the same holomorphic germ. They agree on a smaller neighbourhood. The germ of $1$ is different despite having the same value at 0.
\end{atlasdefinition}

\begin{atlasdefinition}{\texttt{09-02}: Stalk - 줄기 / Stalk}
\label{def:09-02}
X를 위상공간 (topological space), \(\mathcal F\)를 X 위 아벨군 준층 (presheaf of abelian groups), \(x\in X\)라 하자. \(\mathcal F_x\)는 모든 열린 근방 (open neighborhood) \(U\ni x\)와 \(s\in\mathcal F(U)\)의 쌍 \((U,s)\)들을 다음 동치관계 (equivalence relation)로 나눈 집합이다. \((U,s)\sim_x(V,t)\)는 어떤 열린 \(W\)가 존재하여 \(x\in W\subset U\cap V\), \(s|_W=t|_W\)라는 뜻이다. 동치류 (equivalence class)를 \([U,s]_x\)라 쓰며 덧셈은 \[[U,s]_x+[V,t]_x=[U\cap V,\ s|_{U\cap V}+t|_{U\cap V}]_x\]로 정의한다. 영원소는 \([X,0]_x\), 덧셈 역원은 \(-[U,s]_x=[U,-s]_x\)다. 이 아벨군 (abelian group)을 x에서의 줄기 (stalk)라 하며 \(\mathcal F_x=\varinjlim_{x\in U}\mathcal F(U)\)로도 쓴다.

\textbf{조건 확인 / Conditions.} 모든 x의 열린 근방 (open neighborhood)을 사용하고 더 작은 근방에서 같아지는 표현들을 식별한다.

\textbf{직접 선행 / Prerequisites.} \hyperref[def:09-01]{\texttt{09-01} Germ - 싹}

\exampleheading
At $0\in\CP^1$, the stalk $\mathcal O_{X,0}$ is $\C\{z\}$, the convergent power-series ring. The elements $1$ and $1+z$ are distinct in this stalk; they become equal after passing to its residue field $\mathcal O_{X,0}/(z)$.
\end{atlasdefinition}

\begin{atlasdefinition}{\texttt{N03}: 환 달린 공간 / Ringed space}
\label{def:N03}
환 달린 공간은 위상공간 (topological space) X와 X 위 단위원을 가진 가환환 (commutative ring)의 층 \(\mathcal O_X\)의 쌍 \((X,\mathcal O_X)\)이다. 각 열린집합 (open set) U의 \(\mathcal O_X(U)\)는 환이고 제한사상 (restriction map)은 단위원을 보존한다. 이 층을 구조층 (structure sheaf)이라 한다. 모든 점 x의 줄기 (stalk) \(\mathcal O_{X,x}\)가 국소환 (local ring)이면 국소환 달린 공간 (locally ringed space)이라 한다.

\textbf{조건 확인 / Conditions.} 구조층 (structure sheaf)이라는 명칭만으로 그 단면 (section)이 매끄러운 함수인지 정칙함수 (holomorphic function)인지 정해지지 않는다.

\textbf{직접 선행 / Prerequisites.} \hyperref[def:26-03]{\texttt{26-03} 환층}; \hyperref[def:09-02]{\texttt{09-02} Stalk - 줄기}; \hyperref[def:28-03]{\texttt{28-03} 국소환}

\exampleheading
The analytic $\CP^1$ with its sheaf $\mathcal O_X$ is a locally ringed space. At a finite point a the stalk has maximal ideal $(z-a)$; at infinity it has maximal ideal $(w)$. The underlying topological sphere alone does not specify this holomorphic structure sheaf.
\end{atlasdefinition}

\begin{atlasdefinition}{\texttt{N04}: 복소다양체의 정칙 구조층 / Holomorphic structure sheaf}
\label{def:N04}
복소 n차원 다양체 X는 Hausdorff 제2가산 위상공간 (topological space)에 \(\mathbb C^n\)의 열린집합 (open set)으로 가는 좌표계들을 주되, 그 좌표계들이 X를 덮고 겹침에서 좌표변환이 정칙인 자료다. 그 정칙 구조층 (structure sheaf)은 \[\mathcal O_X(U)=\{f:U\to\mathbb C:\text{각 복소좌표에서 }f\text{가 정칙}\}\]이며 연산은 점별 덧셈·곱셈, 제한은 함수의 제한이다. 서로 호환되는 최대 복소 atlas를 같은 복소구조로 본다.

\textbf{조건 확인 / Conditions.} 정칙 구조층 (structure sheaf)과 매끄러운 함수층은 서로 다르다. 여기의 함수는 복소수값이다.

\textbf{직접 선행 / Prerequisites.} \hyperref[def:B04]{\texttt{B04} 복소다양체}; \hyperref[def:08-01]{\texttt{08-01} 층}; \hyperref[def:26-03]{\texttt{26-03} 환층}

\exampleheading
On the z-chart of $\CP^1$, $z^2$ is a section of the holomorphic structure sheaf but $\bar z$ is not. Both are sections of the smooth complex function sheaf. This tests which structure sheaf is being used.
\end{atlasdefinition}

\begin{atlasdefinition}{\texttt{10-01}: 층사상 / Morphism of sheaves}
\label{def:10-01}
X를 위상공간 (topological space), \(\mathcal F,\mathcal G\)를 X 위 아벨군 층 (sheaf of abelian groups)이라 하자. 층사상 (sheaf morphism) \(f:\mathcal F\to\mathcal G\)는 각 열린 \(U\subset X\)에 대한 준동형 (homomorphism) \(f_U:\mathcal F(U)\to\mathcal G(U)\)의 족으로, 모든 열린 \(V\subset U\)에 대해 \[\rho^{\mathcal G}_{UV}\circ f_U=f_V\circ\rho^{\mathcal F}_{UV}\]를 만족한다. \(\rho^{\mathcal F}_{UV}:\mathcal F(U)\to\mathcal F(V)\), \(\rho^{\mathcal G}_{UV}:\mathcal G(U)\to\mathcal G(V)\)는 각각의 층의 제한사상 (restriction map)이다.

\textbf{조건 확인 / Conditions.} 같은 공간 X 위의 두 층 사이에서 모든 열린 포함에 대한 호환성을 요구한다.

\textbf{직접 선행 / Prerequisites.} \hyperref[def:08-01]{\texttt{08-01} 층}; \hyperref[def:02-02]{\texttt{02-02} 아벨군 준동형}

\exampleheading
On $\CP^1$, differentiation gives a sheaf map of complex vector spaces $d:\mathcal O_X\to\mathcal K_X$, with $f\mapsto df$. It commutes with restriction. It is not $\mathcal O_X$-linear, since $d(fg)=f\,dg+g\,df$.
\end{atlasdefinition}

\begin{atlasdefinition}{\texttt{N05}: 구조층의 가군층 / Sheaf of modules}
\label{def:N05}
\((X,\mathcal O_X)\)를 환 달린 공간이라 하자. \(\mathcal O_X\)-가군층 (sheaf of modules) \(\mathcal E\)는 X 위 아벨군 층 (sheaf of abelian groups)이며, 모든 열린 U에 \(\mathcal O_X(U)\)-가군 구조가 주어져 모든 열린 \(V\subseteq U\), \(a\in\mathcal O_X(U),s\in\mathcal E(U)\)에 대해 \[(as)|_V=(a|_V)(s|_V)\]를 만족한다. 가군층 사상은 각 U에서 \(\mathcal O_X(U)\)-선형인 층사상 (sheaf morphism)이다.

\textbf{조건 확인 / Conditions.} 스칼라도 단면 (section)과 함께 제한해야 한다.

\textbf{직접 선행 / Prerequisites.} \hyperref[def:N01]{\texttt{N01} 가군}; \hyperref[def:N03]{\texttt{N03} 환 달린 공간}; \hyperref[def:10-01]{\texttt{10-01} 층사상}

\exampleheading
On $\CP^1$, $\mathcal K_X$ is an $\mathcal O_X$-module sheaf: locally $f\cdot(a\,dz)=(fa)\,dz$, and restricting a product gives the product of the restrictions. For example $z\cdot((1+z)dz)=(z+z^2)dz$.
\end{atlasdefinition}

\begin{atlasdefinition}{\texttt{07-04}: 제한층 / Restriction of a sheaf}
\label{def:07-04}
X를 위상공간 (topological space), \(\mathcal F\)를 X 위 아벨군 층 (sheaf of abelian groups)이라 하자. 열린 \(U\subset X\)에 대해 \(\mathcal F|_U\)는 U 위 층으로, 각 열린 \(V\subset U\)에서 \((\mathcal F|_U)(V)=\mathcal F(V)\)이며 제한사상 (restriction map)도 원래 것을 쓴다.

\textbf{조건 확인 / Conditions.} U는 X의 열린 부분집합이다.

\textbf{직접 선행 / Prerequisites.} \hyperref[def:07-01]{\texttt{07-01} 아벨군의 준층}

\exampleheading
Restrict the holomorphic function sheaf of $\CP^1$ to $U_0$. For every open $V\subset U_0$, $(\mathcal O_X|_{U_0})(V)=\mathcal O_X(V)$; the coordinate z identifies this restricted sheaf with the usual holomorphic function sheaf on $\C$.
\end{atlasdefinition}

\begin{atlasdefinition}{\texttt{N06}: 국소자유층 / Locally free sheaf}
\label{def:N06}
\((X,\mathcal O_X)\)를 환 달린 공간, r을 0 이상의 정수라 하자. \(\mathcal O_X\)-가군층 (sheaf of modules) \(\mathcal E\)가 rank r의 국소자유층 (locally free sheaf)이라는 것은 모든 \(x\in X\)에 대해 열린 근방 (open neighborhood) U와 \(\mathcal O_X|_U\)-가군층 동형 \[\mathcal E|_U\cong(\mathcal O_X|_U)^{\oplus r}\]이 존재한다는 뜻이다. 동형은 U의 모든 열린 부분집합의 단면 (section)과 제한에 호환된다.

\textbf{조건 확인 / Conditions.} 한 점의 줄기 (stalk)가 자유라는 진술만으로 정의를 대체하지 않는다.

\textbf{직접 선행 / Prerequisites.} \hyperref[def:N05]{\texttt{N05} 구조층의 가군층}; \hyperref[def:07-04]{\texttt{07-04} 제한층}; \hyperref[def:13-02]{\texttt{13-02} 동형사상}

\exampleheading
On $\CP^1$, $\mathcal K_X|_{U_0}=\mathcal O_{U_0}\,dz$ and $\mathcal K_X|_{U_\infty}=\mathcal O_{U_\infty}\,dw$. It is locally free of rank one, but not globally free: a global basis would be a nowhere-zero holomorphic one-form, and no nonzero global holomorphic one-form exists on $\CP^1$.
\end{atlasdefinition}

\begin{atlasdefinition}{\texttt{N07}: 가역층 / Invertible sheaf}
\label{def:N07}
\((X,\mathcal O_X)\)를 국소환 (local ring) 달린 공간이라 하자. 이 공간 위 가역층 (invertible sheaf)은 rank 1의 국소자유 \(\mathcal O_X\)-가군층 (sheaf of modules) \(\mathcal L\)이다. 즉 모든 \(x\in X\)에 열린 근방 (open neighborhood) U가 존재하여 \(\mathcal L|_U\cong\mathcal O_X|_U\)인 가군층 동형이 있다. 국소 생성원 e를 고르면 모든 열린 \(V\subseteq U\)의 단면 (section)은 \(ae|_V\), \(a\in\mathcal O_X(V)\)로 유일하게 표현된다.

\textbf{조건 확인 / Conditions.} 국소적으로 자유인 것과 전역적으로 \(\mathcal O_X\)와 동형인 것은 다른 조건이다. 여기서는 구조층 (structure sheaf)의 모든 줄기 (stalk)가 국소환 (local ring)이라고 가정한다.

\textbf{직접 선행 / Prerequisites.} \hyperref[def:N06]{\texttt{N06} 국소자유층}; \hyperref[def:B02]{\texttt{B02} 쌍대·텐서곱·외대수}

\exampleheading
On $\CP^1$, evaluation gives $\mathcal K_X\otimes_{\mathcal O_X}\mathcal T_X\simeq\mathcal O_X$, where $\mathcal T_X$ is the sheaf of holomorphic tangent vectors. Locally $dz\otimes\partial_z\mapsto1$; the reciprocal transition functions make these maps agree. Thus $\mathcal K_X$ is invertible.
\end{atlasdefinition}

\begin{atlasdefinition}{\texttt{07-03}: 전역단면 / Global section}
\label{def:07-03}
X는 위상공간 (topological space)이고 층의 계수는 아벨군 (abelian group)으로 한다. \(X\) 위 층 \(\mathcal F\)의 전역단면군은 \(\Gamma(X,\mathcal F)=\mathcal F(X)\)다.

\textbf{조건 확인 / Conditions.} 공간 X 전체 위의 단면 (section)이다.

\textbf{직접 선행 / Prerequisites.} \hyperref[def:07-02]{\texttt{07-02} 단면}

\exampleheading
For $X=\CP^1$, $\Gamma(X,\mathcal O_X)=\C$, represented by constant holomorphic functions. In contrast $\Gamma(X,\mathcal K_X)=0$. Locally $dz$ is a frame, but it is not a global holomorphic one-form.
\end{atlasdefinition}

\begin{atlasdefinition}{\texttt{10-02}: 줄기사상 / Induced map on stalks}
\label{def:10-02}
X를 위상공간 (topological space), \(\mathcal F,\mathcal G\)를 X 위 아벨군 층 (sheaf of abelian groups), \(f:\mathcal F\to\mathcal G\)를 층사상 (sheaf morphism), \(x\in X\)를 고정한 점이라 하자. 각 열린 \(U\subset X\)에서 f의 성분은 준동형 (homomorphism) \(f_U:\mathcal F(U)\to\mathcal G(U)\)이며, 모든 열린 \(V\subset U\)와 \(s\in\mathcal F(U)\)에 대해 \(f_V(s|_V)=f_U(s)|_V\)를 만족한다. 줄기 (stalk) \(\mathcal F_x\)의 원소는 \(x\in U\), \(s\in\mathcal F(U)\)인 쌍 \((U,s)\)의 동치류 (equivalence class) \([U,s]_x\)다. 두 쌍 \((U,s),(V,t)\)는 어떤 열린 \(W\)가 존재하여 \(x\in W\subset U\cap V\), \(s|_W=t|_W\)일 때 동치다. \(\mathcal G_x\)도 같은 방식으로 정의한다. f가 x에서 유도하는 줄기사상 (induced map on stalks)은 준동형 \[f_x:\mathcal F_x\longrightarrow\mathcal G_x,\qquad[U,s]_x\longmapsto[U,f_U(s)]_x\]이다. 제한과의 호환성에 의해 동치인 대표쌍 (representative pair)은 동치인 상을 주므로 이 대응은 대표쌍의 선택에 무관하다. \(s_x:=[U,s]_x\)로 줄여 쓰면 \(f_x(s_x)=(f_U(s))_x\)다.

\textbf{조건 확인 / Conditions.} U는 x의 열린 근방 (open neighborhood)이고 s는 F(U)의 원소다. 입력과 출력은 각각 F\ensuremath{{}_{x}}와 G\ensuremath{{}_{x}}의 동치류 (equivalence class)다.

\textbf{직접 선행 / Prerequisites.} \hyperref[def:10-01]{\texttt{10-01} 층사상}; \hyperref[def:09-02]{\texttt{09-02} Stalk - 줄기}

\exampleheading
At $0\in\CP^1$, the preceding map induces $d_0:\mathcal O_{X,0}\to\mathcal K_{X,0}$, $[f]_0\mapsto[df]_0$. Concretely $[z^3+2z]_0$ maps to $[(3z^2+2)\,dz]_0$. Agreement on a neighbourhood implies agreement of derivatives on that neighbourhood, so representatives do not change the output.
\end{atlasdefinition}

\begin{atlasdefinition}{\texttt{10-03}: 층의 완전열 / Exact sequence of sheaves}
\label{def:10-03}
X를 위상공간 (topological space)이라 하자. X 위 아벨군 층 (sheaf of abelian groups)들과 층사상 (sheaf morphism)으로 이루어진 수열 \(\cdots\to\mathcal F^{r-1}\xrightarrow{a^{r-1}}\mathcal F^r\xrightarrow{a^r}\mathcal F^{r+1}\to\cdots\)이 \(\mathcal F^r\)에서 완전하다는 것은 모든 x\ensuremath{\in}X에서 \[\operatorname{im}(a_x^{r-1}:\mathcal F_x^{r-1}\to\mathcal F_x^r)=\ker(a_x^r:\mathcal F_x^r\to\mathcal F_x^{r+1})\]가 성립한다는 뜻이다. \(\mathcal F_x^r\)는 x에서의 줄기 (stalk)이고 \(a_x^r\)는 층사상으로부터 유도된 줄기 준동형 (homomorphism on stalks)이다. 수열이 완전하다는 것은 모든 항에서 완전하다는 뜻이다.

\textbf{조건 확인 / Conditions.} 모든 열린집합 (open set)의 단면열 (sequence of sections) 완전성을 요구하는 것이 아니다.

\textbf{직접 선행 / Prerequisites.} \hyperref[def:10-02]{\texttt{10-02} 줄기사상}; \hyperref[def:03-01]{\texttt{03-01} 한 자리에서의 완전성}

\exampleheading
On the analytic $\CP^1$, $0\to\underline\C\to\mathcal O_X\xrightarrow d\mathcal K_X\to0$ is exact as a sequence of complex-vector-space sheaves. The kernel consists of locally constant functions. On a small coordinate disc every holomorphic $a(z)\,dz$ has a holomorphic primitive, proving surjectivity on stalks.
\end{atlasdefinition}

\begin{atlasdefinition}{\texttt{10-04}: 층의 단사 / Monomorphism of sheaves}
\label{def:10-04}
X를 위상공간 (topological space), \(\mathcal F,\mathcal G\)를 X 위 아벨군 층 (sheaf of abelian groups), \(f:\mathcal F\to\mathcal G\)를 층사상 (sheaf morphism)이라 하자. 각 x\ensuremath{\in}X에서 \(\mathcal F_x,\mathcal G_x\)는 줄기 (stalk)이고 \(f_x:\mathcal F_x\to\mathcal G_x\)는 f가 유도하는 줄기 준동형 (homomorphism on stalks)이다. f가 단사라는 조건은 모든 x\ensuremath{\in}X와 모든 \(a,b\in\mathcal F_x\)에 대해 \(f_x(a)=f_x(b)\Rightarrow a=b\)라는 것이다.

\textbf{조건 확인 / Conditions.} 모든 점의 줄기 (stalk)에서 검사한다.

\textbf{직접 선행 / Prerequisites.} \hyperref[def:10-02]{\texttt{10-02} 줄기사상}

\exampleheading
The inclusion $\underline\C\hookrightarrow\mathcal O_X$ on $\CP^1$ is injective on every stalk: a locally constant germ that becomes the zero holomorphic germ is already zero. This supplies a concrete sheaf monomorphism.
\end{atlasdefinition}

\begin{atlasdefinition}{\texttt{10-05}: 층의 전사 / Epimorphism of sheaves}
\label{def:10-05}
X를 위상공간 (topological space), \(\mathcal F,\mathcal G\)를 X 위 아벨군 층 (sheaf of abelian groups), \(f:\mathcal F\to\mathcal G\)를 층사상 (sheaf morphism)이라 하자. 각 x\ensuremath{\in}X에서 \(\mathcal F_x,\mathcal G_x\)는 줄기 (stalk)이고 \(f_x:\mathcal F_x\to\mathcal G_x\)는 f가 유도하는 줄기 준동형 (homomorphism on stalks)이다. f가 전사라는 조건은 모든 x\ensuremath{\in}X와 모든 \(b\in\mathcal G_x\)에 대해 어떤 \(a\in\mathcal F_x\)가 존재하여 \(f_x(a)=b\)라는 것이다.

\textbf{조건 확인 / Conditions.} 각 열린 U의 단면사상 (map on sections) f\ensuremath{{}_{U}}가 전사라는 뜻이 아니다.

\textbf{직접 선행 / Prerequisites.} \hyperref[def:10-02]{\texttt{10-02} 줄기사상}

\exampleheading
The derivative $d:\mathcal O_X\to\mathcal K_X$ on $\CP^1$ is surjective on stalks. It need not be surjective on sections of every open set: on $U_0\cap U_\infty=\C^\times$, the form $dz/z$ has no global holomorphic primitive because its integral around the unit circle is $2\pi i$.
\end{atlasdefinition}

\begin{atlasdefinition}{\texttt{11-01}: 국소상수함수 / Locally constant function}
\label{def:11-01}
U를 위상공간 (topological space), B를 집합이라 하자. 함수 \(f:U\to B\)가 국소상수라는 것은 모든 \(x\in U\)에 대해 \(x\in V\subset U\)인 열린 근방 (open neighborhood) V가 존재하여 \(f|_V\)가 상수라는 뜻이다.

\textbf{조건 확인 / Conditions.} U 전체에서 같은 상수일 필요는 없다.

\textbf{직접 선행 / Prerequisites.} \hyperref[def:01-02]{\texttt{01-02} 열린 근방}

\exampleheading
On the disconnected open subset $U=\{Z>1/2\}\cup\{Z<-1/2\}$ of $S^2$, define f=1 on the north cap and f=2 on the south cap. It is locally constant, but not constant on U. On the connected whole sphere a locally constant function is constant.
\end{atlasdefinition}

\begin{atlasdefinition}{\texttt{11-02}: 상수층 / Constant sheaf}
\label{def:11-02}
X를 위상공간 (topological space), B를 아벨군 (abelian group)이라 하자. 각 열린 U는 X의 열린집합 (open set)이다. 함수 \(g:U\to B\)가 국소상수라는 것은 모든 x\ensuremath{\in}U가 g가 상수인 열린 근방 (open neighborhood)을 갖는다는 뜻이다. 아벨군 B의 상수층 (constant sheaf) \(\underline B\)는 각 열린 U에 국소상수함수 (locally constant function) \(U\to B\)들의 군을 주고, 함수의 제한을 제한사상 (restriction map)으로 쓰는 층이다. 덧셈은 점별 덧셈이다.

\textbf{조건 확인 / Conditions.} 상수함수만이 아니라 국소상수함수 (locally constant function) 모두다.

\textbf{직접 선행 / Prerequisites.} \hyperref[def:11-01]{\texttt{11-01} 국소상수함수}; \hyperref[def:08-01]{\texttt{08-01} 층}

\exampleheading
The constant sheaf $\underline\Z$ on $S^2$ assigns locally constant integer-valued functions to U. On the two-cap open set from the preceding example its sections form $\Z\oplus\Z$, whereas on $S^2$ its sections form $\Z$. This is not the presheaf assigning a single copy of $\Z$ to every nonempty open set.
\end{atlasdefinition}

\begin{atlasdefinition}{\texttt{11-03}: 미분형식층 / Sheaf of differential forms}
\label{def:11-03}
q는 0 이상의 정수이며 \(\Omega^q(U)\)는 U 위 매끄러운 q-형식의 실수벡터공간이다. 매끄러운 다양체 (smooth manifold) M 위의 미분형식층 \(\mathcal A^q\)는 열린 \(U\subset M\)에 \(\mathcal A^q(U)=\Omega^q(U)\)를 주고 형식의 제한을 쓰는 층이다.

\textbf{조건 확인 / Conditions.} A\ensuremath{{}^{0}}는 매끄러운 함수층이며 실수 상수층 (constant sheaf)과 다르다.

\textbf{직접 선행 / Prerequisites.} \hyperref[def:06-07]{\texttt{06-07} 미분 q-형식}; \hyperref[def:08-01]{\texttt{08-01} 층}

\exampleheading
The sheaf $\mathcal A^2$ on $S^2$ assigns smooth two-forms to each open U. Its global section $\Omega_p(u,v)=p\cdot(u\times v)$ restricts in the z-chart to $4(1+|z|^2)^{-2}dx\wedge dy$. Its local coordinate formulas agree by pullback on overlaps.
\end{atlasdefinition}

\begin{atlasdefinition}{\texttt{27-03}: 정칙함수층 / Sheaf of holomorphic functions}
\label{def:27-03}
리만곡면 (Riemann surface) X의 열린 U에 대해 \(\mathcal O_X(U)\)는 U 위 정칙함수 (holomorphic function)들의 환이다. 정칙성은 복소좌표에서 검사하며 제한은 함수의 제한이다.

\textbf{조건 확인 / Conditions.} germ과 stalk는 근방에서의 실제 일치로 정의한다.

\textbf{직접 선행 / Prerequisites.} \hyperref[def:27-02]{\texttt{27-02} 리만곡면}; \hyperref[def:26-03]{\texttt{26-03} 환층}

\exampleheading
On $\CP^1$, $\mathcal O_X(U_0)=\mathcal O(\C)$ contains $e^z$, but $\mathcal O_X(X)=\C$. This distinguishes the sheaf, its sections on a chart, and its global sections in one explicit example.
\end{atlasdefinition}

\begin{atlasdefinition}{\texttt{28-04}: 정칙함수의 줄기 / Stalk of the holomorphic function sheaf}
\label{def:28-04}
X는 리만곡면 (Riemann surface)이고 x는 X의 고정한 점이다. 리만곡면 X에서 \(\mathcal O_{X,x}\)는 x의 근방에서 정의된 정칙함수 (holomorphic function)들의 germ으로 이루어진 환이다. 두 함수는 x의 어떤 공통 열린 근방 (open neighborhood)에서 같을 때 같은 germ이며 연산은 공통 근방에서 더하고 곱해 정의한다.

\textbf{조건 확인 / Conditions.} 함수값 하나로 germ을 식별하지 않는다.

\textbf{직접 선행 / Prerequisites.} \hyperref[def:27-03]{\texttt{27-03} 정칙함수층}; \hyperref[def:09-02]{\texttt{09-02} Stalk - 줄기}

\exampleheading
Using z at $0\in\CP^1$, the stalk of the holomorphic function sheaf consists of convergent germs $\sum_{n\ge0}a_nz^n$. The representative $1/(1-z)$ corresponds to the germ $1+z+z^2+\cdots$, with positive radius of convergence.
\end{atlasdefinition}

\begin{atlasdefinition}{\texttt{28-05}: 정칙 줄기의 극대아이디얼 / Maximal ideal of the holomorphic local ring}
\label{def:28-05}
X를 리만곡면 (Riemann surface), x\ensuremath{\in}X를 고정한 점이라 하자. \(\mathcal O_{X,x}\)는 x 근방의 정칙함수 (holomorphic function)들의 germ으로 이루어진 환이다. U가 x의 열린 근방 (open neighborhood)이고 \(f\in\mathcal O_X(U)\)이면 그 germ을 \([U,f]_x\)라 쓴다. \[\mathfrak m_x:=\{[U,f]_x\in\mathcal O_{X,x}:f(x)=0\}\]로 정의한다. x의 어떤 공통 근방에서 같은 함수들은 x에서 같은 값을 가지므로 이 조건은 대표 함수에 무관하다.

\textbf{조건 확인 / Conditions.} 대표 함수의 x에서의 값이 0이라는 조건이다.

\textbf{직접 선행 / Prerequisites.} \hyperref[def:28-04]{\texttt{28-04} 정칙함수의 줄기}; \hyperref[def:28-02]{\texttt{28-02} 극대아이디얼}

\exampleheading
At $0\in\CP^1$, the maximal ideal of holomorphic germs is $\mathfrak m_0=\{f:f(0)=0\}=(z)$. The residue map takes $a_0+a_1z+\cdots$ to $a_0$, so the residue field is $\C$.
\end{atlasdefinition}

\begin{atlasdefinition}{\texttt{28-06}: 수렴급수환 / Ring of convergent power series}
\label{def:28-06}
t를 한 변수라 하자. 수렴급수환 (ring of convergent power series) \(\mathbb C\{t\}\)는 복소 계수 거듭제곱급수 \(\sum_{n\ge0}a_nt^n\) 중 어떤 r>0이 존재하여 모든 \(z\in\mathbb C\), |z|<r에 대해 \(\sum_{n\ge0}a_nz^n\)이 수렴하는 것들의 집합이다. 두 급수는 모든 계수가 같을 때 같은 원소다. 덧셈은 계수별이고, \((\sum a_nt^n)(\sum b_nt^n)\)의 n차 계수는 \(\sum_{i=0}^n a_i b_{n-i}\)로 정의한다. 단위원은 상수급수 1이다.

\textbf{조건 확인 / Conditions.} 급수마다 양의 수렴 반경이 존재해야 한다.

\textbf{직접 선행 / Prerequisites.} \hyperref[def:27-01]{\texttt{27-01} 정칙함수}; \hyperref[def:26-01]{\texttt{26-01} 환}

\exampleheading
The analytic local ring at $0\in\CP^1$ is $\C\{z\}$. The series $\sum_{n\ge0}z^n$ represents a germ, but $\sum_{n\ge0}n!z^n$ does not: its radius of convergence is zero.
\end{atlasdefinition}

\begin{atlasdefinition}{\texttt{28-07}: 형식급수환 / Formal power series ring}
\label{def:28-07}
t를 형식적 변수라 하자. \(\mathbb C[[t]]\)의 원소는 복소수열 \((a_n)_{n\ge0}\)이며 \(\sum_{n\ge0}a_nt^n\)으로 쓴다. \(\sum a_nt^n=\sum b_nt^n\)은 모든 n\ensuremath{\ge}0에서 a\ensuremath{{}_{n}}=b\ensuremath{{}_{n}}라는 뜻이다. 덧셈의 n차 계수는 a\ensuremath{{}_{n}}+b\ensuremath{{}_{n}}, 곱셈의 n차 계수는 \(\sum_{i=0}^n a_i b_{n-i}\)로 정의한다. 영원소는 모든 계수가 0인 급수, 단위원은 상수급수 1이다. 수렴 조건을 부과하지 않는다.

\textbf{조건 확인 / Conditions.} 수렴 조건을 요구하지 않는다.

\textbf{직접 선행 / Prerequisites.} \hyperref[def:26-01]{\texttt{26-01} 환}

\exampleheading
The formal series $\sum_{n\ge0}n!z^n$ is an element of $\C[[z]]$, the completed formal coordinate ring at $0\in\CP^1$. Its existence as a formal element makes no assertion of a holomorphic function near 0.
\end{atlasdefinition}

\begin{atlasdefinition}{\texttt{28-08}: 아이디얼의 거듭제곱 / Power of an ideal}
\label{def:28-08}
R을 단위원을 가진 가환환 (commutative ring), J\ensuremath{\subset}R을 아이디얼 (ideal), m\ensuremath{\ge}1을 정수라 하자. \[J^m:=\left\{\sum_{k=1}^N a_{k,1}\cdots a_{k,m}:N\ge0,\ a_{k,j}\in J\right\}\]로 정의한다. N=0인 합은 0이다. 이는 J의 원소 m개를 곱한 모든 곱으로 생성된 아이디얼이다. 특히 리만곡면 (Riemann surface)의 점 x에서 \(R=\mathcal O_{X,x}\), \(J=\mathfrak m_x\)로 두면 \(\mathfrak m_x^m\)을 얻는다.

\textbf{조건 확인 / Conditions.} 한 원소의 m제곱들만을 말하는 것이 아니다.

\textbf{직접 선행 / Prerequisites.} \hyperref[def:28-01]{\texttt{28-01} 아이디얼}

\exampleheading
At $0\in\CP^1$, $\mathfrak m_0^k=(z^k)\subset\C\{z\}$. The germ $z^3+z^4$ lies in $\mathfrak m_0^3$ but not $\mathfrak m_0^4$. This records a vanishing order of exactly three.
\end{atlasdefinition}

\begin{atlasdefinition}{\texttt{28-09}: 제트 / Jet}
\label{def:28-09}
X를 리만곡면 (Riemann surface), x\ensuremath{\in}X를 고정한 점, m\ensuremath{\ge}1을 정수라 하자. \(\mathcal O_{X,x}\)는 x에서의 정칙함수 (holomorphic function) germ의 환이고 \(\mathfrak m_x\)는 x에서 값이 0인 germ들의 아이디얼 (ideal)이다. 차수 m-1의 함수 jet는 몫환 (quotient ring) \(\mathcal O_{X,x}/\mathfrak m_x^m\)의 원소다. 즉 germ g의 jet는 잉여류 (coset) \(g+\mathfrak m_x^m\)이며, 두 germ g,h는 \(g-h\in\mathfrak m_x^m\)일 때 같은 jet를 나타낸다. \(\mathfrak m_x^m\)은 그 아이디얼의 원소 m개를 곱한 곱들의 유한합으로 이루어진다.

\textbf{조건 확인 / Conditions.} jet 차수는 m-1이고 아이디얼 (ideal) 지수는 m이다.

\textbf{직접 선행 / Prerequisites.} \hyperref[def:28-08]{\texttt{28-08} 아이디얼의 거듭제곱}; \hyperref[def:28-04]{\texttt{28-04} 정칙함수의 줄기}

\exampleheading
The second jet of $e^z$ at $0\in\CP^1$ is its class in $\mathcal O_{X,0}/\mathfrak m_0^3$: $[e^z]=[1+z+z^2/2]$. Terms of degree at least three are forgotten, while value, first derivative and second derivative remain.
\end{atlasdefinition}

\begin{atlasdefinition}{\texttt{28-10}: 완비화 / Completion}
\label{def:28-10}
X를 리만곡면 (Riemann surface), x\ensuremath{\in}X를 고정한 점이라 하자. \(R:=\mathcal O_{X,x}\)는 x에서의 정칙함수 (holomorphic function) germ의 환이고 \(\mathfrak m:=\{g\in R:g(x)=0\}\)이다. 정수 m\ensuremath{\ge}1에 대해 \(R_m:=R/\mathfrak m^m\)으로 놓고 \[\pi_m:R_{m+1}\to R_m,\qquad g+\mathfrak m^{m+1}\mapsto g+\mathfrak m^m\]으로 정의한다. \(\mathfrak m^{m+1}\subset\mathfrak m^m\)이므로 이 사상은 대표원 (representative) 선택에 무관하다. \(\mathfrak m\)-진 완비화 (completion)는 \[\widehat R:=\left\{(j_m)_{m\ge1}\in\prod_{m\ge1}R_m:\pi_m(j_{m+1})=j_m\text{ for every }m\ge1\right\}\]에 성분별 덧셈과 곱셈을 준 환이다. 이를 \(\varprojlim_{m\ge1}R/\mathfrak m^m\)으로 쓴다.

\textbf{조건 확인 / Conditions.} 서로 호환되는 열이며 수렴하는 함수라는 조건은 없다.

\textbf{직접 선행 / Prerequisites.} \hyperref[def:28-08]{\texttt{28-08} 아이디얼의 거듭제곱}

\exampleheading
At $0\in\CP^1$, the compatible truncations $1+z+\cdots+z^k$ modulo $z^{k+1}$ determine the element $\sum_{n\ge0}z^n$ in $\widehat{\mathcal O}_{X,0}\simeq\C[[z]]$. The truncations of $\sum n!z^n$ also give an element of this completion even though no convergent germ has those coefficients.
\end{atlasdefinition}

\subsection{선다발과 인자 / Line bundles and divisors}

\begin{atlasdefinition}{\texttt{N08}: 정칙 선다발과 국소자명화 / Holomorphic line bundle and local trivialization}
\label{def:N08}
X를 복소다양체 (complex manifold)라 하자. 정칙 선다발 (holomorphic line bundle)은 복소다양체 L, 정칙사상 \(\pi:L\to X\), 각 fiber \(L_x=\pi^{-1}(x)\)의 1차원 복소벡터공간 구조와 다음 국소자명화 (local trivialization)로 이루어진다. X의 열린덮개 (open cover) \((U_i)\)에 대해 쌍정칙사상 \(\tau_i:\pi^{-1}(U_i)\to U_i\times\mathbb C\)가 존재하고 \(\operatorname{pr}_1\tau_i=\pi\)이며 각 fiber에서 선형이다. 자명 선다발 (line bundle)은 사영 \(X\times\mathbb C\to X\)이다.

\textbf{조건 확인 / Conditions.} fiber마다 C와 동형이라는 것만으로는 부족하다. 열린 근방 (open neighborhood) 전체에서 정칙인 자명화가 필요하다.

\textbf{직접 선행 / Prerequisites.} \hyperref[def:B04]{\texttt{B04} 복소다양체}; \hyperref[def:B01]{\texttt{B01} 실수·복소 벡터공간}

\exampleheading
The canonical bundle on $\CP^1$ is trivialized over $U_0$ by $(p,a\,dz_p)\mapsto(p,a)$ and over $U_\infty$ by $(p,b\,dw_p)\mapsto(p,b)$. Both trivializations are holomorphic and linear on fibers; their change is induced by $dw=-z^{-2}dz$.
\end{atlasdefinition}

\begin{atlasdefinition}{\texttt{N09}: 다발의 단면과 국소 틀 / Section and local frame}
\label{def:N09}
\(\pi:L\to X\)를 정칙 선다발 (holomorphic line bundle), U를 열린집합 (open set)이라 하자. U 위 정칙 단면 (section)은 정칙사상 \(s:U\to L\)로 \(\pi\circ s=\operatorname{id}_U\)를 만족하는 것이다. 모든 x\ensuremath{\in}U에서 \(e(x)\ne0\in L_x\)인 정칙 단면 e를 국소 틀 (local frame)이라 한다. 이때 정칙 단면 s는 \(s=fe\), \(f\in\mathcal O_X(U)\)로 유일하게 표현된다. U마다 정칙 단면들을 모아 제한과 점별 연산을 주면 \(\mathcal O_X\)-가군층 (sheaf of modules)을 얻는다.

\textbf{조건 확인 / Conditions.} s(x)는 X의 점 x가 아니라 그 위 fiber의 벡터다.

\textbf{직접 선행 / Prerequisites.} \hyperref[def:N08]{\texttt{N08} 정칙 선다발과 국소자명화}; \hyperref[def:N04]{\texttt{N04} 복소다양체의 정칙 구조층}

\exampleheading
On $U_0\subset\CP^1$, $e=dz$ is a local frame of $K_X$, and $s=(z^2+1)dz$ is a holomorphic section. The section s is not a frame on all of $U_0$, because it vanishes at $z=i$ and $z=-i$.
\end{atlasdefinition}

\begin{atlasdefinition}{\texttt{N10}: 전이함수와 성분의 변환 / Transition function}
\label{def:N10}
X 위 정칙 선다발 (holomorphic line bundle) L의 국소 틀 (local frame) \(e_i\)를 열린덮개 (open cover) \((U_i)\) 위에서 잡는다. 전이함수 (transition function)의 관례를 \[e_j=g_{ij}e_i,\qquad g_{ij}\in\mathcal O_X^*(U_i\cap U_j)\]로 정한다. \(\mathcal O_X^*(U)\)는 U에서 영점이 없는 정칙함수 (holomorphic function)들의 곱셈군이다. 삼중 교집합에서 \(g_{ik}=g_{ij}g_{jk}\), \(g_{ii}=1\)이다. 단면 (section)을 \(s=f_ie_i=f_je_j\)로 쓰면 \(f_i=g_{ij}f_j\)이다.

\textbf{조건 확인 / Conditions.} 틀의 전이와 성분의 전이는 역방향이다. 관례 e\ensuremath{{}_{j}}=g\ensuremath{{}_{i}}\ensuremath{{}_{j}}e\ensuremath{{}_{i}}를 고정한다.

\textbf{직접 선행 / Prerequisites.} \hyperref[def:N09]{\texttt{N09} 다발의 단면과 국소 틀}

\exampleheading
For $K_{\CP^1}$ take $e_0=dz$, $e_\infty=dw$. The frame transition is $g_{0\infty}=-z^{-2}$. If $s=a(z)dz=b(w)dw$, the coefficient relation is $a(z)=-z^{-2}b(1/z)$, or $b(w)=-w^{-2}a(1/w)$.
\end{atlasdefinition}

\begin{atlasdefinition}{\texttt{N11}: 선다발의 동형 / Isomorphism of line bundles}
\label{def:N11}
\(\pi_L:L\to X\), \(\pi_E:E\to X\)를 정칙 선다발 (holomorphic line bundle)이라 하자. 그 동형은 쌍정칙사상 \(\Phi:L\to E\)로 \(\pi_E\Phi=\pi_L\)이며 모든 x\ensuremath{\in}X에서 \(\Phi|_{L_x}:L_x\to E_x\)가 복소선형 동형인 것이다. L이 자명하다는 것은 자명 선다발 (line bundle) \(X\times\mathbb C\)와 이런 동형이 존재한다는 뜻이다.

\textbf{조건 확인 / Conditions.} 같은 기저공간 X 위에서 fiber별 선형성과 정칙성을 모두 요구한다.

\textbf{직접 선행 / Prerequisites.} \hyperref[def:N08]{\texttt{N08} 정칙 선다발과 국소자명화}

\exampleheading
The bundle associated with $\mathcal O_X(-2[\infty])$ is isomorphic to $K_{\CP^1}$ by multiplication by the meromorphic section dz. Its local frames $1,w^2$ map to $dz,-dw$, respectively. Every fiber map is complex-linear and invertible, and the maps cover the identity of $\CP^1$.
\end{atlasdefinition}

\begin{atlasdefinition}{\texttt{N12}: 쌍대 선다발 / Dual line bundle}
\label{def:N12}
X 위 정칙 선다발 (holomorphic line bundle) L의 쌍대는 각 x의 fiber가 \(L_x^*=\operatorname{Hom}_{\mathbb C}(L_x,\mathbb C)\)인 선다발 (line bundle) \(L^*\)이다. 국소 틀 (local frame) e\ensuremath{{}_{i}}에 대해 \(e_i^*(e_i)=1\)인 쌍대 틀을 잡아 정칙구조를 정한다. \(e_j=g_{ij}e_i\)이면 \(e_j^*=g_{ij}^{-1}e_i^*\)로 붙인다.

\textbf{조건 확인 / Conditions.} 복소선형 쌍대이며, 복소켤레 선형 함수들의 공간과 구별한다.

\textbf{직접 선행 / Prerequisites.} \hyperref[def:N08]{\texttt{N08} 정칙 선다발과 국소자명화}; \hyperref[def:B02]{\texttt{B02} 쌍대·텐서곱·외대수}

\exampleheading
For $K_{\CP^1}$, the dual bundle is $T^{1,0}\CP^1$. Since $dw=-z^{-2}dz$, the dual frame transforms as $\partial_w=-z^2\partial_z$. Evaluating the two dual frames gives 1 in either chart.
\end{atlasdefinition}

\begin{atlasdefinition}{\texttt{N13}: 유리형 함수의 위수 / Order of a meromorphic function}
\label{def:N13}
X를 연결 리만곡면 (Riemann surface)이라 하자. 유리형 함수 (meromorphic function)는 각 점 근방에서 정칙함수 (holomorphic function)의 몫으로 표현되며, 분모는 그 근방에서 항등적으로 0이 아니어야 한다. 항등적으로 0이 아닌 유리형 함수 f와 점 p에서 \(t(p)=0\)인 국소좌표를 잡으면 \(f=t^m u\), \(m\in\mathbb Z\), u는 p 근방의 정칙함수이고 \(u(p)\ne0\)로 유일한 m을 갖게 쓸 수 있다. \(\operatorname{ord}_p(f):=m\)으로 정의한다. m>0은 영점, m<0은 극점, m=0은 정칙이고 0이 아닌 경우다.

\textbf{조건 확인 / Conditions.} 0함수는 여기의 위수 정의에서 제외한다. 극점의 위수는 음수다.

\textbf{직접 선행 / Prerequisites.} \hyperref[def:27-02]{\texttt{27-02} 리만곡면}; \hyperref[def:28-03]{\texttt{28-03} 국소환}

\exampleheading
The rational function $f=z^3/(z-2)$ on $\CP^1$ has orders 3 at 0 and $-1$ at 2. At infinity, $f=w^{-2}/(1-2w)$, so its order there is $-2$. At every other point its order is zero.
\end{atlasdefinition}

\begin{atlasdefinition}{\texttt{N14}: 곡선 위의 인자 / Divisor on a compact Riemann surface}
\label{def:N14}
X를 연결 compact 리만곡면 (Riemann surface)이라 하자. X 위 divisor는 유한 형식합 \[D=\sum_{p\in X}n_p[p],\qquad n_p\in\mathbb Z,\quad n_p=0\text{ except finitely many }p\]이다. 덧셈은 계수별이다. 모든 p에서 n\ensuremath{{}_{p}}\ensuremath{\ge}0이면 유효 인자 (effective divisor)라 한다. \(D\ge E\)는 모든 p에서 D의 계수가 E의 계수 이상이라는 뜻이다.

\textbf{조건 확인 / Conditions.} 이 항목은 매끄러운 compact 복소곡선에 한정한다. 일반 고차원 공간의 divisor와 같은 정의로 취급하지 않는다.

\textbf{직접 선행 / Prerequisites.} \hyperref[def:27-02]{\texttt{27-02} 리만곡면}

\exampleheading
The formal sum $D=2[0]-[\infty]$ is a divisor on $\CP^1$, with support $\{0,\infty\}$. It is not effective because its infinity coefficient is negative. The divisor $2[0]+[2]$ is effective.
\end{atlasdefinition}

\begin{atlasdefinition}{\texttt{N15}: 주인자 / Principal divisor}
\label{def:N15}
X를 연결 compact 리만곡면 (Riemann surface), f를 항등적으로 0이 아닌 전역 유리형 함수 (meromorphic function)라 하자. \[\operatorname{div}(f):=\sum_{p\in X}\operatorname{ord}_p(f)[p]\]를 f의 주인자 (principal divisor)라 한다. 이 합은 유한하며 영점의 차수는 양수, 극점의 차수는 음수로 들어간다. 이런 형태의 divisor들을 주인자라고 한다.

\textbf{조건 확인 / Conditions.} 국소 함수 하나의 영점·극점 기록이 아니라 전역 유리형 함수 (meromorphic function)의 divisor다.

\textbf{직접 선행 / Prerequisites.} \hyperref[def:N13]{\texttt{N13} 유리형 함수의 위수}; \hyperref[def:N14]{\texttt{N14} 곡선 위의 인자}

\exampleheading
On $\CP^1$, the meromorphic function z has principal divisor $\operatorname{div}(z)=[0]-[\infty]$. Its zero at 0 and pole at infinity both have order one; its divisor degree is zero.
\end{atlasdefinition}

\begin{atlasdefinition}{\texttt{N16}: 인자의 선형동치 / Linear equivalence of divisors}
\label{def:N16}
연결 compact 리만곡면 (Riemann surface) X의 두 divisor D,E가 선형동치 (linear equivalence)라는 것은 항등적으로 0이 아닌 전역 유리형 함수 (meromorphic function) f가 존재하여 \[D-E=\operatorname{div}(f)\]인 것이다. \(D\sim E\)로 쓴다. Divisor들의 군을 주인자 (principal divisor)들의 부분군 (subgroup)으로 나눈 몫군 (quotient group)이 divisor class group이다.

\textbf{조건 확인 / Conditions.} 같은 degree라는 것만으로 선형동치 (linear equivalence)가 되는 것은 일반적으로 아니다.

\textbf{직접 선행 / Prerequisites.} \hyperref[def:N14]{\texttt{N14} 곡선 위의 인자}; \hyperref[def:N15]{\texttt{N15} 주인자}

\exampleheading
On $\CP^1$, $[0]$ and $[\infty]$ are linearly equivalent because their difference is $\operatorname{div}(z)$. Also $-2[0]$ and $-2[\infty]$ differ by $\operatorname{div}(z^{-2})$. Equality of divisors is stronger than linear equivalence.
\end{atlasdefinition}

\begin{atlasdefinition}{\texttt{N17}: 인자의 차수 / Degree of a divisor}
\label{def:N17}
연결 compact 리만곡면 (Riemann surface) X 위 \(D=\sum_p n_p[p]\)의 차수는 \[\deg D:=\sum_p n_p\in\mathbb Z\]이다. 여기서는 모든 점의 잉여체가 C이므로 잉여체 확장차수의 가중치는 1이다.

\textbf{조건 확인 / Conditions.} degree는 하나의 정수로 정보를 합친 것이다. divisor 자체나 그 동치류 (equivalence class)와 구별한다.

\textbf{직접 선행 / Prerequisites.} \hyperref[def:N14]{\texttt{N14} 곡선 위의 인자}

\exampleheading
On $\CP^1$, $D=3[0]-[2]-2[\infty]$ has degree $3-1-2=0$, while $\operatorname{div}(dz)=-2[\infty]$ has degree $-2$. The first is a function divisor; the second is a one-form's divisor.
\end{atlasdefinition}

\begin{atlasdefinition}{\texttt{N18}: 인자에 대응하는 가역층 / The invertible sheaf O(D)}
\label{def:N18}
연결 compact 리만곡면 (Riemann surface) X와 divisor \(D=\sum_p n_p[p]\)를 잡는다. \(\mathcal O_X(D)(U)\)는 U 위 유리형 함수 (meromorphic function) f 중 모든 p\ensuremath{\in}U에서 \(\operatorname{ord}_p(f)+n_p\ge0\)인 것들의 집합이며 0도 포함한다. U의 어떤 연결성분에서 f가 0이면 그 성분에서는 조건을 자동으로 만족한 것으로 정한다. 덧셈·정칙함수 (holomorphic function) 곱·제한은 함수의 연산이다. p만 D의 지지에 들어오는 작은 좌표근방에서 \(t(p)=0\)이면 이 층은 \(t^{-n_p}\mathcal O_X\)로 표현된다.

\textbf{조건 확인 / Conditions.} 양의 계수 n\ensuremath{{}_{p}}는 f에 차수 n\ensuremath{{}_{p}}까지의 극점을 허용한다. \(\mathcal O(D)\)와 \(\mathcal O(-D)\)의 부호를 구별한다.

\textbf{직접 선행 / Prerequisites.} \hyperref[def:N05]{\texttt{N05} 구조층의 가군층}; \hyperref[def:N07]{\texttt{N07} 가역층}; \hyperref[def:N13]{\texttt{N13} 유리형 함수의 위수}; \hyperref[def:N14]{\texttt{N14} 곡선 위의 인자}

\exampleheading
For $D=[\infty]$ on $\CP^1$, the analytic sheaf $\mathcal O_X(D)$ has local frame 1 on $U_0$ and $w^{-1}$ near infinity. A global section is a meromorphic function with at most a simple pole at infinity and no other poles; the global section space is $\C\langle1,z\rangle$. As a section of this line bundle, 1 has a zero at infinity because $1=w\cdot w^{-1}$.
\end{atlasdefinition}

\begin{atlasdefinition}{\texttt{N19}: 표준다발과 표준층 / Canonical bundle and canonical sheaf}
\label{def:N19}
X를 복소 n차원 다양체라 하자. 표준선다발 (canonical line bundle)은 \(K_X:=\bigwedge^n(T^{1,0}X)^*\)이고, 그 정칙 단면층을 표준층 (canonical sheaf)이라 한다. 국소 정칙좌표 z\ensuremath{{}^{1}},\ldots{},z\ensuremath{{}^{n}}에서 틀은 \(dz^1\wedge\cdots\wedge dz^n\)이며 모든 정칙 단면 (section)은 그 틀에 정칙함수 (holomorphic function)를 곱한 것이다. 좌표변환 w=w(z)에서 \[dw^1\wedge\cdots\wedge dw^n=\det\!\left(\frac{\partial w^i}{\partial z^j}\right)dz^1\wedge\cdots\wedge dz^n.\]여기서 \(T^{1,0}X\)는 정칙좌표의 \(\partial/\partial z^j\)를 틀로 갖는 정칙 접다발이다.

\textbf{조건 확인 / Conditions.} 곡선에서는 K\ensuremath{{}_{x}}가 정칙 1-형식의 다발이다. 접다발 자체와 구별한다.

\textbf{직접 선행 / Prerequisites.} \hyperref[def:B07]{\texttt{B07} 정칙 접다발과 여접다발}; \hyperref[def:B02]{\texttt{B02} 쌍대·텐서곱·외대수}; \hyperref[def:N08]{\texttt{N08} 정칙 선다발과 국소자명화}; \hyperref[def:N10]{\texttt{N10} 전이함수와 성분의 변환}

\exampleheading
For $X=\CP^1$, the complex dimension is one, so $K_X=(T^{1,0}X)^*$. The local frames $dz,dw$ satisfy $dw=-z^{-2}dz$. The meromorphic section dz has a second-order pole at infinity and no other zero or pole, giving $K_X\simeq\mathcal O(-2)$.
\end{atlasdefinition}

\begin{atlasdefinition}{\texttt{N20}: 유리형 단면의 인자와 표준인자 / Divisor of a meromorphic section}
\label{def:N20}
X를 연결 compact 리만곡면 (Riemann surface), L을 정칙 선다발 (holomorphic line bundle)이라 하자. L의 유리형 단면 (section) s는 국소 정칙 틀 e\ensuremath{{}_{i}}에서 \(s=f_ie_i\)이며 f\ensuremath{{}_{i}}가 유리형이고 겹침에서 \(f_i=g_{ij}f_j\)인 자료다. s가 항등적으로 0이 아니면 \(\operatorname{ord}_p(s):=\operatorname{ord}_p(f_i)\)와 \(\operatorname{div}(s)=\sum_p\operatorname{ord}_p(s)[p]\)로 정의한다. 전이함수 (transition function) g\ensuremath{{}_{i}}\ensuremath{{}_{j}}가 정칙이고 영점이 없으므로 위수는 틀의 선택에 무관하다. L=K\ensuremath{{}_{x}}일 때 이런 divisor를 표준인자 (canonical divisor)라 한다.

\textbf{조건 확인 / Conditions.} 표준인자 (canonical divisor)는 선택한 유리형 단면 (section)에 따라 달라진다. 그 선형동치 (linear equivalence)류가 표준적으로 정해진다.

\textbf{직접 선행 / Prerequisites.} \hyperref[def:N10]{\texttt{N10} 전이함수와 성분의 변환}; \hyperref[def:N13]{\texttt{N13} 유리형 함수의 위수}; \hyperref[def:N14]{\texttt{N14} 곡선 위의 인자}; \hyperref[def:N19]{\texttt{N19} 표준다발과 표준층}

\exampleheading
On $\CP^1$, $dz=-w^{-2}dw$ gives $\operatorname{div}(dz)=-2[\infty]$. The different meromorphic section $dz/z=-dw/w$ has divisor $-[0]-[\infty]$. Their difference is the principal divisor of $1/z$, so they give the same canonical divisor class.
\end{atlasdefinition}

\begin{atlasdefinition}{\texttt{N21}: 사영직선의 O(n) 선다발 / The line bundle O(n) on the projective line}
\label{def:N21}
\(\mathbf{CP}^1\)의 열린집합 (open set) \(U_0=\{Z_0\ne0\}\), \(U_\infty=\{Z_1\ne0\}\)에 좌표 \(z=Z_1/Z_0\), \(w=Z_0/Z_1=1/z\)를 둔다. 정수 n에 대해 \(\mathcal O(n)\)은 두 자명 선다발 (line bundle)의 틀을 겹침에서 \[e_\infty=z^n e_0\]로 붙여 얻는 정칙 선다발 (holomorphic line bundle)이다. 동명 기호로 그 정칙 단면층도 쓴다. 전역 단면 (section)의 성분은 \(f_0(z)=z^n f_\infty(1/z)\)를 만족한다.

\textbf{조건 확인 / Conditions.} 여기의 전이 관례를 고정해야 n의 부호를 정확히 읽을 수 있다.

\textbf{직접 선행 / Prerequisites.} \hyperref[def:N10]{\texttt{N10} 전이함수와 성분의 변환}; \hyperref[def:N23]{\texttt{N23} 복소사영공간}

\exampleheading
On $\CP^1$, $\mathcal O(n)=\mathcal O(n[\infty])$ has frames $e_0=1$ and $e_\infty=w^{-n}$, so $e_\infty=z^n e_0$. A global section for $n\ge0$ is represented on $U_0$ by a polynomial of degree at most n. For n=2, the basis is $1,z,z^2$; for n=-2 there are no nonzero global holomorphic sections.
\end{atlasdefinition}

\subsection{계량·연결·곡률 / Metrics, connections and curvature}

\begin{atlasdefinition}{\texttt{N28}: 선다발의 Hermitian 계량 / Hermitian metric on a line bundle}
\label{def:N28}
X 위 정칙 선다발 (holomorphic line bundle) L의 Hermitian 계량은 각 fiber L\ensuremath{{}_{x}}의 양의 정부호 Hermitian 형식 \(h_x:L_x\times L_x\to\mathbb C\)의 족이다. 첫 인자에 선형, 둘째에 켤레선형인 관례를 쓴다. 매끄럽다는 것은 모든 국소 정칙 틀 e에서 \(h_e(x):=h_x(e(x),e(x))\)가 매끄러운 양의 실수함수라는 뜻이다. \(e_j=g_{ij}e_i\)이면 \(h_{e_j}=|g_{ij}|^2h_{e_i}\)여야 한다.

\textbf{조건 확인 / Conditions.} 틀마다 임의의 양의 함수를 고르는 것만으로는 부족하다. 전이함수 (transition function)와의 호환조건이 필요하다.

\textbf{직접 선행 / Prerequisites.} \hyperref[def:N08]{\texttt{N08} 정칙 선다발과 국소자명화}; \hyperref[def:N10]{\texttt{N10} 전이함수와 성분의 변환}; \hyperref[def:B03]{\texttt{B03} Hermitian 내적}

\exampleheading
For $K_{\CP^1}$ induced from the fixed unit-sphere metric, put $h_z=\|dz\|^2=(1+|z|^2)^2/2$. Then $h_w=|z|^{-4}h_z=(1+|w|^2)^2/2$. Both functions are positive and smooth, including w=0, and obey the frame transition law.
\end{atlasdefinition}

\begin{atlasdefinition}{\texttt{N29}: 복소 connection / Complex connection}
\label{def:N29}
복소다양체 (complex manifold) X 위 정칙 선다발 (holomorphic line bundle) L에서, 매끄러운 단면 (section)을 매끄러운 L값 복소 1-형식으로 보내는 복소선형 국소 연산 \(\nabla:\mathcal A^0(L)\to\mathcal A^1(L)\)이 connection이라는 것은 모든 열린 U, \(f\in C^\infty(U,\mathbb C)\), 매끄러운 단면 s에 대해 \[\nabla(fs)=df\otimes s+f\nabla s\]이며 열린 부분집합으로의 제한과 호환된다는 뜻이다. \(\mathcal A^k(L)(U)\)는 \(\bigwedge^kT^*_{\mathbb R}U\otimes_{\mathbb R}\mathbb C\otimes_{\mathbb C}L|_U\)의 매끄러운 단면이다. 국소 틀 (local frame) e에서 \(\nabla e=A_e\otimes e\)로 쓴다.

\textbf{조건 확인 / Conditions.} \(\nabla\)는 매끄러운 함수에 대해 선형이 아니라 Leibniz 법칙을 만족한다.

\textbf{직접 선행 / Prerequisites.} \hyperref[def:N09]{\texttt{N09} 다발의 단면과 국소 틀}; \hyperref[def:B05]{\texttt{B05} 복소값 미분형식}; \hyperref[def:06-10]{\texttt{06-10} 외미분}

\exampleheading
In the holomorphic frame dz of $K_{\CP^1}$, take $A_z=2\bar z(1+|z|^2)^{-1}dz$. The connection acts by $\nabla(f\,dz)=(df+fA_z)\otimes dz$. The corresponding form in the other frame is $A_w=2\bar w(1+|w|^2)^{-1}dw$; it satisfies $A_w=A_z+g^{-1}dg$ for $g=-z^{-2}$.
\end{atlasdefinition}

\begin{atlasdefinition}{\texttt{B08}: 정칙 선다발의 Dolbeault 연산자 / Dolbeault operator on a holomorphic line bundle}
\label{def:B08}
정칙 선다발 L의 정칙 국소 기저 e에 대해 \(\bar\partial_L(fe)=(\bar\partial f)\otimes e\)로 정의한다. 이는 매끄러운 단면을 매끄러운 L값 (0,1)-형식으로 보내는 제한과 호환되는 복소선형 연산이다. e는 정칙이고 f는 매끄러운 복소값 함수이다. 연결 \ensuremath{\nabla}의 (0,1) 성분은 그 L값 1-형식 출력의 (0,1) 성분을 취한 연산이다.

\textbf{조건 확인 / Conditions.} 기저를 임의의 매끄러운 기저로 바꾸어 이 식을 그대로 쓰지 않는다. 정칙 전이함수의 \ensuremath{\bar\partial}가 0이어서 이 정의가 붙는다.

\textbf{직접 선행 / Prerequisites.} \hyperref[def:N09]{\texttt{N09} 다발의 단면과 국소 틀}; \hyperref[def:B06]{\texttt{B06} 형식의 형식차수와 Dolbeault 연산자}

\exampleheading
On the z-chart of $\CP^1$, the smooth section $s=\bar z\,dz$ of $K_X$ satisfies $\bar\partial_Ks=d\bar z\otimes dz$, which is nonzero. The section $z^2dz$ instead has $\bar\partial_K(z^2dz)=0$. A smooth section and a holomorphic section are thus distinguished explicitly.
\end{atlasdefinition}

\begin{atlasdefinition}{\texttt{B09}: 선다발 연결의 곡률 / Curvature of a line bundle connection}
\label{def:B09}
연결 \ensuremath{\nabla}를 L값 형식에 \(\nabla(\alpha\otimes s)=d\alpha\otimes s+(-1)^k\alpha\wedge\nabla s\)로 확장한다. 여기서 \ensuremath{\alpha}는 복소값 k-형식이다. 단면 s에 대해 \(\nabla^2s=F_\nabla\otimes s\)를 만족하는 전역 복소값 2-형식 \(F_\nabla\)를 곡률이라 한다. 선다발의 섬유 자기준동형은 자연스럽게 복소 스칼라이므로 이렇게 표현된다. 국소 기저에서 \(\nabla e=A_e\otimes e\)이면 \(F_\nabla=dA_e\)이다.

\textbf{조건 확인 / Conditions.} A\ensuremath{{}_{e}}는 국소 기저에 의존하고 곡률은 전역 형식이다. rank 1이므로 A\ensuremath{{}_{e}}\ensuremath{\wedge}A\ensuremath{{}_{e}}=0이다.

\textbf{직접 선행 / Prerequisites.} \hyperref[def:N29]{\texttt{N29} 복소 connection}; \hyperref[def:06-09]{\texttt{06-09} Wedge 곱}

\exampleheading
For the preceding connection on $K_{\CP^1}$, differentiation gives $F_K=dA_z=-2(1+|z|^2)^{-2}dz\wedge d\bar z$. Hence $\nabla^2(f\,dz)=fF_K\otimes dz$. The expression in the w-chart is identical after coordinate substitution.
\end{atlasdefinition}

\begin{atlasdefinition}{\texttt{N30}: Chern connection과 곡률 / Chern connection and curvature}
\label{def:N30}
정칙 선다발 (holomorphic line bundle) L과 Hermitian 계량 h에서 Chern connection은 \(\nabla^{0,1}=\bar\partial_L\)이고 계량과 호환되는 connection이다. 계량 호환은 모든 실수 벡터장 V와 매끄러운 단면 (section) s,t에 대해 \[Vh(s,t)=h(\nabla_Vs,t)+h(s,\nabla_Vt)\]인 것이다. \(\bar\partial_L(fe)=(\bar\partial f)\otimes e\)는 정칙 틀 e로 정의한다. 여기서 \(\bar\partial f=\sum_j\frac12(\partial_{x_j}+i\partial_{y_j})f\,d\bar z^j\), \(\partial f=\sum_j\frac12(\partial_{x_j}-i\partial_{y_j})f\,dz^j\)이다. 정칙 틀에서는 \(A_e=\partial\log h_e\), 곡률은 \(F_\nabla=dA_e=-\partial\bar\partial\log h_e\)로 주어진다.

\textbf{조건 확인 / Conditions.} 곡률의 관례는 \(\nabla^2=F_\nabla\)이다. \(\frac{i}{2\pi}F_\nabla\)를 실수 Chern 형식으로 사용하는 관례와 함께 고정한다.

\textbf{직접 선행 / Prerequisites.} \hyperref[def:N28]{\texttt{N28} 선다발의 Hermitian 계량}; \hyperref[def:N29]{\texttt{N29} 복소 connection}; \hyperref[def:B08]{\texttt{B08} 정칙 선다발의 Dolbeault 연산자}; \hyperref[def:B09]{\texttt{B09} 선다발 연결의 곡률}

\exampleheading
The Hermitian metric $h_z=(1+|z|^2)^2/2$ on $K_{\CP^1}$ gives the Chern connection $A_z=\partial\log h_z=2\bar z(1+|z|^2)^{-1}dz$. Its curvature is $F_K=-2(1+|z|^2)^{-2}dz\wedge d\bar z$, a global (1,1)-form. The factor $1/2$ in h disappears when differentiating its logarithm.
\end{atlasdefinition}

\begin{atlasdefinition}{\texttt{B10}: 정준선다발의 Chern 곡률 / Chern curvature of the canonical line bundle}
\label{def:B10}
리만곡면 X의 정준선다발 \(K_X\)에 매끄러운 Hermitian 계량 h를 고른다. 이 쌍의 Chern 연결을 \(\nabla^{\mathrm{Ch}}\)라 할 때 곡률 \(F_{K_X,h}:=F_{\nabla^{\mathrm{Ch}}}\)이다. 정칙좌표 z의 기저 dz에 대해 \(h_z=h(dz,dz)\)라 쓰면 \[F_{K_X,h}=-\partial\bar\partial\log h_z.\]이는 전역 복소값 (1,1)-형식이다.

\textbf{조건 확인 / Conditions.} K\ensuremath{{}_{x}}만으로 하나의 곡률형식이 정해지지는 않는다. h를 지정해야 한다. 가우스 곡률이라는 실수함수와 동일한 대상이 아니다.

\textbf{직접 선행 / Prerequisites.} \hyperref[def:N19]{\texttt{N19} 표준다발과 표준층}; \hyperref[def:N30]{\texttt{N30} Chern connection과 곡률}

\exampleheading
For the canonical bundle of $\CP^1$ with the fixed sphere-induced metric, $F_{K,h}=-2(1+|z|^2)^{-2}dz\wedge d\bar z=i\Omega$. Thus $\frac{i}{2\pi}\int_{\CP^1}F_{K,h}=-\frac1{2\pi}\int_{S^2}\Omega=-2$. This uses the identification of the sphere with $\CP^1$ and its outward orientation.
\end{atlasdefinition}

\begin{atlasdefinition}{\texttt{B11}: 첫째 Chern 형식 / First Chern form}
\label{def:B11}
정칙 선다발 L과 매끄러운 Hermitian 계량 h의 Chern 곡률을 \(F_{L,h}\)라 할 때 첫째 Chern 형식은 \[c_1(L,h)=\frac{i}{2\pi}F_{L,h}\]이다. 이 규약에서 이는 닫힌 실수값 (1,1)-형식이다.

\textbf{조건 확인 / Conditions.} i/(2\ensuremath{\pi})와 1/(2\ensuremath{\pi}i)는 부호가 다르다. 형식 자체와 그 코호몰로지류를 구별한다.

\textbf{직접 선행 / Prerequisites.} \hyperref[def:N30]{\texttt{N30} Chern connection과 곡률}

\exampleheading
For $K_{\CP^1}$, $c_1(K,h)=\frac{i}{2\pi}F_K=-\Omega/(2\pi)$ and its integral is $-2$. For the dual tangent line bundle, the form is $+\Omega/(2\pi)$ and its integral is $+2$. The normalization and the choice of bundle both affect the sign.
\end{atlasdefinition}

\begin{atlasdefinition}{\texttt{B12}: 실수 첫째 Chern 류 / Real first Chern class}
\label{def:B12}
정칙 선다발 L의 실수 첫째 Chern 류는 Hermitian 계량 h를 골라 \[c_1(L)_{\mathbb R}:=[c_1(L,h)]\in H^2_{\mathrm{dR}}(X;\mathbb R)\]로 정의한다. 다른 계량은 같은 류를 준다. 이 항목에서는 정수 계수 위상적 Chern 류가 아니라 그 실수 de Rham 표현을 다룬다.

\textbf{조건 확인 / Conditions.} 계량 독립성은 정의의 정당화에 필요한 사실이다. 연결 정리 T09에서 형식의 차이가 완전형식임을 따로 명시한다.

\textbf{직접 선행 / Prerequisites.} \hyperref[def:B11]{\texttt{B11} 첫째 Chern 형식}; \hyperref[def:06-14]{\texttt{06-14} de Rham 코호몰로지}

\exampleheading
Let $u=[\Omega/(4\pi)]\in H^2_{\mathrm{dR}}(S^2;\R)$, normalized by $\int_{S^2}u=1$. Under $S^2\simeq\CP^1$, $c_1(K)_{\R}=-2u$ and $c_1(T^{1,0}\CP^1)_{\R}=2u$. A change of Hermitian metric changes a representative form by an exact form, not these classes.
\end{atlasdefinition}

\subsection{복합체와 호몰로지 대수 / Complexes and homological algebra}

\begin{atlasdefinition}{\texttt{17-01}: 가법함자 / Additive functor}
\label{def:17-01}
가법적 범주 사이의 함자 T가 가법이라는 것은 각 대상쌍 A,B와 모든 \(f,g\in\operatorname{Hom}(A,B)\)에 대해 \(T(f+g)=T(f)+T(g)\)를 만족한다는 뜻이다.

\textbf{조건 확인 / Conditions.} 같은 Hom 안의 사상 덧셈을 보존한다.

\textbf{직접 선행 / Prerequisites.} \hyperref[def:13-03]{\texttt{13-03} 공변함자}; \hyperref[def:14-07]{\texttt{14-07} 아벨범주}

\exampleheading
The global-section functor on abelian sheaves over $S^2$ is additive. For sheaf maps $\phi,\psi:\mathcal F\to\mathcal G$ and a global section s, $\Gamma(\phi+\psi)(s)=\phi_X(s)+\psi_X(s)$. It also sends $\mathcal F\oplus\mathcal G$ to $\mathcal F(S^2)\oplus\mathcal G(S^2)$.
\end{atlasdefinition}

\begin{atlasdefinition}{\texttt{03-01}: 한 자리에서의 완전성 / Exactness at a term}
\label{def:03-01}
아벨군 (abelian group) 준동형 (homomorphism)의 수열 \(A\xrightarrow f B\xrightarrow g C\)가 \(B\)에서 완전하다는 것은 \(\operatorname{im}f=\ker g\)라는 뜻이다.

\textbf{조건 확인 / Conditions.} 들어오는 사상의 상과 나가는 사상의 핵이 같다.

\textbf{직접 선행 / Prerequisites.} \hyperref[def:02-03]{\texttt{02-03} 핵 - 아벨군}; \hyperref[def:02-04]{\texttt{02-04} 상 - 아벨군}

\exampleheading
On $S^2$, use the constant integer sections. The sequence $\Z\xrightarrow{\times2}\Z\xrightarrow{\bmod2}\Z/2\Z$ is exact at the middle group: its incoming image and outgoing kernel are both $2\Z$. The maps are induced by maps of the corresponding constant sheaves on $S^2$.
\end{atlasdefinition}

\begin{atlasdefinition}{\texttt{03-02}: 짧은 완전열 / Short exact sequence}
\label{def:03-02}
A,B,C를 아벨군 (abelian group), f와 g를 아래에 표시한 정의역과 공역을 갖는 준동형 (homomorphism)이라 하자. 0은 영군을 뜻한다. \(0\to A\xrightarrow f B\xrightarrow g C\to0\)가 모든 자리에서 완전한 수열이다. 즉 \(f\)는 단사, \(g\)는 전사이며 \(\operatorname{im}f=\ker g\)다.

\textbf{조건 확인 / Conditions.} B에서의 조건만으로는 부족하다.

\textbf{직접 선행 / Prerequisites.} \hyperref[def:03-01]{\texttt{03-01} 한 자리에서의 완전성}

\exampleheading
For $p=0\in\CP^1$ and $i:\{p\}\hookrightarrow X$, the sequence $0\to\mathcal I_p\to\mathcal O_X\xrightarrow{\operatorname{ev}_p}i_*\C\to0$ is short exact. At p it is $0\to(z)\to\C\{z\}\to\C\to0$; away from p it is $0\to\mathcal O_{X,q}\xrightarrow{\id}\mathcal O_{X,q}\to0\to0$.
\end{atlasdefinition}

\begin{atlasdefinition}{\texttt{17-02}: 왼쪽 완전함자 / Left exact functor}
\label{def:17-02}
아벨범주 (abelian category) 사이의 공변 가법함자 (additive functor) T가 모든 짧은 완전열 (short exact sequence) \(0\to A\to B\to C\to0\)를 완전한 수열 \(0\to T(A)\to T(B)\to T(C)\)로 보내면 왼쪽 완전이라 한다.

\textbf{조건 확인 / Conditions.} T(C) 뒤의 0, 즉 마지막 전사는 요구하지 않는다.

\textbf{직접 선행 / Prerequisites.} \hyperref[def:17-01]{\texttt{17-01} 가법함자}; \hyperref[def:03-02]{\texttt{03-02} 짧은 완전열}

\exampleheading
The exponential sequence on $\CP^1$, $0\to\underline\Z\to\mathcal O_X\xrightarrow{f\mapsto e^{2\pi if}}\mathcal O_X^\times\to1$, is exact as a sequence of abelian sheaves. Global sections give $0\to\Z\to\C\to\C^\times$, which is exact. This illustrates preservation of the left end. Local logarithms establish sheaf surjectivity; left exactness alone would not imply surjectivity on global sections.
\end{atlasdefinition}

\begin{atlasdefinition}{\texttt{04-01}: 코사슬복합체 / Cochain complex}
\label{def:04-01}
아벨군 (abelian group)의 코사슬복합체 (cochain complex) \((C^\bullet,d)\)는 각 정수 \(q\in\mathbb Z\)에 대한 아벨군 \(C^q\)와 준동형 (homomorphism) \(d^q:C^q\to C^{q+1}\)의 족으로, 모든 \(q\in\mathbb Z\)에 대해 \[d^{q+1}\circ d^q=0:C^q\longrightarrow C^{q+2}\]을 만족한다. 우변의 0은 모든 원소를 영원소로 보내는 영준동형 (zero homomorphism)이다. 추가로 모든 \(q<0\)에 대해 \(C^q=0\)이면 비음수 차수에 집중된 코사슬복합체라 한다. 이때 0은 영군이며 \(d^{-1}:0\to C^0\)도 영준동형이다.

\textbf{조건 확인 / Conditions.} 음의 차수 항이 0이라는 조건은 일반 정의의 필수조건이 아니라 추가 조건이다.

\textbf{직접 선행 / Prerequisites.} \hyperref[def:02-02]{\texttt{02-02} 아벨군 준동형}

\exampleheading
The de Rham cochain complex of $S^2$ is $C^0=C^\infty(S^2,\R)$, $C^1=\Omega^1(S^2)$, $C^2=\Omega^2(S^2)$, and $C^q=0$ otherwise, with differential d. The condition $d^{q+1}d^q=0$ is the equality of mixed partials in local coordinates.
\end{atlasdefinition}

\begin{atlasdefinition}{\texttt{16-01}: 분해 / Resolution}
\label{def:16-01}
아벨범주 (abelian category)에서 F의 resolution은 \[0\to F\to L^0\xrightarrow{d^0}L^1\xrightarrow{d^1}L^2\to\cdots\]라는 완전열 (exact sequence)이다.

\textbf{조건 확인 / Conditions.} F를 포함한 모든 자리에서 완전해야 한다.

\textbf{직접 선행 / Prerequisites.} \hyperref[def:03-01]{\texttt{03-01} 한 자리에서의 완전성}; \hyperref[def:04-01]{\texttt{04-01} 코사슬복합체}

\exampleheading
On $S^2$, the augmented de Rham sequence $0\to\underline\R\to\mathcal A^0\xrightarrow d\mathcal A^1\xrightarrow d\mathcal A^2\to0$ is a resolution. Exactness is a stalkwise statement: near every point a small coordinate disc admits primitives for closed forms (Poincar\'e lemma). Global area forms need not have global primitives.
\end{atlasdefinition}

\begin{atlasdefinition}{\texttt{15-02}: Injective 대상 / Injective object}
\label{def:15-02}
범주 \(\mathcal C\)와 그 대상 I를 고정한다. 아래의 대상과 사상은 모두 \(\mathcal C\) 안에서 취한다. 대상 I가 injective라는 것은 모든 단사상 (monomorphism) \(j:A\to B\)와 모든 사상 \(u:A\to I\)에 대해 사상 \(v:B\to I\)가 존재하여 \(v\circ j=u\)가 되는 것이다.

\textbf{조건 확인 / Conditions.} 연장의 존재를 요구하며 유일성은 요구하지 않는다.

\textbf{직접 선행 / Prerequisites.} \hyperref[def:15-01]{\texttt{15-01} 단사상 - 범주론}

\exampleheading
At a fixed $p\in S^2$, the skyscraper $i_*\Q$ is injective in abelian sheaves. Indeed $\operatorname{Hom}(\mathcal F,i_*\Q)\simeq\operatorname{Hom}(\mathcal F_p,\Q)$, the stalk functor is exact, and the divisible group $\Q$ is injective. Concretely the map $n\Z\to\Q$, $n\mapsto a$, extends to $\Z\to\Q$, $1\mapsto a/n$. This example is in abelian sheaves, not in $\mathcal O_X$-modules.
\end{atlasdefinition}

\begin{atlasdefinition}{\texttt{16-02}: 단사적 분해 / Injective resolution}
\label{def:16-02}
아벨범주 (abelian category) \(\mathcal C\)와 그 대상 F를 고정한다. 아래의 모든 항과 사상은 \(\mathcal C\) 안에서 취하며 0은 영대상 (zero object)이다. \(0\to F\to I^0\to I^1\to\cdots\)가 완전열 (exact sequence)이고 모든 \(q\ge0\)에 대해 \(I^q\)가 injective인 resolution이다.

\textbf{조건 확인 / Conditions.} 맨 앞의 F 자체는 injective일 필요가 없다.

\textbf{직접 선행 / Prerequisites.} \hyperref[def:16-01]{\texttt{16-01} 분해}; \hyperref[def:15-02]{\texttt{15-02} Injective 대상}

\exampleheading
For $i:\{p\}\hookrightarrow S^2$, $0\to i_*\Z\to i_*\Q\to i_*(\Q/\Z)\to0$ is an injective resolution of $i_*\Z$, concentrated in resolution degrees 0 and 1. Both $\Q$ and $\Q/\Z$ are divisible, hence injective abelian groups. The differential is the quotient map, and its kernel is precisely $\Z$.
\end{atlasdefinition}

\begin{atlasdefinition}{\texttt{04-02}: 코사이클 / Cocycle}
\label{def:04-02}
아벨군 (abelian group)의 코사슬복합체 (cochain complex) \((C^\bullet,d)\)와 정수 q에 대해, \(Z^q(C^\bullet):=\ker(d^q:C^q\to C^{q+1})\)의 원소를 q-cocycle이라 한다. 즉 \(c\in C^q\)가 \(d^q(c)=0\)을 만족하는 것이다.

\textbf{조건 확인 / Conditions.} 차수 q의 원소이며 q차 미분의 핵에 속한다.

\textbf{직접 선행 / Prerequisites.} \hyperref[def:04-01]{\texttt{04-01} 코사슬복합체}; \hyperref[def:02-03]{\texttt{02-03} 핵 - 아벨군}

\exampleheading
In the de Rham complex of $S^2$, the area form $\Omega$ is a degree-two cocycle: it is in $C^2$ and d sends it to zero in $C^3=0$. Being a cocycle does not say it is a coboundary.
\end{atlasdefinition}

\begin{atlasdefinition}{\texttt{04-03}: 코경계 / Coboundary}
\label{def:04-03}
아벨군 (abelian group)의 코사슬복합체 (cochain complex) \((C^\bullet,d)\)와 정수 q에 대해, \(B^q(C^\bullet):=\operatorname{im}(d^{q-1}:C^{q-1}\to C^q)\)의 원소를 q-coboundary라 한다. 즉 어떤 \(b\in C^{q-1}\)에 대해 \(c=d^{q-1}(b)\)로 표현되는 \(c\in C^q\)다.

\textbf{조건 확인 / Conditions.} 비음수 차수에 집중된 복합체에서는 C\ensuremath{{}^{-}}\ensuremath{{}^{1}}=0이므로 B\ensuremath{{}^{0}}=0이다. 일반 복합체에서는 B\ensuremath{{}^{0}}=0일 필요가 없다.

\textbf{직접 선행 / Prerequisites.} \hyperref[def:04-01]{\texttt{04-01} 코사슬복합체}; \hyperref[def:02-04]{\texttt{02-04} 상 - 아벨군}

\exampleheading
In the de Rham complex of $S^2$, $dH\in C^1$ is a degree-one coboundary because H is a global element of $C^0$. It is also a cocycle, since $d(dH)=0$.
\end{atlasdefinition}

\begin{atlasdefinition}{\texttt{04-04}: 복합체의 코호몰로지 / Cohomology of a complex}
\label{def:04-04}
아벨군 (abelian group)의 코사슬복합체 (cochain complex) \((C^\bullet,d)\)와 정수 q에 대해 \[H^q(C^\bullet):=\frac{\ker(d^q:C^q\to C^{q+1})}{\operatorname{im}(d^{q-1}:C^{q-1}\to C^q)}=Z^q/B^q.\]\(d^q\circ d^{q-1}=0\)이므로 분모는 분자의 부분군 (subgroup)이다. 두 cocycle의 차가 coboundary이면 같은 코호몰로지 (cohomology) 원소를 나타낸다.

\textbf{조건 확인 / Conditions.} 분모는 q-1차 미분의 상이다. C\ensuremath{{}^{-}}\ensuremath{{}^{1}}=0인 경우 H\ensuremath{{}^{0}}=ker d\ensuremath{{}^{0}}이다.

\textbf{직접 선행 / Prerequisites.} \hyperref[def:04-02]{\texttt{04-02} 코사이클}; \hyperref[def:04-03]{\texttt{04-03} 코경계}; \hyperref[def:02-05]{\texttt{02-05} 몫군}

\exampleheading
For the de Rham complex of $S^2$, $H^2(C^\bullet)=\Omega^2(S^2)/d\Omega^1(S^2)\simeq\R$. The class of $\Omega/(4\pi)$ is normalized to have integral one; the class of a global form $d\alpha$ is zero.
\end{atlasdefinition}

\begin{atlasdefinition}{\texttt{17-03}: 오른쪽 유도함자 / Right derived functor}
\label{def:17-03}
\(\mathcal A,\mathcal B\)를 아벨범주 (abelian category), \(T:\mathcal A\to\mathcal B\)를 왼쪽 완전 공변 가법함자 (additive functor)라 하고, \(\mathcal A\)가 enough injectives를 갖는다고 가정한다. 대상 F\ensuremath{\in}A의 injective resolution \(0\to F\xrightarrow{\epsilon_F}I^0\xrightarrow{d^0}I^1\to\cdots\)를 택한다. 각 정수 q\ensuremath{\ge}0에 대해 \[R^qT(F):=H^q\bigl(T(I^\bullet),T(d^\bullet)\bigr)\]로 정의한다. 대상 G의 분해 \(0\to G\xrightarrow{\epsilon_G}J^0\to\cdots\)와 사상 \(u:F\to G\)에 대해서는, \(\widetilde u^0\epsilon_F=\epsilon_Gu\)인 복합체 사상 (morphism of complexes) \(\widetilde u:I^\bullet\to J^\bullet\)를 택하여 \(R^qT(u):=H^q(T(\widetilde u))\)로 정한다. 분해와 연장사상의 선택에 대한 독립성은 이 정의를 정당화하는 정리다.

\textbf{조건 확인 / Conditions.} T(I\ensuremath{\bullet})는 T(I\ensuremath{{}^{0}})에서 시작하며 T(F)를 항에 넣지 않는다.

\textbf{직접 선행 / Prerequisites.} \hyperref[def:17-02]{\texttt{17-02} 왼쪽 완전함자}; \hyperref[def:16-02]{\texttt{16-02} 단사적 분해}; \hyperref[def:04-04]{\texttt{04-04} 복합체의 코호몰로지}

\exampleheading
For $F=\Gamma(S^2,-)$, the right derived functors satisfy $R^0F(\underline\R)=\R$, $R^1F(\underline\R)=0$, and $R^2F(\underline\R)=\R$. The last group is computed by the acyclic de Rham resolution and represented by $\Omega/(4\pi)$. Using an acyclic resolution is a theorem about derived functors, not their definition.
\end{atlasdefinition}

\begin{atlasdefinition}{\texttt{04-05}: Acyclic 복합체 / Acyclic complex}
\label{def:04-05}
아벨군 (abelian group)의 코사슬복합체 (cochain complex) \((C^\bullet,d)\)를 고정한다. 차수 q는 정수 전체를 돈다. 복합체 \(C^\bullet\)가 acyclic이라는 것은 모든 차수 \(q\)에서 \(H^q(C^\bullet)=0\)이라는 뜻이다.

\textbf{조건 확인 / Conditions.} 0차도 포함한다. Acyclic 층의 정의와 구별한다.

\textbf{직접 선행 / Prerequisites.} \hyperref[def:04-04]{\texttt{04-04} 복합체의 코호몰로지}

\exampleheading
Let $A=\Gamma(S^2,\underline\R)=\R$. The two-term complex $C^0=A\xrightarrow{\id}C^1=A$, zero in other degrees, is acyclic: $H^0=\ker\id=0$ and $H^1=A/\operatorname{im}\id=0$. Its individual terms are nevertheless nonzero.
\end{atlasdefinition}

\begin{atlasdefinition}{\texttt{05-01}: 복합체 사상 / Morphism of cochain complexes}
\label{def:05-01}
\((C^\bullet,d_C)\), \((D^\bullet,d_D)\)를 아벨군 (abelian group)의 코사슬복합체 (cochain complex)라 하자. \(d_C^q:C^q\to C^{q+1}\), \(d_D^q:D^q\to D^{q+1}\)는 각 복합체의 미분이고 q는 정수다. 복합체 사상 (morphism of complexes) \(f:C^\bullet\to D^\bullet\)는 각 \(q\)의 준동형 (homomorphism) \(f^q:C^q\to D^q\)로 이루어지며 \(d_D^q\circ f^q=f^{q+1}\circ d_C^q\)를 모든 차수에서 만족한다.

\textbf{조건 확인 / Conditions.} 양변 모두 C의 q차에서 D의 q+1차로 간다.

\textbf{직접 선행 / Prerequisites.} \hyperref[def:04-01]{\texttt{04-01} 코사슬복합체}

\exampleheading
Complex conjugation on $\CP^1$ is the smooth map $c([S:T])=[\bar S:\bar T]$, corresponding on the sphere to $(X,Y,Z)\mapsto(X,-Y,Z)$. Pullback $c^*$ is a map of de Rham complexes, because $d(c^*\alpha)=c^*(d\alpha)$. It sends the area form to $-\Omega$.
\end{atlasdefinition}

\begin{atlasdefinition}{\texttt{15-03}: 충분한 단사 대상 / Enough injectives}
\label{def:15-03}
범주 \(\mathcal C\)를 고정한다. 아래의 대상과 사상은 모두 \(\mathcal C\) 안에서 취한다. 모든 대상 F에 대해 어떤 injective 대상 I와 단사상 (monomorphism) \(F\hookrightarrow I\)가 존재하면 그 범주가 enough injectives를 갖는다고 한다.

\textbf{조건 확인 / Conditions.} injective 대상 하나의 존재만을 뜻하지 않는다.

\textbf{직접 선행 / Prerequisites.} \hyperref[def:15-02]{\texttt{15-02} Injective 대상}

\exampleheading
For any abelian sheaf $\mathcal F$ on $S^2$, choose group embeddings $\mathcal F_p\hookrightarrow I_p$ into injective abelian groups. Sending a section to all its germs defines a monomorphism $\mathcal F\hookrightarrow\prod_{p\in S^2}i_{p*}I_p$. Each factor is injective by stalk adjunction, and a product of injective objects is injective here. This is a concrete construction on the fixed sphere, using the standard enough-injectives theorem for abelian groups.
\end{atlasdefinition}

\begin{atlasdefinition}{\texttt{23-01}: 사슬복합체 / Chain complex}
\label{def:23-01}
아벨군 (abelian group)의 사슬복합체 (chain complex) \((C_\bullet,\partial)\)는 각 정수 \(q\in\mathbb Z\)에 대한 아벨군 \(C_q\)와 준동형 (homomorphism) \(\partial_q:C_q\to C_{q-1}\)의 족으로, 모든 정수 q에 대해 \[\partial_{q-1}\circ\partial_q=0:C_q\longrightarrow C_{q-2}\]을 만족한다. 우변의 0은 영준동형 (zero homomorphism)이다. 추가로 모든 q<0에서 \(C_q=0\)이면 비음수 차수에 집중되었다고 하며, 이때 \(\partial_0:C_0\to0\)은 영준동형이다.

\textbf{조건 확인 / Conditions.} 경계는 차수를 하나 내린다. 음의 차수 항의 소멸은 추가 조건이다.

\textbf{직접 선행 / Prerequisites.} \hyperref[def:02-02]{\texttt{02-02} 아벨군 준동형}

\exampleheading
Triangulate $S^2$ by the boundary K of a tetrahedron. Its simplicial chain complex has $C_2=\Z^4$, $C_1=\Z^6$, $C_0=\Z^4$, with $\partial[i,j,k]=[j,k]-[i,k]+[i,j]$. All groups in degrees above two vanish, and $\partial^2=0$.
\end{atlasdefinition}

\begin{atlasdefinition}{\texttt{23-02}: 사이클 / Cycle}
\label{def:23-02}
\((C_\bullet,\partial)\)를 아벨군 (abelian group)의 사슬복합체 (chain complex), q를 정수라 하자. 사슬복합체에서 \(\ker(\partial_q:C_q\to C_{q-1})\)의 원소를 q-cycle이라 한다.

\textbf{조건 확인 / Conditions.} q차 경계가 0인 원소다.

\textbf{직접 선행 / Prerequisites.} \hyperref[def:23-01]{\texttt{23-01} 사슬복합체}; \hyperref[def:02-03]{\texttt{02-03} 핵 - 아벨군}

\exampleheading
In the tetrahedral model of $S^2$, $c=[1,2,3]-[0,2,3]+[0,1,3]-[0,1,2]$ is a two-cycle. Applying the displayed boundary formula cancels every oriented edge exactly twice with opposite signs.
\end{atlasdefinition}

\begin{atlasdefinition}{\texttt{23-03}: 경계 / Boundary}
\label{def:23-03}
\((C_\bullet,\partial)\)를 아벨군 (abelian group)의 사슬복합체 (chain complex), q를 정수라 하자. 사슬복합체에서 \(\operatorname{im}(\partial_{q+1}:C_{q+1}\to C_q)\)의 원소를 q-boundary라 한다.

\textbf{조건 확인 / Conditions.} q+1차에서 내려온 원소다.

\textbf{직접 선행 / Prerequisites.} \hyperref[def:23-01]{\texttt{23-01} 사슬복합체}; \hyperref[def:02-04]{\texttt{02-04} 상 - 아벨군}

\exampleheading
On the tetrahedrally triangulated $S^2$, the chain $b=[1,2]-[0,2]+[0,1]$ is the boundary $\partial[0,1,2]$. It is a one-cycle, since the endpoint contributions in $\partial b$ cancel.
\end{atlasdefinition}

\begin{atlasdefinition}{\texttt{23-04}: 호몰로지 / Homology}
\label{def:23-04}
\((C_\bullet,\partial)\)를 아벨군 (abelian group)의 사슬복합체 (chain complex), q를 정수라 하자. q차 호몰로지 (homology)는 \[H_q(C_\bullet):=\frac{\ker(\partial_q:C_q\to C_{q-1})}{\operatorname{im}(\partial_{q+1}:C_{q+1}\to C_q)}\]다. \(\partial_q\circ\partial_{q+1}=0\)이므로 분모는 분자의 부분군 (subgroup)이다. 두 cycle의 차가 boundary이면 같은 호몰로지 원소를 나타낸다.

\textbf{조건 확인 / Conditions.} 분모의 경계는 q+1차에서 출발한다.

\textbf{직접 선행 / Prerequisites.} \hyperref[def:23-02]{\texttt{23-02} 사이클}; \hyperref[def:23-03]{\texttt{23-03} 경계}; \hyperref[def:02-05]{\texttt{02-05} 몫군}

\exampleheading
For the tetrahedral model of $S^2$, the class of c from the cycle example generates $H_2(K;\Z)\simeq\Z$. It is not a boundary in this simplicial complex because $C_3(K)=0$. The one-cycle bounding a face represents zero in $H_1(K;\Z)$.
\end{atlasdefinition}

\subsection{코호몰로지의 계산 / Computing cohomology}

\begin{atlasdefinition}{\texttt{18-01}: 전역단면함자 / Global section functor}
\label{def:18-01}
X를 위상공간 (topological space)이라 하자. \(\operatorname{Sh}(X;\mathbf{Ab})\)는 X 위 아벨군 층 (sheaf of abelian groups)과 층사상 (sheaf morphism)의 범주, \(\mathbf{Ab}\)는 아벨군 (abelian group)과 준동형 (homomorphism)의 범주다. 전역단면함자 (global section functor) \(\Gamma(X,-):\operatorname{Sh}(X;\mathbf{Ab})\to\mathbf{Ab}\)는 대상 \(\mathcal F\)를 \(\mathcal F(X)\)로, 층사상 \(f:\mathcal F\to\mathcal G\)를 그 X에서의 성분 \(f_X:\mathcal F(X)\to\mathcal G(X)\)로 보낸다.

\textbf{조건 확인 / Conditions.} 공간 X를 고정한 함자다.

\textbf{직접 선행 / Prerequisites.} \hyperref[def:07-03]{\texttt{07-03} 전역단면}; \hyperref[def:10-01]{\texttt{10-01} 층사상}; \hyperref[def:13-03]{\texttt{13-03} 공변함자}

\exampleheading
For $X=\CP^1$, the functor $\Gamma(X,-)$ gives $\Gamma(X,\mathcal O_X)=\C$ and $\Gamma(X,\mathcal K_X)=0$. It gives $\Gamma(X,i_*\C)=\C$ for a skyscraper at any fixed p. The inputs are different sheaves on the same space.
\end{atlasdefinition}

\begin{atlasdefinition}{\texttt{18-02}: 층 코호몰로지 / Sheaf cohomology}
\label{def:18-02}
X를 위상공간 (topological space), \(\mathcal F\)를 X 위 아벨군 층 (sheaf of abelian groups)이라 하자. 층의 injective resolution \(0\to\mathcal F\to I^0\xrightarrow{d^0}I^1\to\cdots\)를 택한다. \(D^r:=\Gamma(X,d^r):I^r(X)\to I^{r+1}(X)\)로 놓고, \(I^{-1}(X)=0\), \(D^{-1}:0\to I^0(X)\)으로 둔다. 각 정수 q\ensuremath{\ge}0의 층 코호몰로지 (sheaf cohomology)는 \[H^q(X,\mathcal F):=\frac{\ker(D^q:I^q(X)\to I^{q+1}(X))}{\operatorname{im}(D^{q-1}:I^{q-1}(X)\to I^q(X))}=R^q\Gamma(X,\mathcal F)\]다. \(R^q\Gamma\)는 전역단면함자 (global section functor)의 q차 오른쪽 유도함자 (right derived functor)를 뜻한다.

\textbf{조건 확인 / Conditions.} 전역단면 (global section) 복합체는 \ensuremath{\Gamma}(X,I\ensuremath{{}^{0}})에서 시작한다.

\textbf{직접 선행 / Prerequisites.} \hyperref[def:18-01]{\texttt{18-01} 전역단면함자}; \hyperref[def:17-03]{\texttt{17-03} 오른쪽 유도함자}

\exampleheading
For the real constant sheaf on $S^2$, $H^2(S^2,\underline\R)\simeq\R$, computed by the fine de Rham resolution and represented there by $\Omega/(4\pi)$. On the analytic $\CP^1$, $H^1(X,\mathcal O_X)=0$ but $H^1(X,\mathcal K_X)\simeq\C$. The coefficient sheaf is essential.
\end{atlasdefinition}

\begin{atlasdefinition}{\texttt{N22}: 층 코호몰로지의 차원과 Euler characteristic / Cohomology dimension and Euler characteristic}
\label{def:N22}
X를 복소다양체 (complex manifold), E를 \(\mathcal O_X\)-가군층 (sheaf of modules)이라 하자. \(H^q(X,E)\)가 유한차원 복소벡터공간일 때 \(h^q(X,E):=\dim_{\mathbb C}H^q(X,E)\)라 쓴다. 모든 q\ensuremath{\ge}0에서 이 차원이 유한하고 충분히 큰 q에서 0이면 \[\chi(X,E):=\sum_{q\ge0}(-1)^q h^q(X,E)\]로 정의한다.

\textbf{조건 확인 / Conditions.} 위상적 Betti 수와 층 코호몰로지 (sheaf cohomology) 차원 h\ensuremath{{}^{q}}를 구별한다. 이 합은 주어진 E에 의존한다.

\textbf{직접 선행 / Prerequisites.} \hyperref[def:18-02]{\texttt{18-02} 층 코호몰로지}; \hyperref[def:B01]{\texttt{B01} 실수·복소 벡터공간}

\exampleheading
On $\CP^1$, $h^0(\mathcal O(n))=\max(n+1,0)$ and $h^1(\mathcal O(n))=\max(-n-1,0)$, so $\chi(\mathcal O(n))=n+1$. In particular $\chi(K)=-1$ while $\deg K=-2$; for the tangent line, $\chi(T^{1,0}X)=3$ while its degree is 2. These are different invariants.
\end{atlasdefinition}

\begin{atlasdefinition}{\texttt{06-14}: de Rham 코호몰로지 / de Rham cohomology}
\label{def:06-14}
M을 매끄러운 다양체 (smooth manifold), q\ensuremath{\ge}0을 정수라 하자. \(\Omega^r(M)\)는 M 위 매끄러운 r-형식의 공간, \(d^r:\Omega^r(M)\to\Omega^{r+1}(M)\)는 외미분 (exterior derivative)이다. 음의 차수 항은 영공간이고 \(d^{-1}:0\to\Omega^0(M)\)는 영사상 (zero morphism)으로 둔다. q차 de Rham 코호몰로지 (cohomology)는 \[H^q_{\mathrm{dR}}(M):=\frac{\ker(d^q:\Omega^q(M)\to\Omega^{q+1}(M))}{\operatorname{im}(d^{q-1}:\Omega^{q-1}(M)\to\Omega^q(M))}\]다. \(d^q d^{q-1}=0\)이므로 분모는 분자의 부분공간이다.

\textbf{조건 확인 / Conditions.} 같은 M 위의 closed 형식을 exact 형식으로 나눈다.

\textbf{직접 선행 / Prerequisites.} \hyperref[def:06-11]{\texttt{06-11} Closed 형식}; \hyperref[def:06-12]{\texttt{06-12} Exact 형식}; \hyperref[def:04-04]{\texttt{04-04} 복합체의 코호몰로지}

\exampleheading
Integration induces $H^2_{\mathrm{dR}}(S^2;\R)\simeq\R$, $[\alpha]\mapsto\int_{S^2}\alpha$. The class $[\Omega/(4\pi)]$ maps to 1. In degrees 0 and 1 the groups are $\R$ and 0, respectively. These are the computed groups of this sphere, not part of the general definition.
\end{atlasdefinition}

\begin{atlasdefinition}{\texttt{12-04}: 함수의 지지 / Support of a function}
\label{def:12-04}
X를 위상공간 (topological space), \(f:X\to\mathbb R\)를 함수라 하자. f의 지지는 \[\operatorname{supp}f:=\overline{\{x\in X:f(x)\ne0\}}^{\,X}\]다. 폐포 (closure)는 정의역 X의 위상에서 취한다. 복소수값 함수에도 같은 정의를 사용한다.

\textbf{조건 확인 / Conditions.} 0이 아닌 점들의 집합 자체가 아니라 그 폐포 (closure)다.

\textbf{직접 선행 / Prerequisites.} \hyperref[def:01-05]{\texttt{01-05} 폐포}; \hyperref[def:06-04]{\texttt{06-04} 매끄러운 다양체}

\exampleheading
Use the smooth function $\eta(t)=e^{-1/t}$ for $t>0$ and 0 for $t\le0$. On $S^2$, $f(p)=\eta(Z(p))$ has support the closed northern hemisphere $\{Z\ge0\}$: support includes the equator although f vanishes there.
\end{atlasdefinition}

\begin{atlasdefinition}{\texttt{12-05}: 매끄러운 단위분할 / Smooth partition of unity}
\label{def:12-05}
매끄러운 다양체 (smooth manifold) M의 열린덮개 (open cover) \(\{U_i\}\)에 대해, 매끄러운 함수 \(\rho_i:M\to\mathbb R\)들이 \[0\le\rho_i\le1,\quad\operatorname{supp}\rho_i\subset U_i,\quad\sum_i\rho_i=1\]을 만족하고 지지들의 족이 국소유한 (locally finite)이면 이 덮개에 종속된 매끄러운 단위분할 (partition of unity)이다.

\textbf{조건 확인 / Conditions.} 지지들의 국소유한 (locally finite)성이 합을 정의하게 한다.

\textbf{직접 선행 / Prerequisites.} \hyperref[def:12-04]{\texttt{12-04} 함수의 지지}; \hyperref[def:12-01]{\texttt{12-01} 국소유한 집합족}; \hyperref[def:01-03]{\texttt{01-03} 열린덮개}

\exampleheading
On $S^2$, put $\rho_0=\eta(Z+1/2)/(\eta(Z+1/2)+\eta(1/2-Z))$ and $\rho_\infty=1-\rho_0$. The denominator is everywhere positive. These smooth functions sum to 1, with supports in $\{Z\ge-1/2\}\subset U_0$ and $\{Z\le1/2\}\subset U_\infty$. This is an explicit partition subordinate to the two standard charts.
\end{atlasdefinition}

\begin{atlasdefinition}{\texttt{12-06}: 층 자기 사상의 지지 / Support of a sheaf endomorphism}
\label{def:12-06}
X를 위상공간 (topological space), \(\mathcal F\)를 X 위 아벨군 층 (sheaf of abelian groups), \(h:\mathcal F\to\mathcal F\)를 층사상 (sheaf morphism)이라 하자. 각 x\ensuremath{\in}X에서 \(h_x:\mathcal F_x\to\mathcal F_x\)는 \([U,s]_x\mapsto[U,h_U(s)]_x\)로 유도되는 줄기 준동형 (homomorphism on stalks)이다. h의 지지를 \[\operatorname{supp}h:=\overline{\{x\in X:h_x\ne0\}}^{\,X}\]로 정의한다. 여기서 \(h_x\ne0\)은 줄기 준동형이 영준동형 (zero homomorphism)이 아니라는 뜻이며, 폐포 (closure)는 X에서 취한다.

\textbf{조건 확인 / Conditions.} h\ensuremath{{}_{x}}는 줄기사상 (induced map on stalks)이며 그 비영점 집합의 폐포 (closure)를 취한다.

\textbf{직접 선행 / Prerequisites.} \hyperref[def:10-01]{\texttt{10-01} 층사상}; \hyperref[def:09-02]{\texttt{09-02} Stalk - 줄기}; \hyperref[def:01-05]{\texttt{01-05} 폐포}

\exampleheading
On the smooth function sheaf $\mathcal A^0$ of $S^2$, multiplication $m_{\rho_0}$ is a sheaf endomorphism. Its support is $\{Z\ge-1/2\}$, the closure of the points where the germ of $\rho_0$ is nonzero. At a boundary point the value may be zero while the germ is not the zero germ.
\end{atlasdefinition}

\begin{atlasdefinition}{\texttt{12-07}: Fine 층 / Fine sheaf}
\label{def:12-07}
X를 paracompact Hausdorff 공간, \(\mathcal F\)를 X 위 아벨군 층 (sheaf of abelian groups)이라 하자. \(\mathcal F\)가 fine이라는 것은 X의 모든 국소유한 (locally finite) 열린덮개 (open cover) \(\{U_i\}_{i\in I}\)에 대해 층사상족 \((h_i:\mathcal F\to\mathcal F)_{i\in I}\)이 존재하여 \[\operatorname{supp}h_i\subset U_i\quad(i\in I),\qquad\sum_{i\in I}h_i=\mathrm{id}_{\mathcal F}\]를 만족한다는 뜻이다. 각 점 근방에서 유한 개의 h\ensuremath{{}_{i}}만 0이 아닐 수 있으므로 이 합이 정의된다.

\textbf{조건 확인 / Conditions.} 주어진 모든 국소유한 (locally finite) 열린덮개 (open cover)에 대해 그 덮개에 종속된 층 자기 사상족의 존재를 요구한다.

\textbf{직접 선행 / Prerequisites.} \hyperref[def:12-06]{\texttt{12-06} 층 자기 사상의 지지}; \hyperref[def:12-01]{\texttt{12-01} 국소유한 집합족}

\exampleheading
For the two-chart cover of $S^2$, the endomorphisms $m_{\rho_0}$ and $m_{\rho_\infty}$ of $\mathcal A^k$ have supports inside the corresponding charts and sum to the identity. For any locally finite cover, a subordinate smooth partition yields the same construction, proving that smooth differential-form sheaves are fine. One fixed partition alone illustrates, but does not exhaust, the quantifier over covers.
\end{atlasdefinition}

\begin{atlasdefinition}{\texttt{16-03}: 파인 분해 / Fine resolution}
\label{def:16-03}
X를 paracompact Hausdorff 공간, \(\mathcal F\)를 X 위 아벨군 층 (sheaf of abelian groups)이라 하자. 아래의 \(\mathcal L^q\)와 사상은 X 위 아벨군 층의 범주에서 취한다. \(0\to\mathcal F\to\mathcal L^0\to\mathcal L^1\to\cdots\)가 층의 완전열 (exact sequence)이고 모든 \(q\ge0\)의 \(\mathcal L^q\)가 fine인 resolution이다. Fine의 공간 가정은 paracompact Hausdorff다.

\textbf{조건 확인 / Conditions.} 해결되는 F 자체가 fine일 필요는 없다.

\textbf{직접 선행 / Prerequisites.} \hyperref[def:16-01]{\texttt{16-01} 분해}; \hyperref[def:12-07]{\texttt{12-07} Fine 층}

\exampleheading
The sequence $0\to\underline\R\to\mathcal A^0\to\mathcal A^1\to\mathcal A^2\to0$ on $S^2$ is a fine resolution: Poincar\'e lemma gives exactness, and subordinate partitions act by multiplication to make each $\mathcal A^k$ fine.
\end{atlasdefinition}

\begin{atlasdefinition}{\texttt{19-01}: Acyclic 층 / Acyclic sheaf}
\label{def:19-01}
X를 위상공간 (topological space), \(\mathcal F\)를 X 위 아벨군 층 (sheaf of abelian groups)이라 하자. \(H^q(X,\mathcal F)\)는 이 층의 q차 층 코호몰로지 (sheaf cohomology)이며 q는 정수다. 층 \(\mathcal F\)가 X에서 acyclic이라는 것은 모든 \(q>0\)에 대해 \(H^q(X,\mathcal F)=0\)이라는 뜻이다.

\textbf{조건 확인 / Conditions.} 0차 소멸은 요구하지 않고 공간 X를 명시한다.

\textbf{직접 선행 / Prerequisites.} \hyperref[def:18-02]{\texttt{18-02} 층 코호몰로지}

\exampleheading
The smooth function sheaf $\mathcal A^0$ on $S^2$ is acyclic for global sections: $H^q(S^2,\mathcal A^0)=0$ for every q>0, by the fine-sheaf acyclicity theorem on paracompact Hausdorff spaces. Its global sections are not zero; for instance it contains H and the constant 1.
\end{atlasdefinition}

\begin{atlasdefinition}{\texttt{19-02}: 비순환 분해 / Acyclic resolution}
\label{def:19-02}
X를 위상공간 (topological space), \(\mathcal F\)를 X 위 아벨군 층 (sheaf of abelian groups)이라 하자. X에서 acyclic인 층 L이란 모든 정수 q>0에 대해 \(H^q(X,L)=0\)인 층이다. X 위 층의 resolution \(0\to\mathcal F\to\mathcal L^0\to\mathcal L^1\to\cdots\)에서 모든 \(p\ge0\)의 \(\mathcal L^p\)가 X에서 acyclic인 것이다.

\textbf{조건 확인 / Conditions.} 각 p에 대해 모든 q>0의 H\ensuremath{{}^{q}}(X,L\ensuremath{{}^{p}})가 0이다.

\textbf{직접 선행 / Prerequisites.} \hyperref[def:19-01]{\texttt{19-01} Acyclic 층}; \hyperref[def:16-01]{\texttt{16-01} 분해}

\exampleheading
The fine de Rham resolution on $S^2$ is also an acyclic resolution, since its terms have no higher sheaf cohomology. Taking global sections therefore computes $H^q(S^2,\underline\R)$. No assertion that each term is an injective sheaf is needed.
\end{atlasdefinition}

\begin{atlasdefinition}{\texttt{20-01}: Čech 코사슬 / Čech cochain}
\label{def:20-01}
X를 위상공간 (topological space), \(\mathcal F\)를 X 위 아벨군 층 (sheaf of abelian groups), \(\mathfrak U=(U_i)_{i\in I}\)를 X의 열린덮개 (open cover)라 하자. 지표집합 I에 전순서를 고정한다. 정수 p\ensuremath{\ge}0에 대해 \(U_{i_0\cdots i_p}:=\bigcap_{a=0}^pU_{i_a}\)로 놓고 \[\check C^p(\mathfrak U,\mathcal F):=\prod_{i_0<\cdots<i_p}\mathcal F(U_{i_0\cdots i_p})\]로 정의한다. p-코사슬은 모든 증가 지표열에 대해 단면 (section) \(c_{i_0\cdots i_p}\in\mathcal F(U_{i_0\cdots i_p})\)를 하나씩 지정한 족이다. 덧셈은 성분별로 하며, 빈 교집합의 성분군은 \(\mathcal F(\varnothing)=0\)이다.

\textbf{조건 확인 / Conditions.} 직합이 아니라 곱이며 지표는 증가 순서다.

\textbf{직접 선행 / Prerequisites.} \hyperref[def:01-03]{\texttt{01-03} 열린덮개}; \hyperref[def:08-01]{\texttt{08-01} 층}

\exampleheading
For $\mathfrak U=(U_0,U_\infty)$ on $\CP^1$ and $\mathcal K_X$, a zero-cochain is $(a(z)dz,b(w)dw)$ with a,b entire. A one-cochain is $f(z)dz$ with f holomorphic on $\C^\times$. Thus the overlap has its own sections; it is not just the intersection of two vector spaces of global sections.
\end{atlasdefinition}

\begin{atlasdefinition}{\texttt{20-02}: Čech 미분 / Čech differential / Coboundary operator}
\label{def:20-02}
X를 위상공간 (topological space), \(\mathcal F\)를 X 위 아벨군 층 (sheaf of abelian groups), \(\mathfrak U=(U_i)_{i\in I}\)를 지표 I에 전순서가 주어진 열린덮개 (open cover)라 하자. \(U_{i_0\cdots i_r}:=\bigcap_{b=0}^rU_{i_b}\)로 쓴다. p\ensuremath{\ge}0에 대해 \(\check C^p(\mathfrak U,\mathcal F)=\prod_{i_0<\cdots<i_p}\mathcal F(U_{i_0\cdots i_p})\)의 원소를 \(c=(c_{i_0\cdots i_p})\)라 하자. Čech 미분 \(\delta^p:\check C^p(\mathfrak U,\mathcal F)\to\check C^{p+1}(\mathfrak U,\mathcal F)\)는 모든 \(i_0<\cdots<i_{p+1}\)에서 \[(\delta^pc)_{i_0\cdots i_{p+1}}:=\sum_{a=0}^{p+1}(-1)^a c_{i_0\cdots\widehat{i_a}\cdots i_{p+1}}|_{U_{i_0\cdots i_{p+1}}}\]로 정의한다. 모자는 그 지표의 생략, 세로줄은 층의 제한사상 (restriction map)을 뜻한다. 음의 차수 코사슬군은 0이고 \(\delta^{-1}:0\to\check C^0\)는 영준동형 (zero homomorphism)이다.

\textbf{조건 확인 / Conditions.} 모든 항을 동일한 전체 교집합으로 제한한다.

\textbf{직접 선행 / Prerequisites.} \hyperref[def:20-01]{\texttt{20-01} Čech 코사슬}

\exampleheading
On the two standard charts of $\CP^1$, using the standard alternating convention $(\delta s)_{0\infty}=s_\infty-s_0$, the preceding zero-cochain maps, in the dz frame, to $(-z^{-2}b(1/z)-a(z))dz$. With two members there are no increasing triple indices, so the next alternating differential is zero.
\end{atlasdefinition}

\begin{atlasdefinition}{\texttt{21-01}: 한 덮개의 Čech 코호몰로지 / Čech cohomology of a cover}
\label{def:21-01}
X를 위상공간 (topological space), \(\mathcal F\)를 X 위 아벨군 층 (sheaf of abelian groups), \(\mathfrak U\)를 X의 열린덮개 (open cover)라 하자. 이 덮개의 Čech 코사슬복합체 (cochain complex)를 \((\check C^\bullet(\mathfrak U,\mathcal F),\delta)\)라 쓰며, \(\delta^r:\check C^r\to\check C^{r+1}\)는 교집합으로의 제한들을 교대부호로 합한 Čech 미분이다. 정수 p\ensuremath{\ge}0에 대해 \[\check H^p(\mathfrak U,\mathcal F):=\frac{\ker(\delta^p:\check C^p\to\check C^{p+1})}{\operatorname{im}(\delta^{p-1}:\check C^{p-1}\to\check C^p)}\]로 정의한다. 모든 코사슬군은 같은 \((\mathfrak U,\mathcal F)\)에 대한 것이며 \(\check C^{-1}=0\), \(\delta^{-1}=0\)이다.

\textbf{조건 확인 / Conditions.} 고정한 한 덮개에 대한 군이다.

\textbf{직접 선행 / Prerequisites.} \hyperref[def:20-02]{\texttt{20-02} Čech 미분}; \hyperref[def:04-04]{\texttt{04-04} 복합체의 코호몰로지}

\exampleheading
For the analytic two-chart cover of $\CP^1$, $\check H^1(\mathfrak U,\mathcal K_X)$ is $\mathcal O(\C^\times)/(\mathcal O(\C)+z^{-2}\mathcal O(\C)\vert_{w=1/z})$, times dz. Laurent expansion leaves precisely the coefficient of $z^{-1}$, so the group is $\C\langle[dz/z]\rangle$. Positive and sufficiently negative Laurent tails extend to the respective charts.
\end{atlasdefinition}

\begin{atlasdefinition}{\texttt{21-02}: 전체 Čech 코호몰로지 / Čech cohomology of a space}
\label{def:21-02}
X를 위상공간 (topological space), \(\mathcal F\)를 X 위 아벨군 층 (sheaf of abelian groups), p\ensuremath{\ge}0을 정수라 하자. 열린덮개 (open cover) \(\mathfrak V\)가 \(\mathfrak U\)를 세분하면 층의 제한사상 (restriction map)들이 유도하는 사상을 \(r_{\mathfrak U\mathfrak V}:\check H^p(\mathfrak U,\mathcal F)\to\check H^p(\mathfrak V,\mathcal F)\)라 쓴다. 모든 열린덮개와 이 사상들에 대한 직극한 \[\check H^p(X,\mathcal F):=\varinjlim_{\mathfrak U}\check H^p(\mathfrak U,\mathcal F)\]를 전체 Čech 코호몰로지 (cohomology)라 한다. 구체적으로 쌍 \((\mathfrak U,\alpha)\), \(\alpha\in\check H^p(\mathfrak U,\mathcal F)\)들의 동치류 (equivalence class)이며, \((\mathfrak U,\alpha)\sim(\mathfrak V,\beta)\)는 어떤 공통 세분 \(\mathfrak W\)에서 \(r_{\mathfrak U\mathfrak W}(\alpha)=r_{\mathfrak V\mathfrak W}(\beta)\)라는 뜻이다. 덧셈은 공통 세분으로 옮긴 두 원소를 더하여 정의한다.

\textbf{조건 확인 / Conditions.} 세분하면 원래 덮개의 코호몰로지 (cohomology)에서 세분된 덮개의 것으로 간다.

\textbf{직접 선행 / Prerequisites.} \hyperref[def:21-01]{\texttt{21-01} 한 덮개의 Čech 코호몰로지}; \hyperref[def:12-02]{\texttt{12-02} 열린덮개의 세분}

\exampleheading
Take the open-star good cover of the tetrahedral $S^2$ from the common conventions. Its degree-two constant-real Čech cocycle assigning 1 to the face (1,2,3) and 0 to the other faces gives the nonzero orientation class, after passing through refinement maps to full Čech cohomology. Its value on the oriented fundamental cycle is 1, so it cannot be a coboundary. Good-cover comparison identifies this with the generator of $H^2(S^2;\R)$.
\end{atlasdefinition}

\begin{atlasdefinition}{\texttt{22-01}: F-acyclic 덮개 / F-acyclic cover}
\label{def:22-01}
X를 위상공간 (topological space), \(\mathcal F\)를 X 위 아벨군 층 (sheaf of abelian groups), \(\mathfrak U=(U_i)_{i\in I}\)를 X의 열린덮개 (open cover)라 하자. 모든 정수 p\ensuremath{\ge}0, 모든 유한 지표열 \(i_0,\ldots,i_p\)에 대해 \(W:=U_{i_0}\cap\cdots\cap U_{i_p}\ne\varnothing\)이면 모든 정수 q>0에서 \[H^q(W,\mathcal F|_W)=0\]일 때 \(\mathfrak U\)를 \(\mathcal F\)-acyclic 덮개라 한다. \(\mathcal F|_W\)는 W로의 제한층이며 H\ensuremath{{}^{q}}는 층 코호몰로지 (sheaf cohomology)다. p=0은 W=U\ensuremath{{}_{i}}인 경우다.

\textbf{조건 확인 / Conditions.} U\ensuremath{{}_{i}} 하나인 교집합도 포함하며 계수층 F를 명시한다.

\textbf{직접 선행 / Prerequisites.} \hyperref[def:19-01]{\texttt{19-01} Acyclic 층}; \hyperref[def:01-03]{\texttt{01-03} 열린덮개}

\exampleheading
On analytic $\CP^1$, $U_0\simeq\C$, $U_\infty\simeq\C$, and their overlap $\C^\times$ are Stein. The restrictions of the locally free sheaf $\mathcal K_X$ have vanishing higher cohomology by Cartan's theorem B. Consequently this cover is $\mathcal K_X$-acyclic and its Čech calculation computes sheaf cohomology. The Stein acyclicity theorem is an explicitly used theorem, not a property of every open cover.
\end{atlasdefinition}

\begin{atlasdefinition}{\texttt{24-01}: 표준 단체 / Standard simplex}
\label{def:24-01}
q를 0 이상의 정수라 하자. \[\Delta^q=\{(t_0,\ldots,t_q)\in\mathbb R^{q+1}:t_i\ge0\text{ for all }i,\ \sum_{i=0}^q t_i=1\}.\]

\textbf{조건 확인 / Conditions.} 좌표는 q+1개이며 부등호는 \ensuremath{\ge}다.

\textbf{직접 선행 / Prerequisites.} \hyperref[def:B01]{\texttt{B01} 실수·복소 벡터공간}

\exampleheading
For the triangle used to parametrize a face of the tetrahedral $S^2$, $\Delta^2=\{(t_0,t_1,t_2)\in\R^3:t_i\ge0,\ t_0+t_1+t_2=1\}$. Its barycentre is (1/3,1/3,1/3), and its three vertices are the coordinate vectors. This is the parameter simplex, not the spherical image itself.
\end{atlasdefinition}

\begin{atlasdefinition}{\texttt{24-02}: 특이 단체 / Singular simplex}
\label{def:24-02}
q를 0 이상의 정수라 하자. \(\Delta^q\)는 \(t_i\ge0\), \(\sum_{i=0}^q t_i=1\)을 만족하는 \((t_0,\ldots,t_q)\in\mathbb R^{q+1}\)의 공간이며 부분공간 위상을 준다. 위상공간 (topological space) X의 특이 q-단체는 연속사상 \(\sigma:\Delta^q\to X\)다.

\textbf{조건 확인 / Conditions.} X에서 단체로 가는 사상이 아니며 단사일 필요도 없다.

\textbf{직접 선행 / Prerequisites.} \hyperref[def:24-01]{\texttt{24-01} 표준 단체}; \hyperref[def:06-01]{\texttt{06-01} 연속사상}

\exampleheading
Let $v_0,v_1,v_2$ be three vertices of the inscribed tetrahedron. The map $\sigma(t)=\bigl(\sum_{j=0}^2t_jv_j\bigr)/\|\sum_{j=0}^2t_jv_j\|$ from $\Delta^2$ to $S^2$ is a singular two-simplex. Its denominator is nonzero because a tetrahedral face lies in a plane not through the origin.
\end{atlasdefinition}

\begin{atlasdefinition}{\texttt{24-03}: 특이 사슬군 / Singular chain group}
\label{def:24-03}
X를 위상공간 (topological space), q를 0 이상의 정수라 하자. \(\Delta^q\)는 표준 q-단체다. 아래의 \(a_k\)는 실수이고 N은 0 이상의 정수다. \(C_q(X;\mathbb R)\)는 X의 모든 특이 q-단체 \(\sigma:\Delta^q\to X\)를 기저로 하는 실수벡터공간이다. 원소는 \(\sum_{k=1}^N a_k\sigma_k\)인 유한 선형결합이다.

\textbf{조건 확인 / Conditions.} 개별 사슬은 유한합이다.

\textbf{직접 선행 / Prerequisites.} \hyperref[def:24-02]{\texttt{24-02} 특이 단체}; \hyperref[def:02-01]{\texttt{02-01} 아벨군}

\exampleheading
The singular chain group $C_2(S^2;\Z)$ consists of finite formal integer combinations of continuous maps $\Delta^2\to S^2$. If $\sigma_{012}$ and $\sigma_{013}$ are radial face maps, $2\sigma_{012}-\sigma_{013}$ is an element. The coefficients combine maps formally; they do not add points of $S^2$.
\end{atlasdefinition}

\begin{atlasdefinition}{\texttt{24-04}: 특이 경계 / Singular boundary operator}
\label{def:24-04}
X를 위상공간 (topological space)이라 하자. \(C_q(X;\mathbb R)\)는 연속사상 \(\sigma:\Delta^q\to X\)들을 기저로 하는 실수벡터공간이고 \(\partial_q:C_q(X;\mathbb R)\to C_{q-1}(X;\mathbb R)\)를 정의한다. \(C_{-1}(X;\mathbb R)=0\)으로 둔다. q>0에서 \(\epsilon_i:\Delta^{q-1}\to\Delta^q\)를 i번째 좌표에 0을 넣는 순서 보존 면 포함이라 할 때 \[\partial_q\sigma=\sum_{i=0}^q(-1)^i\sigma\circ\epsilon_i\]로 정의하고 선형으로 확장한다. \(\partial_0=0\)이다.

\textbf{조건 확인 / Conditions.} i번째 꼭짓점의 반대 면에 제한한다.

\textbf{직접 선행 / Prerequisites.} \hyperref[def:24-03]{\texttt{24-03} 특이 사슬군}

\exampleheading
For the radial face map $\sigma_{012}:\Delta^2\to S^2$, its singular boundary is $\partial\sigma_{012}=\sigma_{12}-\sigma_{02}+\sigma_{01}$, where the edge maps are restrictions along the ordered face inclusions. Applying the endpoint boundary gives zero.
\end{atlasdefinition}

\begin{atlasdefinition}{\texttt{24-05}: 특이호몰로지 / Singular homology}
\label{def:24-05}
위상공간 (topological space) X와 정수 \(q\ge0\)에 대해, 실수 계수 특이 사슬복합체 (chain complex) \((C_\bullet(X;\mathbb R),\partial)\)의 호몰로지 (homology)를 \[H_q^{\mathrm{sing}}(X;\mathbb R):=\ker\partial_q/\operatorname{im}\partial_{q+1}\]로 정의한다. \(C_{-1}(X;\mathbb R)=0\), \(\partial_0=0\)으로 둔다.

\textbf{조건 확인 / Conditions.} 계수체 \ensuremath{\mathbb R}를 명시한 정의다. 0차의 분자는 C\ensuremath{{}_{0}}(X;\ensuremath{\mathbb R}) 전체다.

\textbf{직접 선행 / Prerequisites.} \hyperref[def:24-04]{\texttt{24-04} 특이 경계}; \hyperref[def:23-04]{\texttt{23-04} 호몰로지}

\exampleheading
The alternating sum of the four radial face maps is a singular two-cycle on $S^2$. Its homology class generates $H_2(S^2;\Z)\simeq\Z$. Unlike the simplicial complex, the singular chain complex does have three-chains; nontriviality here follows from the simplicial-to-singular comparison theorem (or the degree pairing), not from an assertion that $C_3^{\rm sing}=0$.
\end{atlasdefinition}

\begin{atlasdefinition}{\texttt{24-06}: 특이 코사슬 / Singular cochain}
\label{def:24-06}
X를 위상공간 (topological space), q를 0 이상의 정수라 하자. \(C_q(X;\mathbb R)\)는 표준 q-단체에서 X로 가는 연속사상들의 유한 실수 선형결합으로 이루어진 특이 사슬군이다. \(C^q(X;\mathbb R)=\operatorname{Hom}_{\mathbb R}(C_q(X;\mathbb R),\mathbb R)\)다. 즉 특이 q-사슬에 실수를 주는 모든 실수 선형함수의 공간이다.

\textbf{조건 확인 / Conditions.} 코사슬 자체에는 유한 지지 조건을 붙이지 않는다.

\textbf{직접 선행 / Prerequisites.} \hyperref[def:24-03]{\texttt{24-03} 특이 사슬군}; \hyperref[def:02-02]{\texttt{02-02} 아벨군 준동형}

\exampleheading
Fix $N\in S^2$. A singular real zero-cochain can be defined by $a(p)=1$ if p=N and 0 otherwise, and extended linearly to finite zero-chains. Singular cochains are arbitrary homomorphisms from chain groups; this example need not be a continuous function of p.
\end{atlasdefinition}

\begin{atlasdefinition}{\texttt{24-07}: 특이 코사슬 미분 / Singular coboundary operator}
\label{def:24-07}
X를 위상공간 (topological space), q\ensuremath{\ge}0을 정수라 하자. \(C_r(X;\mathbb R)\)를 실수 특이 사슬군, \(\partial_{r}:C_r\to C_{r-1}\)를 특이 경계라 하고 \(C^r(X;\mathbb R):=\operatorname{Hom}_{\mathbb R}(C_r(X;\mathbb R),\mathbb R)\)로 둔다. 특이 코사슬 미분 \(\delta^q:C^q(X;\mathbb R)\to C^{q+1}(X;\mathbb R)\)는 \[\delta^q(c):=c\circ\partial_{q+1}\]로 정의한다. 즉 모든 \(z\in C_{q+1}(X;\mathbb R)\)에 대해 \((\delta^qc)(z)=c(\partial_{q+1}z)\)다. 특히 z가 특이 (q+1)-단체 \ensuremath{\sigma}인 경우에도 이 식이 성립한다.

\textbf{조건 확인 / Conditions.} 여기서 \ensuremath{\sigma}는 q+1차 특이 단체 (singular simplex)다.

\textbf{직접 선행 / Prerequisites.} \hyperref[def:24-06]{\texttt{24-06} 특이 코사슬}; \hyperref[def:24-04]{\texttt{24-04} 특이 경계}

\exampleheading
For that zero-cochain a and a path $\gamma:[0,1]\to S^2$ from S to N, $(\delta a)(\gamma)=a(\partial\gamma)=a(N)-a(S)=1$. Thus the cochain differential is dual to the chain boundary, with the endpoint signs visible.
\end{atlasdefinition}

\begin{atlasdefinition}{\texttt{24-08}: 특이코호몰로지 / Singular cohomology}
\label{def:24-08}
X를 위상공간 (topological space), q\ensuremath{\ge}0을 정수라 하자. 실수 특이 사슬복합체 (chain complex) \((C_\bullet(X;\mathbb R),\partial)\)의 선형쌍대 \(C^r(X;\mathbb R):=\operatorname{Hom}_{\mathbb R}(C_r(X;\mathbb R),\mathbb R)\)에 \(\delta^r(c):=c\circ\partial_{r+1}\)를 준다. q차 실수 특이코호몰로지 (singular cohomology)는 \[H^q_{\mathrm{sing}}(X;\mathbb R):=\frac{\ker(\delta^q:C^q\to C^{q+1})}{\operatorname{im}(\delta^{q-1}:C^{q-1}\to C^q)}\]다. 모든 코사슬군은 \((X;\mathbb R)\)에 대한 것이며 \(C^{-1}=0\), \(\delta^{-1}=0\)으로 둔다.

\textbf{조건 확인 / Conditions.} 특이호몰로지 (singular homology)와 달리 위 첨자를 사용한다.

\textbf{직접 선행 / Prerequisites.} \hyperref[def:24-07]{\texttt{24-07} 특이 코사슬 미분}; \hyperref[def:04-04]{\texttt{04-04} 복합체의 코호몰로지}

\exampleheading
For the sphere, $H^0(S^2;\R)=\R$, $H^1(S^2;\R)=0$, and $H^2(S^2;\R)=\R$. The degree-two class is normalized to evaluate to 1 on the oriented fundamental cycle. Under the de Rham comparison theorem it corresponds to $[\Omega/(4\pi)]$.
\end{atlasdefinition}

\begin{atlasdefinition}{\texttt{25-01}: 추상 단체복합체 / Abstract simplicial complex}
\label{def:25-01}
꼭짓점 집합의 유한 부분집합들로 이루어진 모임 K로, \(\sigma\in K\)이고 \(\tau\subset\sigma\)이면 \(\tau\in K\)다. 각 원소를 단체라 한다.

\textbf{조건 확인 / Conditions.} 단체의 모든 부분집합도 단체여야 한다.

\textbf{직접 선행 / Prerequisites.} 집합·함수·실수·복소수 등 공통 배경.

\exampleheading
Let K consist of all proper subsets of $\{0,1,2,3\}$. It is an abstract simplicial complex: subsets of proper subsets are proper. Its realization is the boundary of a tetrahedron and hence homeomorphic to $S^2$. The full four-vertex subset is deliberately absent.
\end{atlasdefinition}

\begin{atlasdefinition}{\texttt{25-02}: 단체 사슬복합체 / Simplicial chain complex}
\label{def:25-02}
K를 추상 단체복합체 (simplicial complex), q를 0 이상의 정수라 하자. 꼭짓점 순서는 K의 꼭짓점 집합 전체에 정한다. \(C_{-1}(K;\mathbb R)=0\)으로 둔다. 꼭짓점에 전순서를 정하고 q+1개 꼭짓점의 단체 \([v_0,\ldots,v_q]\), \(v_0<\cdots<v_q\)를 기저로 하는 실수벡터공간 \(C_q(K;\mathbb R)\)를 취한다. q>0에서 \[\partial[v_0,\ldots,v_q]=\sum_{i=0}^q(-1)^i[v_0,\ldots,\widehat v_i,\ldots,v_q],\]\(\partial_0=0\)이며 경계를 선형 확장한다.

\textbf{조건 확인 / Conditions.} 사슬은 기저 단체들의 유한 실수 선형결합이다.

\textbf{직접 선행 / Prerequisites.} \hyperref[def:25-01]{\texttt{25-01} 추상 단체복합체}; \hyperref[def:23-01]{\texttt{23-01} 사슬복합체}

\exampleheading
For this K modeling $S^2$, choose increasing orientations on simplices. The chain groups are free on four vertices, six edges and four faces. For example $\partial[0,1,2]=[1,2]-[0,2]+[0,1]$ and $\partial[0,1]=[1]-[0]$.
\end{atlasdefinition}

\begin{atlasdefinition}{\texttt{25-03}: 단체호몰로지 / Simplicial homology}
\label{def:25-03}
K를 추상 단체복합체 (simplicial complex), q를 0 이상의 정수라 하자. \(C_r(K;\mathbb R)\)는 순서 있는 r-단체들의 유한 실수 선형결합 공간이고, \(\partial_r:C_r\to C_{r-1}\)는 한 꼭짓점을 지우는 면들을 교대부호로 합한 단체 경계다. \(C_{-1}=0\), \(\partial_0=0\)이다. \(H_q(K;\mathbb R)=H_q(C_\bullet(K;\mathbb R),\partial)\)로 정의하며, 단체 경계의 \(\ker\partial_q/\operatorname{im}\partial_{q+1}\)다.

\textbf{조건 확인 / Conditions.} 주어진 단체복합체 (simplicial complex)의 실수 사슬복합체 (chain complex)를 쓴다.

\textbf{직접 선행 / Prerequisites.} \hyperref[def:25-02]{\texttt{25-02} 단체 사슬복합체}; \hyperref[def:23-04]{\texttt{23-04} 호몰로지}

\exampleheading
For the tetrahedral boundary modeling $S^2$, the simplicial homology groups are $H_0(K;\Z)=\Z$, $H_1(K;\Z)=0$, $H_2(K;\Z)=\Z$. The oriented sum of the four faces is the degree-two generator; a face boundary is a zero class in degree one.
\end{atlasdefinition}

\begin{atlasdefinition}{\texttt{25-04}: 단체코호몰로지 / Simplicial cohomology}
\label{def:25-04}
K를 추상 단체복합체 (simplicial complex), q를 0 이상의 정수라 하자. \((C_\bullet(K;\mathbb R),\partial)\)는 K의 실수 단체 사슬복합체 (chain complex)이며 \(C_{-1}=0\), \(\partial_0=0\)이다. 음의 차수 코사슬도 0으로 둔다. \(C^q(K;\mathbb R)=\operatorname{Hom}_{\mathbb R}(C_q(K;\mathbb R),\mathbb R)\), \(\delta c=c\circ\partial_{q+1}\)로 둔 코사슬복합체 (cochain complex)의 코호몰로지 (cohomology) \(H^q(K;\mathbb R)=\ker\delta^q/\operatorname{im}\delta^{q-1}\)다.

\textbf{조건 확인 / Conditions.} 쌍대 미분은 차수를 올린다.

\textbf{직접 선행 / Prerequisites.} \hyperref[def:25-02]{\texttt{25-02} 단체 사슬복합체}; \hyperref[def:04-04]{\texttt{04-04} 복합체의 코호몰로지}

\exampleheading
For K modeling $S^2$, define a two-cochain a by $a([1,2,3])=1$ and a=0 on the other increasing faces. It is a cocycle since $C^3(K)=0$. It evaluates to 1 on the fundamental cycle, whereas every coboundary evaluates to zero, so its class is nonzero in $H^2(K;\Z)$.
\end{atlasdefinition}

\begin{atlasdefinition}{\texttt{25-05}: 덮개의 너브 / Nerve of a cover}
\label{def:25-05}
X를 위상공간 (topological space), \(\mathfrak U=(U_i)_{i\in I}\)를 X의 열린덮개 (open cover)라 하자. 너브 \(N(\mathfrak U)\)는 빈 단체와 다음 유한 지표집합들로 이루어진 추상 단체복합체 (simplicial complex)다. 비어 있지 않은 유한 \(J\subset I\)에 대해 \[J\in N(\mathfrak U)\quad\Longleftrightarrow\quad\bigcap_{i\in J}U_i\ne\varnothing.\]특히 꼭짓점은 U\ensuremath{{}_{i}}\ensuremath{\ne}\ensuremath{\varnothing}인 지표 i다.

\textbf{조건 확인 / Conditions.} 꼭짓점은 덮개의 지표다.

\textbf{직접 선행 / Prerequisites.} \hyperref[def:25-01]{\texttt{25-01} 추상 단체복합체}; \hyperref[def:01-03]{\texttt{01-03} 열린덮개}

\exampleheading
Cover the tetrahedral $S^2$ by the four open vertex stars $U_0,\ldots,U_3$. A collection of stars has nonempty intersection exactly when its vertex set is a face of K. Thus its nerve is K itself: all edges and triangular faces occur, but the fourfold intersection is empty.
\end{atlasdefinition}

\begin{atlasdefinition}{\texttt{25-06}: 수축 가능한 공간 / Contractible space}
\label{def:25-06}
공간 W에 어떤 \(x_0\in W\)와 연속사상 \(H:W\times[0,1]\to W\)가 존재하여 모든 \(x\in W\)에 대해 \(H(x,0)=x\), \(H(x,1)=x_0\)이면 W가 수축 가능하다.

\textbf{조건 확인 / Conditions.} H는 W 안에서 움직이는 연속사상이어야 한다.

\textbf{직접 선행 / Prerequisites.} \hyperref[def:06-01]{\texttt{06-01} 연속사상}

\exampleheading
The coordinate chart $U_0\subset\CP^1$ is contractible: in its coordinate z, $F(z,t)=(1-t)z$ contracts it to 0. This does not contract all of $\CP^1$, since the formula does not define that homotopy at infinity when t=1. The whole sphere is not contractible.
\end{atlasdefinition}

\begin{atlasdefinition}{\texttt{25-07}: 좋은 덮개 / Good cover}
\label{def:25-07}
X를 위상공간 (topological space)이라 하자. 열린덮개 (open cover)는 X 전체를 덮는 열린집합족을 뜻한다. 비어 있지 않은 모든 유한 교집합이 수축 가능한 열린덮개를 좋은 덮개라 한다.

\textbf{조건 확인 / Conditions.} 덮개 원소 하나도 유한 교집합에 포함된다.

\textbf{직접 선행 / Prerequisites.} \hyperref[def:25-06]{\texttt{25-06} 수축 가능한 공간}; \hyperref[def:01-03]{\texttt{01-03} 열린덮개}

\exampleheading
The four open vertex stars of the tetrahedral sphere form a good cover. Each nonempty finite intersection is an open star of a simplex and contracts to that simplex's barycentre by linear interpolation in barycentric coordinates, followed by the fixed homeomorphism to $S^2$. The standard two-chart cover is not good: its overlap is $\C^\times$, which is not contractible.
\end{atlasdefinition}

\subsection{후속 연구의 별도 가지 / Further study branches}

\begin{atlasdefinition}{\texttt{N24}: 사영 초곡면 / Projective hypersurface}
\label{def:N24}
n\ensuremath{\ge}1이고 F가 \(\mathbb C[Z_0,\ldots,Z_n]\)의 양의 차수 동차다항식이라 하자. \(V(F):=\{[Z]\in\mathbf{CP}^n:F(Z)=0\}\)를 그 사영 영점집합이라 한다. F가 square-free이면 이 reduced 초곡면의 차수는 deg F이다. 매끄럽다는 것은 V(F)의 모든 점에서 \((\partial F/\partial Z_0,\ldots,\partial F/\partial Z_n)\)가 동시에 0이 되지 않는 것이다.

\textbf{조건 확인 / Conditions.} \(V(F^2)=V(F)\)이지만 F\ensuremath{{}^{2}}의 중복도를 기록한 대상은 달라진다. 다항식을 제곱했다고 매끄러운 새 초곡면이 생기지는 않는다.

\textbf{직접 선행 / Prerequisites.} \hyperref[def:N23]{\texttt{N23} 복소사영공간}; \hyperref[def:26-01]{\texttt{26-01} 환}

\exampleheading
In $\CP^1$ with coordinates [S:T], the homogeneous polynomial T of degree one cuts out the projective hypersurface $\{T=0\}=\{[1:0]\}$. The polynomial $T^2$ has the same point set, but scheme-theoretically gives a double point; the set and its multiplicity data must be distinguished.
\end{atlasdefinition}

\begin{atlasdefinition}{\texttt{N25}: Segre 매장 / Segre embedding}
\label{def:N25}
m,n\ensuremath{\ge}0에 대해 Segre 매장은 \[\sigma:\mathbf{CP}^m\times\mathbf{CP}^n\longrightarrow\mathbf{CP}^{(m+1)(n+1)-1},\qquad([a],[b])\longmapsto[a_i b_j]_{0\le i\le m,\ 0\le j\le n}\]이다. a와 b의 대표를 비영 스칼라 배로 바꾸면 출력 전체에 비영 스칼라가 곱해져 같은 사영점을 준다. 그 상은 비영 행렬 \((c_{ij})\) 중 rank 1인 것들의 사영류이며 모든 2\ensuremath{\times}2 소행렬식이 0이라는 조건으로 표현된다.

\textbf{조건 확인 / Conditions.} 이 정의는 두 인자의 곱에 대한 것이다. 여러 인자는 tensor 성분의 곱으로 확장한다.

\textbf{직접 선행 / Prerequisites.} \hyperref[def:N23]{\texttt{N23} 복소사영공간}; \hyperref[def:B02]{\texttt{B02} 쌍대·텐서곱·외대수}

\exampleheading
The Segre map for two copies of $\CP^1$ sends $([a:b],[c:d])$ to $[ac:ad:bc:bd]\in\CP^3$. Restricting the second factor to [1:0] gives the explicit copy $[a:b]\mapsto[a:0:b:0]$ of $\CP^1$. General image coordinates satisfy $x_0x_3-x_1x_2=0$.
\end{atlasdefinition}

\begin{atlasdefinition}{\texttt{N26}: 유한차원 순수 양자상태 / Pure quantum state}
\label{def:N26}
H를 유한차원 복소 Hilbert 공간이라 하자. 순수상태는 H의 1차원 복소 부분공간, 즉 \(\mathbf P(H)\)의 원소다. 단위벡터 \(\psi\in H\)로 표현하면 \(\psi\)와 \(e^{i\theta}\psi\), \(\theta\in\mathbb R\)는 같은 상태다. 동등하게 직교사영 \(\rho_\psi(v)=\psi\langle\psi,v\rangle\)로 표현할 수 있다. 여기서는 내적이 첫 인자에 켤레선형, 둘째 인자에 선형이다.

\textbf{조건 확인 / Conditions.} 순수상태와 모든 density operator를 포함하는 혼합상태 공간을 구별한다.

\textbf{직접 선행 / Prerequisites.} \hyperref[def:N23]{\texttt{N23} 복소사영공간}; \hyperref[def:B03]{\texttt{B03} Hermitian 내적}

\exampleheading
For a two-level system, rays $[a:b]\in\CP^1$ represent pure states. On $a\ne0$ choose $\psi_z=(1,z)/\sqrt{1+|z|^2}$. Its Bloch vector is $(2\Re z,2\Im z,1-|z|^2)/(1+|z|^2)\in S^2$. Multiplying $\psi_z$ by a unit complex phase changes neither its ray nor this sphere point.
\end{atlasdefinition}

\begin{atlasdefinition}{\texttt{N27}: 곱상태와 순수상태의 얽힘 / Product state and pure-state entanglement}
\label{def:N27}
두 유한차원 복소 Hilbert 공간 H\ensuremath{{}_{a}},H\ensuremath{{}_{b}}의 복합계는 \(H_A\otimes_{\mathbb C}H_B\)로 나타낸다. Tensor product는 쌍선형사상 \((a,b)\mapsto a\otimes b\)을 가지며, 임의의 복소벡터공간 V로 가는 모든 쌍선형사상이 이 사상을 통해 유일한 선형사상으로 인수분해되는 벡터공간이다. 순수상태 \([\psi]\)가 곱상태라는 것은 비영 a,b가 존재하여 \([\psi]=[a\otimes b]\)인 것이다. 그렇지 않은 순수상태를 얽힌 상태라 한다.

\textbf{조건 확인 / Conditions.} 이분 순수상태에서는 곱상태들이 Segre의 상이다. 혼합상태의 separability 조건으로 그대로 옮기지 않는다.

\textbf{직접 선행 / Prerequisites.} \hyperref[def:N25]{\texttt{N25} Segre 매장}; \hyperref[def:N26]{\texttt{N26} 유한차원 순수 양자상태}; \hyperref[def:B02]{\texttt{B02} 쌍대·텐서곱·외대수}

\exampleheading
The family $[a:b]\in\CP^1\mapsto[a:0:0:b]\in\CP^3$ represents bipartite pure states $a|00\rangle+b|11\rangle$. Its coefficient matrix is $\operatorname{diag}(a,b)$. It is a product state exactly when ab=0; otherwise it has rank two and is entangled. Thus one fixed projective-line family contains both product endpoints and entangled interior points.
\end{atlasdefinition}

\clearpage
\section{연결 정리 / Optional connecting results}
정의와 정리는 구별하여 읽습니다. 아래 예시는 정리의 적용이며 일반 증명을 대체하지 않습니다.

\begin{atlasdefinition}{T01: 완전열을 잇는 정리 / Snake lemma and Mayer-Vietoris sequence}
아벨군의 가환도표에서 두 행이 짧은 완전열이고 세 수직사상이 \(a:A'\to A\), \(b:B'\to B\), \(c:C'\to C\)이면, snake lemma는 연결 준동형을 포함하는 완전열 \[0\to\ker a\to\ker b\to\ker c\to\operatorname{coker}a\to\operatorname{coker}b\to\operatorname{coker}c\to0\]을 준다. 코사슬복합체의 짧은 완전열은 코호몰로지의 긴 완전열을 유도한다. Mayer-Vietoris 완전열은 열린덮개에 관련된 복합체의 짧은 완전열에 이 사실을 적용하여 얻는다.
\exampleheading
For the standard two-chart cover of $\CP^1$, the de Rham connecting map sends the overlap class $[dz/z]$ to a class in $H^2(\CP^1;\C)$. Choosing the lift-difference convention $B_0-B_\infty=dz/z$, Stokes gives $\int\delta[dz/z]=\oint_{|z|=1}dz/z=2\pi i$.
\end{atlasdefinition}

\begin{atlasdefinition}{T02: de Rham 층의 완전성과 fine / Poincaré lemma and fine resolution}
M이 Hausdorff이고 제2가산인 매끄러운 다양체이면 실수 상수층과 미분형식층 사이의 열 \[0\to\underline{\mathbb R}\to\Omega_M^0\xrightarrow d\Omega_M^1\xrightarrow d\cdots\]은 층의 완전열이다. 여기서 \(\Omega_M^q(U)=\Omega^q(U)\)이다. 양의 차수에서의 국소 완전성은 Poincaré lemma로, 차수 0에서는 미분이 0인 매끄러운 함수가 국소상수라는 사실로 주어진다. 매끄러운 단위분할을 이용하면 각 미분형식층은 fine이다.
\exampleheading
The sequence $0\to\underline\R\to\mathcal A^0\to\mathcal A^1\to\mathcal A^2\to0$ on $S^2$ is a fine resolution: Poincar\'e lemma gives exactness, and subordinate partitions act by multiplication to make each $\mathcal A^k$ fine.
\end{atlasdefinition}

\begin{atlasdefinition}{T03: 층 코호몰로지와 acyclic 분해 / Foundations and acyclic resolutions}
위상공간 X 위 아벨군 층의 범주는 아벨범주이고 충분한 injective 대상을 갖는다. 전역단면함자 \(\Gamma(X,-)\)는 왼쪽 완전하다. Injective resolution으로 얻는 \(H^q(X,\mathcal F)\)는 분해의 선택에 표준적인 동형으로 독립적이며 \(H^0(X,\mathcal F)=\mathcal F(X)\)이다. 완전열 \(0\to\mathcal F\to\mathcal A^0\to\mathcal A^1\to\cdots\)에서 모든 \(\mathcal A^p\)가 \(\Gamma(X,-)\)-acyclic이면 그 전역단면 복합체의 코호몰로지는 \(H^q(X,\mathcal F)\)와 동형이다. Paracompact Hausdorff 공간의 fine 아벨군 층은 이 acyclicity를 만족한다.
\exampleheading
The fine de Rham resolution on $S^2$ is also an acyclic resolution, since its terms have no higher sheaf cohomology. Taking global sections therefore computes $H^q(S^2,\underline\R)$. No assertion that each term is an injective sheaf is needed.
\end{atlasdefinition}

\begin{atlasdefinition}{T04: Leray 비교정리 / Leray comparison theorem}
X 위 아벨군 층 \(\mathcal F\)와 열린덮개 \(\mathfrak U=\{U_i\}_{i\in I}\)를 잡는다. 모든 비어 있지 않은 유한 교집합 \(U_{i_0}\cap\cdots\cap U_{i_p}\)과 모든 \(q>0\)에 대해 \(H^q(U_{i_0}\cap\cdots\cap U_{i_p},\mathcal F|_{U_{i_0}\cap\cdots\cap U_{i_p}})=0\)이면 모든 \(p\ge0\)에서 \[\check H^p(\mathfrak U,\mathcal F)\cong H^p(X,\mathcal F).\]이 acyclicity는 줄기가 UFD라는 조건으로 대체되지 않는다.
\exampleheading
On analytic $\CP^1$, $U_0\simeq\C$, $U_\infty\simeq\C$, and their overlap $\C^\times$ are Stein. The restrictions of the locally free sheaf $\mathcal K_X$ have vanishing higher cohomology by Cartan's theorem B. Consequently this cover is $\mathcal K_X$-acyclic and its Čech calculation computes sheaf cohomology. The Stein acyclicity theorem is an explicitly used theorem, not a property of every open cover.
\end{atlasdefinition}

\begin{atlasdefinition}{T05: Nerve와 de Rham 비교 / Nerve comparison and de Rham theorem}
열린덮개의 모든 비어 있지 않은 유한 교집합이 연결되어 있으면, 실수 상수층의 Čech 코사슬복합체는 그 덮개의 nerve에 대한 실수 단체 코사슬복합체와 일치한다. Hausdorff 제2가산 매끄러운 다양체 M의 좋은 덮개에서는 각 교집합의 수축 가능성을 이용해 M의 코호몰로지를 nerve로 계산한다. de Rham 정리는 미분형식의 적분이 모든 \(q\ge0\)에서 \[H^q_{\mathrm{dR}}(M)\cong H^q_{\mathrm{sing}}(M;\mathbb R)\]를 유도한다는 정리다.
\exampleheading
The four open vertex stars of the tetrahedral sphere form a good cover. Each nonempty finite intersection is an open star of a simplex and contracts to that simplex's barycentre by linear interpolation in barycentric coordinates, followed by the fixed homeomorphism to $S^2$. The standard two-chart cover is not good: its overlap is $\C^\times$, which is not contractible. For the sphere, $H^0(S^2;\R)=\R$, $H^1(S^2;\R)=0$, and $H^2(S^2;\R)=\R$. The degree-two class is normalized to evaluate to 1 on the oriented fundamental cycle. Under the de Rham comparison theorem it corresponds to $[\Omega/(4\pi)]$.
\end{atlasdefinition}

\begin{atlasdefinition}{T06: 정칙 싹·Taylor 급수·완비화 / Holomorphic germs, Taylor series and completion}
리만곡면 X의 점 x에서 \(t(x)=0\)인 정칙 국소좌표 t를 고르면 \(\mathcal O_{X,x}\cong\mathbb C\{t\}\)이고, 극대아이디얼 \(\mathfrak m_x\)는 이 동형 아래 \((t)\)에 대응한다. 따라서 \(\mathfrak m_x\)-진 완비화는 \(\mathbb C[[t]]\)와 동형이다. 정칙 싹은 수렴 Taylor 급수로 결정되지만, 매끄러운 함수의 싹은 일반적으로 그 Taylor 급수만으로 결정되지 않는다. 여기서 급수 표현은 좌표의 선택에 의존한다.
\exampleheading
At $0\in\CP^1$, the compatible truncations $1+z+\cdots+z^k$ modulo $z^{k+1}$ determine the element $\sum_{n\ge0}z^n$ in $\widehat{\mathcal O}_{X,0}\simeq\C[[z]]$. The truncations of $\sum n!z^n$ also give an element of this completion even though no convergent germ has those coefficients.
\end{atlasdefinition}

\begin{atlasdefinition}{T09: 계량 독립성과 곡률 적분 / Metric independence and degree}
두 Hermitian 계량이 \(h'=e^\psi h\)로 연결되면 \ensuremath{\psi}는 전역 매끄러운 실함수이며 \(F_{h'}-F_h=d(\partial\psi)\)이다. 따라서 첫째 Chern 형식의 차이는 실수 완전형식이다. 연결된 콤팩트 리만곡면에서 선다발 L의 영이 아닌 유리형 단면 s를 택하면 \(\frac{i}{2\pi}\int_XF_{L,h}=\sum_p\operatorname{ord}_p(s)\). 첫째 등식은 국소 공식으로, 적분 등식은 영점·극을 제거한 영역에서 Stokes 정리로 확인한다.
\exampleheading
For $K_{\CP^1}$, $c_1(K,h)=\frac{i}{2\pi}F_K=-\Omega/(2\pi)$ and its integral is $-2$. For the dual tangent line bundle, the form is $+\Omega/(2\pi)$ and its integral is $+2$. The normalization and the choice of bundle both affect the sign.
\end{atlasdefinition}

\begin{atlasdefinition}{T10: P\ensuremath{{}^{1}} 발표의 세 표현 / Three descriptions on the projective line}
\(w=1/z\)에서 \(dw=-z^{-2}dz\)이고 \(\operatorname{div}(dz)=-2[\infty]\)이다. 따라서 \(K_{\mathbf P^1}\simeq\mathcal O(-2)\)이며 \(\frac{i}{2\pi}\int_{\mathbf P^1}F_{K,h}=-2\). 전이함수의 winding, 인자의 차수, 곡률 적분을 연결하는 계산 결과이며 세 용어의 정의를 동일시하지 않는다.
\exampleheading
For the canonical bundle of $\CP^1$ with the fixed sphere-induced metric, $F_{K,h}=-2(1+|z|^2)^{-2}dz\wedge d\bar z=i\Omega$. Thus $\frac{i}{2\pi}\int_{\CP^1}F_{K,h}=-\frac1{2\pi}\int_{S^2}\Omega=-2$. This uses the identification of the sphere with $\CP^1$ and its outward orientation.
\end{atlasdefinition}

\clearpage
\section{범위와 원고 정리 기록 / Scope and editorial notes}
\begin{itemize}
\item Supplied manuscript: 22 main definition boxes and 3 subdefinition boxes; every box now has an explicit sphere or projective-line example. The single [PROMPT] is fulfilled by S01.
\item Definition atlas: all 161 tagged entries, including N01--N30 and B01--B12, each have their own example. The common sphere conventions are fixed once to avoid conflicting signs and normalizations.
\item Definitions requiring categories or resolutions use sheaves, groups, chain complexes or functors on the fixed sphere. These are genuine instances of the abstract definitions, not assertions that a category is itself a point of the sphere.
\item The Segre and entanglement examples use explicit families parametrized by $\CP^1$; their target can be $\CP^3$. This is stated explicitly rather than pretending that every construction has target $\CP^1$.
\item The source's undefined reference to Definition 23 in the K\"ahler paragraph has been replaced by the existing type-decomposition reference. The open-cover subscript typo and inconsistent open-ball radius notation were corrected. The missing Demailly bibliography entry is supplied below.
\item The canonical curvature and the Gaussian curvature are distinguished; the dual tangent-line curvature has the opposite sign. Metric normalization matches the supplied $g_{i\bar j}=g_\C(\partial_{z^i},\partial_{\bar z^j})$ convention.
\end{itemize}
\begin{thebibliography}{9}
\bibitem{demailly} Jean-Pierre Demailly, \emph{Complex Analytic and Differential Geometry}, online monograph. Reference for Hermitian bundles, Chern connections and Chern--Weil theory.
\bibitem{stacks-injective} The Stacks Project, \emph{Abelian sheaves on a space}, Tag 01DF, \url{https://stacks.math.columbia.edu/tag/01DF}. Used for injective skyscrapers and enough injectives.
\bibitem{stacks-stalk} The Stacks Project, \emph{Skyscraper sheaves and stalks}, Tag 0099, \url{https://stacks.math.columbia.edu/tag/0099}.
\bibitem{stacks-line} The Stacks Project, \emph{Invertible modules}, Tag 01CR, \url{https://stacks.math.columbia.edu/tag/01CR}. The locally ringed hypothesis is retained when identifying invertible sheaves with locally free rank-one sheaves.
\end{thebibliography}
\end{document}

